% Pratique du secourisme pour la prévention des accidents cardiovasculaires — SVT Tle D — Cours signé OnBuch+
% Programme officiel MINESEC. Compiler avec : tectonic N14-secourisme-accidents-cardiovasculaires.tex
\newcommand{\DOCMATIERE}{SVT}
\newcommand{\DOCNIVEAU}{Tle D}
\newcommand{\DOCMODULE}{Module II : L'Éducation à la Santé — Le secourisme}
\newcommand{\DOCLECON}{N14}
\newcommand{\DOCTITRE}{Pratique du secourisme pour la prévention des accidents cardiovasculaires}
% OnBuch+ — préambule « cours signés » ENRICHI (compilé avec tectonic / XeLaTeX)
% Système visuel « Pop » d'OnBuch+ : fond crème (jamais blanc), bordures encre
% épaisses, ombres « dures » décalées, titres Archivo Black en capitales.
% Version MATHÉMATIQUES : graphes (pgfplots/tikz), boîtes pédagogiques
% (définition, propriété, méthode, attention, le savais-tu, exercice résolu,
% exercice, TP, bilan), page de garde et signature OnBuch+ ancrée en bas.
%
% Macros attendues dans main.tex AVANT \input{preamble} :
%   \DOCMATIERE  \DOCNIVEAU  \DOCMODULE  \DOCLECON (numéro)  \DOCTITRE
\documentclass[11pt]{article}

\usepackage{fontspec}
\usepackage[french]{babel}
\usepackage[a4paper,margin=2cm,top=3.4cm,bottom=2.8cm,headheight=1.4cm,footskip=1.3cm]{geometry}
\usepackage{amsmath,amssymb,amsfonts}
\usepackage{mathtools} % \xmapsto, \coloneqq... souvent utilisés par les modèles
\usepackage{cancel} % simplifications barrées (bilans de forces)
% \abs{z} (module, valeur absolue) et \norm{u} (norme), utilisés par les modèles
\providecommand{\abs}{}\let\abs\relax\DeclarePairedDelimiter{\abs}{\lvert}{\rvert}
\providecommand{\norm}{}\let\norm\relax\DeclarePairedDelimiter{\norm}{\lVert}{\rVert}
\usepackage[table]{xcolor} % [table] : fournit aussi \rowcolors (alternance de lignes)
\usepackage{pagecolor}
\usepackage{tikz}
\usetikzlibrary{arrows.meta,positioning,calc,shapes.geometric,shapes.misc,shapes.arrows,decorations.pathmorphing,decorations.pathreplacing,decorations.markings,patterns,shadows,mindmap,trees,fit,backgrounds,matrix,intersections,angles,quotes}
\usepackage{pgfplots}
\usepackage[version=4]{mhchem} % équations nucléaires \ce{^{235}_{92}U}
\usepackage[european,siunitx]{circuitikz} % schémas électriques
\pgfplotsset{compat=1.18}
\usepgfplotslibrary{fillbetween,groupplots} % aires entre courbes (calcul intégral)
\usepackage{enumitem}
\usepackage{fancyhdr}
% [most] charge toutes les bibliothèques tcolorbox (listings, minted, raster,
% documentation...) et gonfle inutilement la mémoire TeX sur un document long
% (~300 boîtes) : on ne charge que ce qui est réellement utilisé.
\usepackage[skins,breakable]{tcolorbox}
\usepackage{graphicx}
\usepackage{colortbl}
\usepackage{booktabs}
\usepackage{tabularx}
\usepackage{diagbox} % case d'en-tête diagonale des tableaux à double entrée
% Colonnes extensibles centrée / gauche / droite (C, L, R), souvent utilisées par les modèles
\newcolumntype{C}{>{\centering\arraybackslash}X}
\newcolumntype{L}{>{\raggedright\arraybackslash}X}
\newcolumntype{R}{>{\raggedleft\arraybackslash}X}
\usepackage{array}
\usepackage{multicol}
\usepackage{siunitx}
% parse-numbers=false : accepte aussi \SI{2,5 \times 10^{-3}}{...}
\sisetup{locale=FR,output-decimal-marker={,},group-minimum-digits=4,parse-numbers=false}
\DeclareSIUnit\ohm{\ensuremath{\symup{\Omega}}} % U+2126 absent de la police
\DeclareSIUnit\division{div} % réglages d'oscilloscope (ms/div, V/div)
\DeclareSIUnit\tr{tr} % tours (vitesse de rotation, ex. \SI{1500}{\tr\per\minute})
\DeclareSIUnit\molar{M} % concentration molaire (ex. \SI{3}{\molar}, \SI{1}{\micro\molar})
\DeclareSIUnit\an{an} % année (le modèle confond parfois avec \year, un compteur TeX) — ex. \SI{5}{\kilo\meter\per\an}
\DeclareSIUnit\FCFA{FCFA} % franc CFA (ex. \SI{500}{\FCFA\per\kilo\gram})
\DeclareSIUnit\degreeBrix{\textdegree Brix} % teneur en sucre d'un sirop/jus (réfractométrie)
\DeclareSIUnit\month{mois} % le modèle utilise parfois \month, un compteur TeX, comme unité de durée
\usepackage{qrcode}
% \degree : commande fréquemment utilisée par les modèles pour un angle ou une
% température en dehors de \SI{}{} (non fournie par les paquets chargés).
\providecommand{\degree}{^{\circ}}
% \qed (fin de démonstration), utilisé par les modèles sans amsthm
\providecommand{\qed}{\hfill$\square$}
% \barre : double barre des valeurs interdites dans un tableau de variations (tkz-tab)
\providecommand{\barre}{\ensuremath{\|}}
% environnement proof (amsthm non chargé) utilisé par les modèles
\ifdefined\proof\else\newenvironment{proof}[1][Démonstration]{\par\noindent\textit{#1.}\ }{\qed\par}\fi
\usepackage{lastpage}
\usepackage{hyperref}
\hypersetup{hidelinks}

% ============================================================
% Palette « Pop » OnBuch+ — mêmes hex que la classe P (plus_ui.dart)
% ============================================================
\definecolor{poporange}{HTML}{F59321}
\definecolor{popdark}{HTML}{E07A0C}
\definecolor{popcream}{HTML}{FFF6E8}
\definecolor{popink}{HTML}{1C1714}
\definecolor{poppurple}{HTML}{7A5AE0}
\definecolor{popgreen}{HTML}{2E9E6A}
\definecolor{poppink}{HTML}{E0457B}
\definecolor{popblue}{HTML}{2F7DE1}
\definecolor{popgold}{HTML}{C98A12}
\definecolor{popmuted}{HTML}{8A7F76}
\definecolor{popline}{HTML}{EADFD2}
\definecolor{popgray}{HTML}{8A7F76}
\colorlet{popgrey}{popgray}
% Alias tolérants (le modèle nomme parfois les couleurs différemment)
\colorlet{popwhite}{white}
\colorlet{popblack}{popink}
\colorlet{popred}{poppink}
\colorlet{popyellow}{popgold}
% Teintes claires (fonds de tableaux/encadrés) — le modèle nomme parfois
% les couleurs avec un suffixe L (ex. popblueL) sans les avoir définies.
\colorlet{poporangeL}{poporange!20}
\colorlet{popdarkL}{popdark!20}
\colorlet{poppurpleL}{poppurple!20}
\colorlet{popgreenL}{popgreen!20}
\colorlet{poppinkL}{poppink!20}
\colorlet{popblueL}{popblue!20}
\colorlet{popgoldL}{popgold!20}
\colorlet{popredL}{popred!20}
\colorlet{popgrayL}{popgray!20}
% Teintes claires pour fonds de tableaux / zones
\colorlet{popblueL}{popblue!10!white}
\colorlet{poporangeL}{poporange!14!white}
\colorlet{popgreenL}{popgreen!12!white}
\colorlet{poppurpleL}{poppurple!12!white}
\colorlet{poppinkL}{poppink!12!white}
\colorlet{popgoldL}{popgold!14!white}

\pagecolor{popcream}

% ============================================================
% Typographie
% ============================================================
\defaultfontfeatures{Ligatures=TeX}
\setmainfont{PlusJakartaSans-Medium.ttf}[Path=../../../fonts/,BoldFont=PlusJakartaSans-Bold.ttf]
% Maths en sans-serif (Fira Math) avec chiffres et lettres droites de Plus
% Jakarta, pour que les formules se fondent dans le texte.
\usepackage[math-style=ISO,bold-style=ISO]{unicode-math}
\setmathfont{FiraMath-Regular.otf}[Path=../../../fonts/]
\setmathfont{PlusJakartaSans-Medium.ttf}[Path=../../../fonts/,range={up/num,up/{Latin,latin},\mathup}]
\usepackage{icomma} % virgule décimale sans espace en mode math
% Caractères mathématiques Unicode absents des polices de texte (Plus Jakarta,
% Archivo Black) : rendus par leur équivalent mathématique, en texte comme en
% formule — sinon ils s'affichent en carré vide (ex. « Arithmétique dans ℤ »).
\usepackage{newunicodechar}
\newunicodechar{ℝ}{\ensuremath{\mathbb{R}}}
\newunicodechar{ℤ}{\ensuremath{\mathbb{Z}}}
\newunicodechar{ℂ}{\ensuremath{\mathbb{C}}}
\newunicodechar{ℕ}{\ensuremath{\mathbb{N}}}
\newunicodechar{ℚ}{\ensuremath{\mathbb{Q}}}
\newunicodechar{𝔻}{\ensuremath{\mathbb{D}}}
\newunicodechar{α}{\ensuremath{\alpha}}
\newunicodechar{β}{\ensuremath{\beta}}
\newunicodechar{γ}{\ensuremath{\gamma}}
\newunicodechar{δ}{\ensuremath{\delta}}
\newunicodechar{ε}{\ensuremath{\varepsilon}}
\newunicodechar{θ}{\ensuremath{\theta}}
\newunicodechar{λ}{\ensuremath{\lambda}}
\newunicodechar{μ}{\ensuremath{\mu}}
\newunicodechar{σ}{\ensuremath{\sigma}}
\newunicodechar{τ}{\ensuremath{\tau}}
\newunicodechar{φ}{\ensuremath{\varphi}}
\newunicodechar{ψ}{\ensuremath{\psi}}
\newunicodechar{ω}{\ensuremath{\omega}}
\newunicodechar{Σ}{\ensuremath{\Sigma}}
% majuscules produites par \MakeUppercase dans les titres (α-aminés -> Α)
\newunicodechar{Α}{\ensuremath{\alpha}}
\newunicodechar{Β}{\ensuremath{\beta}}
\newunicodechar{Γ}{\ensuremath{\Gamma}}
\newunicodechar{Δ}{\ensuremath{\Delta}}
\newunicodechar{Ε}{\ensuremath{\varepsilon}}
\newunicodechar{Θ}{\ensuremath{\Theta}}
\newunicodechar{Λ}{\ensuremath{\Lambda}}
\newunicodechar{Μ}{\ensuremath{\mu}}
\newunicodechar{Τ}{\ensuremath{\tau}}
\newunicodechar{Φ}{\ensuremath{\Phi}}
\newunicodechar{Ψ}{\ensuremath{\Psi}}
\newunicodechar{Ω}{\ensuremath{\Omega}}
\newunicodechar{↦}{\ensuremath{\mapsto}}
\newunicodechar{∈}{\ensuremath{\in}}
\newunicodechar{∉}{\ensuremath{\notin}}
\newunicodechar{∧}{\ensuremath{\wedge}}
\newunicodechar{⇔}{\ensuremath{\Leftrightarrow}}
\newunicodechar{⟺}{\ensuremath{\Longleftrightarrow}}
\newunicodechar{⇒}{\ensuremath{\Rightarrow}}
\newunicodechar{⟹}{\ensuremath{\Longrightarrow}}
\newunicodechar{∘}{\ensuremath{\circ}}
\newunicodechar{∪}{\ensuremath{\cup}}
\newunicodechar{∩}{\ensuremath{\cap}}
\newunicodechar{≡}{\ensuremath{\equiv}}
\newunicodechar{⊂}{\ensuremath{\subset}}
\newunicodechar{⊥}{\ensuremath{\perp}}
\newunicodechar{∃}{\ensuremath{\exists}}
\newunicodechar{∀}{\ensuremath{\forall}}
\newunicodechar{∛}{\ensuremath{\sqrt[3]{\,}}}
\newunicodechar{∜}{\ensuremath{\sqrt[4]{\,}}}
\newunicodechar{⃗}{\ensuremath{{}^{\to}}}
\newunicodechar{ⁿ}{\ifmmode^{n}\else\textsuperscript{\ensuremath{n}}\fi}
\newunicodechar{ˣ}{\ifmmode^{x}\else\textsuperscript{\ensuremath{x}}\fi}
\newunicodechar{ᵇ}{\ifmmode^{b}\else\textsuperscript{\ensuremath{b}}\fi}
\newunicodechar{ʳ}{\ifmmode^{r}\else\textsuperscript{\ensuremath{r}}\fi}
\newunicodechar{ᵘ}{\ifmmode^{u}\else\textsuperscript{\ensuremath{u}}\fi}
\newunicodechar{ʸ}{\ifmmode^{y}\else\textsuperscript{\ensuremath{y}}\fi}
\newunicodechar{ᵃ}{\ifmmode^{a}\else\textsuperscript{\ensuremath{a}}\fi}
\newunicodechar{ᵉ}{\ifmmode^{e}\else\textsuperscript{\ensuremath{e}}\fi}
\newunicodechar{ᵏ}{\ifmmode^{k}\else\textsuperscript{\ensuremath{k}}\fi}
\newunicodechar{ᵖ}{\ifmmode^{p}\else\textsuperscript{\ensuremath{p}}\fi}
\newunicodechar{ᴺ}{\ifmmode^{N}\else\textsuperscript{\ensuremath{N}}\fi}
\newunicodechar{ᶜ}{\ifmmode^{c}\else\textsuperscript{\ensuremath{c}}\fi}
\newunicodechar{ʰ}{\ifmmode^{h}\else\textsuperscript{\ensuremath{h}}\fi}
\newunicodechar{ᵠ}{\ifmmode^{q}\else\textsuperscript{\ensuremath{q}}\fi}
\newunicodechar{ₙ}{\ifmmode_{n}\else\textsubscript{\ensuremath{n}}\fi}
\newunicodechar{ₐ}{\ifmmode_{a}\else\textsubscript{\ensuremath{a}}\fi}
\newunicodechar{ᵢ}{\ifmmode_{i}\else\textsubscript{\ensuremath{i}}\fi}
\newunicodechar{ᵣ}{\ifmmode_{r}\else\textsubscript{\ensuremath{r}}\fi}
\newunicodechar{ₖ}{\ifmmode_{k}\else\textsubscript{\ensuremath{k}}\fi}
\newunicodechar{ᵦ}{\ifmmode_{\beta}\else\textsubscript{\ensuremath{\beta}}\fi}
\newunicodechar{ₓ}{\ifmmode_{x}\else\textsubscript{\ensuremath{x}}\fi}
\newfontfamily\popdisplay{ArchivoBlack-Regular.ttf}[Path=../../../fonts/]
\newfontfamily\poplogo{SpaceGrotesk-Bold.ttf}[Path=../../../fonts/]
\newfontfamily\popsemi{PlusJakartaSans-SemiBold.ttf}[Path=../../../fonts/]
\newfontfamily\popxbold{PlusJakartaSans-ExtraBold.ttf}[Path=../../../fonts/]
\newfontfamily\popxboldls{PlusJakartaSans-ExtraBold.ttf}[Path=../../../fonts/,LetterSpace=3.0]

\setlist[enumerate]{leftmargin=1.7em,itemsep=5pt,topsep=4pt}
\setlist[itemize]{leftmargin=1.7em,itemsep=4pt,topsep=4pt,label={\color{poporange}\textbullet}}
\setlength{\parindent}{0pt}
\setlength{\parskip}{5pt}
\renewcommand{\arraystretch}{1.25}

% pgfplots : style Pop par défaut
\pgfplotsset{
  every axis/.append style={axis line style={popink,line width=1pt},tick style={popink},
    grid style={popline,line width=0.5pt},label style={font=\small},tick label style={font=\footnotesize}},
  popcurve/.style={poporange,line width=1.8pt,no markers,smooth},
}
% Même style disponible dans un \draw TikZ simple (hors pgfplots)
\tikzset{popcurve/.style={poporange,line width=1.8pt}}
\tikzset{popgrid/.style={popline,very thin}}
\colorlet{popcurve}{poporange} % le nom du style est parfois utilisé comme couleur

% ============================================================
% Pill / squiggle
% ============================================================
\newcommand{\poppill}[3]{%
  \tikz[baseline=-0.55ex]\node[draw=popink,line width=1.2pt,rounded corners=2.6pt,%
    fill=#2,inner xsep=6.5pt,inner ysep=3pt]{%
    \popxboldls\fontsize{7}{7}\selectfont\color{#3}\MakeUppercase{#1}};%
}
\newcommand{\popsquiggle}[1]{%
  \tikz[baseline=-0.4ex]{
    \draw[#1,line width=2.6pt,line cap=round]
      plot [smooth] coordinates {(0,0.09) (0.34,0.01) (0.68,0.21) (1.02,0.01) (1.36,0.10)};
  }%
}
% Mot-clé surligné (marqueur orange sous le texte)
\newcommand{\cle}[1]{{\popxbold\color{popdark}#1}}

% ============================================================
% En-tête / pied
% ============================================================
\pagestyle{fancy}
\fancyhf{}
\renewcommand{\headrulewidth}{0pt}
\renewcommand{\footrulewidth}{0pt}
\lhead{\raisebox{-0.15ex}{\poplogo\fontsize{13}{13}\selectfont\color{popink}OnBuch\color{poporange}+}}
\rhead{\poppill{\DOCMATIERE\ \textbullet\ \DOCNIVEAU\ \textbullet\ Leçon \DOCLECON}{white}{popink}}
\lfoot{\popxbold\fontsize{7.6}{7.6}\selectfont\color{popmuted}\MakeUppercase{Cours signé OnBuch+}\ \textbullet\ {\popsemi\DOCTITRE}}
\rfoot{\tikz[baseline=-0.6ex]\node[rounded corners=8.5pt,fill=popink,minimum height=17pt,minimum width=17pt,inner xsep=5pt,inner sep=0]{\popxbold\fontsize{8}{8}\selectfont\color{popcream}\thepage\,/\,\pageref*{LastPage}};}

% ============================================================
% Titres
% ============================================================
\newcommand{\chaptitre}[2]{%
  \begin{tcolorbox}[enhanced,colback=poporange,colframe=popink,boxrule=2.6pt,arc=11pt,
      drop shadow=popink,left=15pt,right=15pt,top=12pt,bottom=11pt,before skip=3pt,after skip=18pt]
    \poppill{#2}{popink}{popcream}\par\vspace{9pt}
    {\raggedright\hyphenpenalty=10000\exhyphenpenalty=10000\popdisplay\fontsize{22}{24}\selectfont\color{popcream}\MakeUppercase{#1}\par}
  \end{tcolorbox}
}
% Section numérotée : pastille orange avec le numéro + titre Archivo
\newcounter{popsec}
\newcommand{\coursec}[1]{%
  \stepcounter{popsec}\setcounter{popsub}{0}%
  \par\vspace{16pt}\noindent
  \begin{minipage}{\linewidth}
  \tikz[baseline=-0.7ex]\node[draw=popink,line width=1.6pt,fill=poporange,rounded corners=4pt,minimum size=22pt,inner sep=0]{\popdisplay\fontsize{12}{12}\selectfont\color{popcream}\thepopsec};%
  \hspace{8pt}{\popdisplay\fontsize{15}{17}\selectfont\color{popink}\MakeUppercase{#1}}\par
  \vspace{4pt}\noindent{\color{poporange}\rule{36pt}{3.6pt}}
  \end{minipage}\par\nopagebreak\vspace{8pt}
}
\newcounter{popsub}
\newcommand{\courssub}[1]{%
  \stepcounter{popsub}%
  \par\vspace{9pt}\noindent{\popxbold\fontsize{11.5}{13}\selectfont\color{popink}{\color{poporange}\thepopsec.\thepopsub}\hspace{6pt}#1}\par\nopagebreak\vspace{4pt}
}

% ============================================================
% Boîtes pédagogiques — même recette : fond blanc, bordure épaisse colorée,
% ombre dure encre, badge plein. Titre optionnel [#1].
% ============================================================
\tcbset{popbase/.style={breakable,enhanced,colback=white,boxrule=2.3pt,arc=9pt,
  shadow={3pt}{-3pt}{0pt}{popink},left=14pt,right=14pt,top=11pt,bottom=11pt,before skip=11pt,after skip=15pt,
  fontupper=\popsemi\fontsize{10.6}{14.5}\selectfont\color{popink},
  fontlower=\popsemi\fontsize{10.6}{14.5}\selectfont\color{popink}}}
% Coche dessinée (le glyphe ✓ n'existe pas dans les polices du document)
\newcommand{\popcheck}[1]{\tikz[baseline=-0.5ex]\draw[#1,line width=1.6pt,line cap=round,line join=round](-0.13,0.0)--(-0.04,-0.09)--(0.13,0.1);}
\providecommand{\coche}{\popcheck{popgreen}}
\providecommand{\resultat}[1]{\par\smallskip\noindent\fcolorbox{popgreen}{popgreenL}{\parbox{\dimexpr\linewidth-2\fboxsep-2\fboxrule\relax}{\textbf{Résultat.} #1}}\par\smallskip}
\newcommand{\popboxhead}[3]{\poppill{#1}{#2}{white}\ifx\relax#3\relax\else\hspace{6pt}{\popxbold\fontsize{10.6}{12}\selectfont\color{popink}#3}\fi\par\vspace{7pt}}

\newtcolorbox{aretenir}[1][]{popbase,colframe=popgreen,before upper={\popboxhead{À retenir}{popgreen}{#1}}}
\newtcolorbox{exemplebox}[1][]{popbase,colframe=poporange,before upper={\popboxhead{Exemple}{poporange}{#1}}}
\newtcolorbox{definition}[1][]{popbase,colframe=popblue,before upper={\popboxhead{Définition}{popblue}{#1}}}
\newtcolorbox{propriete}[1][]{popbase,colframe=poppurple,before upper={\popboxhead{Propriété}{poppurple}{#1}}}
\newtcolorbox{methode}[1][]{popbase,colframe=popgold,before upper={\popboxhead{Méthode}{popgold}{#1}}}
\newtcolorbox{attention}[1][]{popbase,colframe=poppink,before upper={\popboxhead{Attention}{poppink}{#1}}}
\newtcolorbox{experience}[1][]{popbase,colframe=popblue,colback=popblueL,before upper={\popboxhead{Expérience}{popblue}{#1}}}
% « Le savais-tu ? » : carte violette pleine (contexte, culture, Cameroun)
\newtcolorbox{savaistu}[1][]{popbase,colback=poppurple,colframe=popink,
  fontupper=\popsemi\fontsize{10.4}{14.2}\selectfont\color{white},
  before upper={\poppill{Le savais-tu\,?}{white}{poppurple}\ifx\relax#1\relax\else\hspace{6pt}{\popxbold\color{white}#1}\fi\par\vspace{7pt}}}
% Exercice résolu : énoncé en haut, \tcblower puis la solution sur fond crème
\newcounter{popexo}\newcounter{popexr}
\newtcolorbox{exoresolu}[1][]{popbase,colframe=poporange,colbacklower=popcream,
  segmentation style={popink,line width=1.2pt,dashed},
  before upper={\stepcounter{popexr}\popboxhead{Exercice résolu \thepopexr}{poporange}{#1}},
  before lower={\poppill{Solution}{popink}{popcream}\par\vspace{6pt}}}
% Exercice (non résolu en place) — la correction est dans la partie Corrigés
\newtcolorbox{exercice}[1][]{popbase,colframe=popink,
  before upper={\stepcounter{popexo}\popboxhead{Exercice \thepopexo}{popink}{#1}}}
% Corrigé
\newtcolorbox{corrige}[1][]{popbase,colframe=popgreen,colback=popgreenL,
  before upper={\popboxhead{Corrigé}{popgreen}{#1}}}
% Objectifs / prérequis / bilan
\newtcolorbox{objectifs}{popbase,colframe=popink,colback=poporangeL,
  before upper={\popboxhead{Objectifs de la leçon}{popink}{}}}
\newtcolorbox{prerequis}{popbase,colframe=popink,colback=popblueL,
  before upper={\popboxhead{Prérequis}{popblue}{}}}
\newtcolorbox{bilan}{popbase,colframe=popink,colback=white,boxrule=2.8pt,
  before upper={\popboxhead{Fiche bilan}{poporange}{L'essentiel à connaître pour l'examen}}}
% Figure encadrée avec légende
\newtcolorbox{popfigure}[1][]{enhanced,breakable=false,colback=white,colframe=popline,boxrule=1.4pt,arc=8pt,
  left=10pt,right=10pt,top=10pt,bottom=8pt,before skip=10pt,after skip=12pt,halign=center,
  lower separated=false,halign lower=center,
  fontlower=\popsemi\fontsize{9}{11.5}\selectfont\color{popmuted},
  #1}
\newcommand{\legende}[1]{\par\vspace{6pt}{\popsemi\fontsize{9}{11.5}\selectfont\color{popmuted}#1}}

% ============================================================
% Signature OnBuch+
% ============================================================
\newcommand{\poplogoinline}[1]{{\poplogo\fontsize{#1}{#1}\selectfont\color{popink}OnBuch\color{poporange}+}}
\newcommand{\signatureonbuch}{%
  \begin{tcolorbox}[enhanced,colback=popink,colframe=popink,boxrule=2.2pt,arc=10pt,
      drop shadow=poporange,left=14pt,right=14pt,top=10pt,bottom=10pt,before skip=0pt,after skip=0pt]
    \tikz[baseline=-0.7ex]\node[circle,fill=popgreen,draw=popcream,line width=1.3pt,minimum size=22pt,inner sep=0]{\popcheck{white}};%
    \hspace{9pt}\begin{minipage}[c]{0.62\linewidth}
      {\popxbold\fontsize{12}{13}\selectfont\color{popcream}Signé\ \textbullet\ {\poplogo OnBuch\color{poporange}+}}\par\vspace{2pt}
      {\raggedright\popsemi\fontsize{8}{10.5}\selectfont\color{popline}\MakeUppercase{Équipe pédagogique OnBuch+}\\\MakeUppercase{Conforme au programme officiel MINESEC}\par}
    \end{minipage}\hfill
    {\poplogo\fontsize{22}{22}\selectfont\color{popcream}OnBuch\color{poporange}+}
  \end{tcolorbox}%
}

% ============================================================
% Page de garde
% #1 titre · #2 matière · #3 niveau · #4 module · #5 n° leçon · #6 durée ·
% #7 description / objectifs (texte ou itemize)
% ============================================================
\newcommand{\pagegarde}[7]{%
  \thispagestyle{empty}
  \newgeometry{a4paper,margin=1.7cm,top=1.6cm,bottom=1.4cm}%
  \begin{tikzpicture}[remember picture,overlay]
    \fill[poporange!18] ([xshift=-3.2cm,yshift=-3.6cm]current page.north east) circle (4.6cm);
    \fill[poppurple!14] ([xshift=2.4cm,yshift=7.6cm]current page.south west) circle (3.8cm);
  \end{tikzpicture}%
  {\poplogo\fontsize{30}{30}\selectfont\color{popink}OnBuch\color{poporange}+}\hfill
  \raisebox{0.6ex}{\poppill{Cours signé \textbullet\ Premium}{popink}{popcream}}\par
  \vspace{1pt}\popsquiggle{poppurple}\par
  \vspace{8pt}
  \begin{tcolorbox}[enhanced,colback=poporange,colframe=popink,boxrule=2.8pt,arc=13pt,
      drop shadow=popink,left=18pt,right=18pt,top=12pt,bottom=13pt,after skip=10pt]
    \poppill{#2 \textbullet\ #3}{popink}{popcream}\hspace{5pt}\poppill{#4}{white}{popink}\par\vspace{8pt}
    {\popdisplay\fontsize{14}{14}\selectfont\color{popink}LEÇON #5}\par\vspace{4pt}
    {\raggedright\hyphenpenalty=10000\exhyphenpenalty=10000\popdisplay\fontsize{25}{27}\selectfont\color{popcream}\MakeUppercase{#1}\par}
    \vspace{8pt}
    {\popxbold\fontsize{9.5}{9.5}\selectfont\color{popink}\MakeUppercase{Durée conseillée : #6}}
  \end{tcolorbox}
  \begin{tcolorbox}[enhanced,colback=white,colframe=popink,boxrule=2.2pt,arc=10pt,
      drop shadow=popink,left=16pt,right=16pt,top=10pt,bottom=10pt,after skip=8pt]
    \poppill{Dans cette leçon}{poppurple}{white}\par\vspace{5pt}
    \popsemi\fontsize{9.3}{12.6}\selectfont\color{popink}#7
  \end{tcolorbox}
  \noindent\begin{minipage}[t]{0.487\linewidth}
    \begin{tcolorbox}[enhanced,colback=popgreen,colframe=popink,boxrule=2.2pt,arc=10pt,
        drop shadow=popink,left=11pt,right=11pt,top=10pt,bottom=10pt,height=3.1cm,valign=center]
      \tcbox[colback=white,colframe=popink,boxrule=1.3pt,arc=5pt,boxsep=0pt,nobeforeafter,
        left=3pt,right=3pt,top=3pt,bottom=3pt,valign=center]{\qrcode[height=1.9cm]{https://play.google.com/store/apps/details?id=com.luvvix.onbuch&hl=fr}}%
      \hspace{8pt}\begin{minipage}[c]{0.5\linewidth}
        {\popdisplay\fontsize{11}{12.5}\selectfont\color{white}\MakeUppercase{Télécharge OnBuch}}\par\vspace{4pt}
        {\raggedright\popsemi\fontsize{8.4}{11}\selectfont\color{white}Gratuit \textbullet\ Google Play\\Tuteur IA, annales, concours, résultats\par}
      \end{minipage}
    \end{tcolorbox}
  \end{minipage}\hfill
  \begin{minipage}[t]{0.487\linewidth}
    \begin{tcolorbox}[enhanced,colback=poppurple,colframe=popink,boxrule=2.2pt,arc=10pt,
        drop shadow=popink,left=11pt,right=11pt,top=10pt,bottom=10pt,height=3.1cm,valign=center]
      \tikz[baseline=-0.7ex]\node[circle,fill=white,draw=popink,line width=1.3pt,minimum size=1.6cm,inner sep=0]{\popdisplay\fontsize{24}{24}\selectfont\color{poppurple}+};%
      \hspace{8pt}\begin{minipage}[c]{0.6\linewidth}
        {\popdisplay\fontsize{11}{12.5}\selectfont\color{white}\MakeUppercase{Passe à OnBuch+}}\par\vspace{4pt}
        {\raggedright\popsemi\fontsize{8.4}{11}\selectfont\color{white}Tous les cours signés, Tuteur IA illimité, fiches PDF — onglet OnBuch+ de l'app\par}
      \end{minipage}
    \end{tcolorbox}
  \end{minipage}
  \vspace{8pt}
  \popcontenu
  \vfill
  \signatureonbuch
  \restoregeometry
  \newpage
}

% Rangée « Ce cours contient » (page de garde)
\newcommand{\poptile}[3]{% #1 couleur #2 gros texte #3 légende
  \begin{tcolorbox}[enhanced,colback=#1,colframe=popink,boxrule=2pt,arc=9pt,shadow={2.5pt}{-2.5pt}{0pt}{popink},
      left=6pt,right=6pt,top=9pt,bottom=9pt,halign=center,valign=center,height=1.75cm,width=\linewidth,nobeforeafter]
    {\popdisplay\fontsize{12.5}{13}\selectfont\color{white}\MakeUppercase{#2}}\par\vspace{3pt}
    {\centering\popsemi\fontsize{7.8}{9.5}\selectfont\color{white}#3\par}
  \end{tcolorbox}}
\newcommand{\popcontenu}{%
  \par\noindent{\popxboldls\fontsize{7.5}{7.5}\selectfont\color{popmuted}CE COURS SIGNÉ CONTIENT}\par\vspace{7pt}
  \noindent\begin{minipage}[t]{0.235\linewidth}\poptile{popblue}{Cours}{complet, illustré, pas à pas}\end{minipage}\hfill
  \begin{minipage}[t]{0.235\linewidth}\poptile{poporange}{Méthodes}{et exercices résolus}\end{minipage}\hfill
  \begin{minipage}[t]{0.235\linewidth}\poptile{poppink}{Exercices}{type examen + corrigés}\end{minipage}\hfill
  \begin{minipage}[t]{0.235\linewidth}\poptile{popgreen}{Fiche bilan}{l'essentiel à retenir}\end{minipage}\par}

% Sommaire
\newtcolorbox{sommaire}{enhanced,colback=white,colframe=popink,boxrule=2.2pt,arc=10pt,
  drop shadow=popink,left=18pt,right=18pt,top=13pt,bottom=13pt,before skip=6pt,after skip=20pt,
  before upper={\poppill{Sommaire}{poppurple}{white}\par\vspace{9pt}\popsemi\fontsize{10.6}{14.5}\selectfont\color{popink}}}

% Signature de fin de cours — poussée tout en bas de la dernière page
\newcommand{\findecours}{%
  \par\vfill\nopagebreak
  \begin{center}\popsquiggle{poporange}\end{center}\vspace{4pt}
  \signatureonbuch
}
\AtBeginDocument{\def\llbracket{\mathopen{[\mkern-3.5mu[}}\def\rrbracket{\mathclose{]\mkern-3.5mu]}}} % intervalles d'entiers (glyphe ⟦ absent de la police)
% Glyphes absents de Fira Math : redéfinis à partir de glyphes disponibles
\AtBeginDocument{%
  \def\setminus{\mathbin{\backslash}}\let\smallsetminus\setminus
  \def\vdots{\mathord{\vbox{\baselineskip3.2pt\lineskiplimit0pt\kern4pt\hbox{.}\hbox{.}\hbox{.}}}}%
  \def\ddots{\mathinner{\mkern1mu\raise6.5pt\hbox{.}\mkern2mu\raise3.6pt\hbox{.}\mkern2mu\raise0.7pt\hbox{.}\mkern1mu}}%
  \def\triangle{\mathord{\scalebox{0.9}{$\symup{\Delta}$}}}%
  \def\nabla{\mathord{\raisebox{\depth}{\scalebox{1}[-1]{$\symup{\Delta}$}}}}%
  \def\checkmark{\popcheck{popdark}}%
  \def\star{\mathbin{\ast}}\def\bigstar{\mathord{\ast}}%
  \def\diamond{\mathbin{\scalebox{0.6}{$\Diamond$}}}%
  \def\bigcup{\mathop{\mathchoice{\raisebox{-0.3em}{\scalebox{1.7}{$\cup$}}}{\scalebox{1.25}{$\cup$}}{\cup}{\cup}}\displaylimits}%
  \def\bigcap{\mathop{\mathchoice{\raisebox{-0.3em}{\scalebox{1.7}{$\cap$}}}{\scalebox{1.25}{$\cap$}}{\cap}{\cap}}\displaylimits}%
  \def\lesseqqgtr{\mathrel{\lessgtr}}\def\bigoplus{\mathop{\oplus}\displaylimits}%
}
\providecommand{\ssi}{si et seulement si} % abréviation fréquente
\AtBeginDocument{\providecommand{\wideparen}{\overparen}} % arc de cercle

%% FIGFIX-BEGIN v5 — correctifs globaux des figures (docs/onbuch-generation/tools/figfix)
\usepackage{adjustbox}\usepackage{etoolbox}\tcbuselibrary{raster}
% toute figure TikZ/pgfplots plus large que la ligne est réduite à la largeur disponible
\BeforeBeginEnvironment{tikzpicture}{\begin{adjustbox}{max width=\linewidth}}
\AfterEndEnvironment{tikzpicture}{\end{adjustbox}}
% légendes pgfplots lisibles par-dessus les courbes
\pgfplotsset{every axis/.append style={legend style={fill=white,fill opacity=0.92,text opacity=1,draw=black!15,inner sep=3pt}}}
% étiquettes d'arêtes (syntaxe edge["..."]) avec fond blanc
\tikzset{every edge quotes/.append style={fill=white,inner sep=1.2pt}}
%% FIGFIX-END
\begin{document}

\pagegarde{Pratique du secourisme pour la prévention des accidents cardiovasculaires}{SVT}{Tle D}{Module II : L'Éducation à la Santé — Le secourisme}{N14}{4 h}{Cette leçon forme le citoyen-secouriste : reconnaissance et conduite à tenir face aux trois urgences vitales les plus fréquentes (arrêt cardiaque, étouffement total, perte de connaissance). Elle articule les bases physiologiques de la circulation et de la ventilation aux gestes techniques de premiers secours, dans le respect de la chaîne de survie.
\vspace{4pt}
\begin{multicols}{2}\small
\begin{itemize}
\item À la fin de cette leçon, je sais expliquer pourquoi l'arrêt cardiaque entraîne la mort cérébrale en quelques minutes (à partir de faits expérimentaux).
\item À la fin de cette leçon, je sais reconnaître un arrêt cardiaque (inconscience, absence de respiration).
\item À la fin de cette leçon, je sais pratiquer la chaîne de survie : alerter, masser (RCP), défibriller.
\item À la fin de cette leçon, je sais reconnaître un étouffement total et le distinguer d'un étouffement partiel.
\item À la fin de cette leçon, je sais pratiquer les gestes qui sauvent en cas d'étouffement total chez l'adulte (claques dans le dos, manœuvre de Heimlich).
\item À la fin de cette leçon, je sais adapter ces gestes au nourrisson et à la femme enceinte.
\end{itemize}\end{multicols}}

\begin{sommaire}
\begin{multicols}{2}
\begin{enumerate}
\item Les accidents cardiovasculaires : ampleur et enjeu de la chaîne de survie
\item L'arrêt cardiaque : bases physiologiques et reconnaissance
\item Les gestes qui sauvent en cas d'arrêt cardiaque : la RCP
\item L'étouffement total : mécanisme et gestes selon la victime
\item La perte de connaissance : reconnaissance et position latérale de sécurité
\item Synthèse : la conduite à tenir du citoyen-secouriste
\item Activité d'intégration
\item Méthodes et exercices résolus
\item Exercices
\item Corrigés des exercices
\item Fiche bilan
\end{enumerate}
\end{multicols}
\end{sommaire}

\begin{objectifs}
\begin{itemize}
\item À la fin de cette leçon, je sais expliquer pourquoi l'arrêt cardiaque entraîne la mort cérébrale en quelques minutes (à partir de faits expérimentaux).
\item À la fin de cette leçon, je sais reconnaître un arrêt cardiaque (inconscience, absence de respiration).
\item À la fin de cette leçon, je sais pratiquer la chaîne de survie : alerter, masser (RCP), défibriller.
\item À la fin de cette leçon, je sais reconnaître un étouffement total et le distinguer d'un étouffement partiel.
\item À la fin de cette leçon, je sais pratiquer les gestes qui sauvent en cas d'étouffement total chez l'adulte (claques dans le dos, manœuvre de Heimlich).
\item À la fin de cette leçon, je sais adapter ces gestes au nourrisson et à la femme enceinte.
\item À la fin de cette leçon, je sais reconnaître une perte de connaissance et placer la victime en position latérale de sécurité (PLS).
\item À la fin de cette leçon, je sais alerter efficacement les secours (message d'alerte complet) et assister la victime avec empathie.
\item À la fin de cette leçon, je sais adopter les comportements de prévention des accidents cardiovasculaires.
\end{itemize}
\end{objectifs}

\begin{prerequis}
\begin{itemize}
\item Anatomie et fonctionnement du cœur et de la circulation sanguine (classe de Première)
\item Mécanismes de la ventilation pulmonaire et trajet de l'air (classe de Première)
\item Besoins en dioxygène des cellules et rôle de la respiration cellulaire
\item Notions de base sur le système nerveux et la conscience
\item Conduite à tenir générale face à un accident : protéger, alerter, secourir (début d'année, module secourisme)
\end{itemize}
\end{prerequis}

\newpage

\coursec{Les accidents cardiovasculaires : ampleur et enjeu de la chaîne de survie}

\courssub{Situation d'accroche et problématique}

Imagine la scène : un après-midi de week-end au stade Omnisports de Yaoundé. Les Lions Indomptables jouent un match décisif. Soudain, en plein effort, un joueur s'effondre brutalement au milieu du terrain. Il ne bouge plus, ne respire plus. Le silence tombe sur les gradins, puis la panique : des cris, des téléphones portables sortis pour filmer, mais \emph{personne n'ose toucher la victime}, personne ne sait quoi faire. Les minutes s'égrènent, interminables. L'ambulance arrive enfin, mais le temps perdu est irréversible.

Cette scène, hélas trop réelle, se rejoue chaque jour au Cameroun : dans nos marchés, nos écoles, nos domiciles, sur nos routes. Selon l'Organisation Mondiale de la Santé (OMS), les \cle{maladies cardiovasculaires} (MCV) sont la première cause de mortalité dans le monde, avec environ \num{17,9} millions de décès par an, soit 32 \% de l'ensemble des décès mondiaux. Au Cameroun, la prévalence de l'hypertension artérielle, principal facteur de risque, touche près de \textbf{30 \% des adultes} (environ un sur trois), souvent à leur insu.

\begin{attention}[Problématique]
Dans une urgence vitale (arrêt cardiaque, étouffement total, perte de connaissance), \textbf{que doit faire un témoin ordinaire dans les premières minutes, avant l'arrivée des secours médicalisés, pour maximiser les chances de survie de la victime sans lui nuire} ?
\end{attention}

Cette section pose les bases épidémiologiques et physiologiques indispensables pour comprendre \emph{pourquoi} chaque seconde compte et \emph{comment} s'organise la \cle{chaîne de survie}.

\courssub{Qu'est-ce qu'un accident cardiovasculaire ? Définitions et distinctions cruciales}

Le terme « accident cardiovasculaire » (ACV) regroupe en réalité plusieurs pathologies distinctes par leur mécanisme, leur gravité immédiate et leur prise en charge. En secourisme, \textbf{ne pas les confondre conditionne le geste juste}.

\begin{definition}[Arrêt cardiaque (AC)]
L'\cle{arrêt cardiaque} est l'\textbf{arrêt brutal et inattendu de l'activité mécanique du cœur} (contraction ventriculaire efficace), entraînant l'\textbf{arrêt immédiat de la circulation sanguine}. La victime est \textbf{inconsciente}, \textbf{ne respire pas} (ou présente une respiration agonique, bruyante, inefficace) et \textbf{ne réagit pas} à la stimulation. C'est une \textbf{urgence vitale absolue} : sans intervention dans les minutes qui suivent, le décès est certain.
\end{definition}

\begin{definition}[Infarctus du myocarde (IDM) ou crise cardiaque]
L'\cle{infarctus du myocarde} est la \textbf{nécrose (mort) d'une partie du muscle cardiaque} (myocarde) due à l'obstruction brutale d'une artère coronaire (le plus souvent par un thrombus sur plaque d'athérome). Le cœur \textbf{bat encore} (souvent de façon anarchique), la victime est \textbf{généralement consciente}, se plaint d'une \textbf{douleur thoracique intense, constrictive, irradiant vers le bras gauche, la mâchoire ou le dos}, accompagnée de sueurs froides, nausées, angoisse. C'est une \textbf{urgence médicale} nécessitant une hospitalisation rapide (unité de soins intensifs cardiologiques), mais \emph{pas} une réanimation cardio-pulmonaire (RCP) tant que le cœur bat.
\end{definition}

\begin{definition}[Accident vasculaire cérébral (AVC) — \emph{Complément utile au Bac}]
L'\cle{accident vasculaire cérébral} (AVC) est une \textbf{atteinte brutale d'une zone du cerveau} par interruption de son irrigation sanguine (AVC ischémique, 80 \% des cas : thrombose ou embolie) ou par rupture d'un vaisseau (AVC hémorragique). La victime présente des \textbf{signes neurologiques focaux soudains} : paralysie d'un côté du corps (hémiparésie), trouble de la parole (aphasie), déviation de la bouche, trouble de la vision, céphalée brutale intense. Le test \textbf{FAST} (Face, Arm, Speech, Time) permet un dépistage rapide. C'est une \textbf{urgence médicale} : « \emph{Time is brain} » (le temps, c'est du cerveau).
\end{definition}

\begin{attention}[Piège majeur : Arrêt cardiaque vs Infarctus du myocarde]
\begin{itemize}
    \item \textbf{Arrêt cardiaque} : Cœur \textbf{arrêté} $\rightarrow$ Victime \textbf{inconsciente, ne respire pas} $\rightarrow$ \textbf{RCP IMMÉDIATE (massage + DAE)}.
    \item \textbf{Infarctus} : Cœur \textbf{bat (mal)} $\rightarrow$ Victime \textbf{consciente, douleur thoracique} $\rightarrow$ \textbf{APPEL 112/119, repos absolu, ASPIRINE si non contre-indiquée, PAS DE MASSAGE}.
\end{itemize}
Confondre les deux conduit soit à ne pas masser un cœur arrêté (décès certain), soit à masser un cœur qui bat (risque de lésions graves, fibrillation ventriculaire déclenchée).
\end{attention}

\courssub{Épidémiologie : une prévalence croissante au Cameroun et dans le monde}

Les maladies cardiovasculaires ne sont plus l'apanage des pays riches. La transition épidémiologique frappe durement l'Afrique subsaharienne.

\begin{savaistu}[Le « tueur silencieux » au Cameroun]
Au Cameroun, l'hypertension artérielle (HTA) touche \textbf{environ un adulte sur trois (29,7 \% selon l'enquête STEPS 2013, confirmée par des études récentes)}. Elle est souvent \textbf{asymptomatique} pendant des années, d'où son surnom de « \emph{tueur silencieux} ». Elle est la première cause d'AVC et la deuxième cause d'insuffisance cardiaque dans nos services de cardiologie (Hôpital Général de Yaoundé, Laquintinie de Douala). Le diabète (prévalence $\approx$ 6--8 \% en milieu urbain) et l'obésité (surpoids/obésité $\approx$ 30 \% des adultes urbains) alimentent cette épidémie.
\end{savaistu}

\begin{center}
\begin{tabularx}{\linewidth}{|l|X|X|}
\hline
\rowcolor{popblueL} \textbf{Facteur de risque} & \textbf{Prévalence estimée au Cameroun (adultes)} & \textbf{Impact pathophysiologique principal} \\
\hline
Hypertension artérielle (HTA) & $\sim$ 30 \% & Surcharge de post-charge $\rightarrow$ hypertrophie ventriculaire gauche, rigidité artérielle, rupture plaque d'athérome \\
Tabagisme & $\sim$ 10--15 \% (homme) & Endothélium dysfonctionnel, pro-thrombotique, spasme coronaire, oxydation LDL \\
Diabète / Dysglycémie & $\sim$ 6--8 \% & Glycation protéines, micro/macro-angiopathie, dyslipidémie atherogène \\
Sédentarité & $>$ 50 \% (urbain) & Déconditionnement cardiovasculaire, prise de poids, résistance à l'insuline \\
Alimentation riche sel/graisses & Fréquente & Rétention hydrosodée (HTA), athérogenèse (LDL-c élevé) \\
Alcool excessif & Variable & Cardiomyopathie dilatée, HTA, arythmies \\
Stress psychosocial & Difficile à quantifier & Activation sympathique chronique, inflammation, troubles du rythme \\
\hline
\end{tabularx}
\end{center}
\legende{Tableau 1 : Principaux facteurs de risque cardiovasculaires au Cameroun et leur impact. Source : OMS, Enquête STEPS Cameroun 2013, registres hospitaliers.}

\courssub{Bases physiologiques : pourquoi l'arrêt cardiaque tue en quelques minutes}

Pour saisir l'urgence, rappelons la physiologie de la circulation.

\begin{experience}[Observation : Conséquences de l'arrêt du flux sanguin cérébral]
\textbf{Matériel :} Modèle physiologique (schéma) ou document : courbe de l'EEG (électroencéphalogramme) et de la conscience lors d'une syncope vasovagale (arrêt transitoire du flux) vs arrêt cardiaque.
\textbf{Observation :}
\begin{itemize}
    \item $t = 0$ : Cœur s'arrête $\rightarrow$ Pression artérielle chute à 0 mmHg.
    \item $t = 10$--$15$ s : Perte de conscience (arrêt fonction corticale par anoxie).
    \item $t = 30$--$60$ s : Arrêt respiratoire (défaillance centres bulbaire par anoxie).
    \item $t = 3$--$5$ min : Début des lésions neuronales \textbf{irréversibles} (mort cellulaire par excitotoxicité, défaillance ATP).
    \item $t > 10$ min : Mort encéphalique quasi certaine.
\end{itemize}
\textbf{Interprétation :} Le cerveau ne stocke \textbf{pas} d'oxygène ni de glucose. Il consomme 20 \% de l'O$_2$ pour 2 \% de la masse corporelle. L'arrêt de la perfusion stoppe l'apport d'O$_2$ et l'évacuation du CO$_2$ $\rightarrow$ défaillance de la phosphorylation oxydative $\rightarrow$ chute de l'ATP $\rightarrow$ défaillance pompes ioniques (Na$^+$/K$^+$) $\rightarrow$ dépolarisation massive, influx Ca$^{2+}$, activation enzymes lytiques $\rightarrow$ nécrose/apoptose.
\textbf{Conclusion :} \textbf{La survie neurologique intacte dépend de la restauration d'un flux sanguin minimal (perfusion cérébrale) dans les 3 à 5 premières minutes.} D'où l'impératif du massage cardiaque externe (MCE) immédiat.
\end{experience}

\courssub{La chaîne de survie : quatre maillons pour une chance}

Face à l'arrêt cardiaque, la survie ne repose pas sur un acte unique, mais sur une \textbf{suite d'actions coordonnées} : la \cle{chaîne de survie}. Chaque maillon est indispensable ; le maillon faible détermine l'issue.

\begin{popfigure}
\begin{tikzpicture}[
    node distance=1.5cm and 1.2cm,
    maillon/.style={draw=popblue, thick, fill=popblueL, rounded corners=8pt, minimum width=3.2cm, minimum height=1.8cm, align=center, font=\sffamily\bfseries},
    fleche/.style={->, >=Stealth, thick, popdark, shorten >=2pt, shorten <=2pt},
    numero/.style={circle, fill=poporange, text=white, font=\bfseries\small, minimum width=0.7cm, minimum height=0.7cm, inner sep=0pt}
]
% Nœuds
\node[maillon, align=center] (m1){Alerter\\ \small{Appeler le 112 / 119  Donner infos précises}};
\node[maillon, right=of m1, align=center] (m2){Masser\\ \small{RCP : 30 compressions / 2 insufflations  100--120 /min, 5--6 cm}};
\node[maillon, right=of m2, align=center] (m3){Défibriller\\ \small{DAE public (DEA)  Choc si indiqué}};
\node[maillon, right=of m3, align=center] (m4){Secours médicalisés\\ \small{SAMU / Pompiers  Soins avancés, transport}};
\node[numero, anchor=north west, xshift=2pt, yshift=-2pt] at (m1.north west) {1};
\node[numero, anchor=north west, xshift=2pt, yshift=-2pt] at (m2.north west) {2};
\node[numero, anchor=north west, xshift=2pt, yshift=-2pt] at (m3.north west) {3};
\node[numero, anchor=north west, xshift=2pt, yshift=-2pt] at (m4.north west) {4};

% Flèches
\draw[fleche] (m1.east) -- (m2.west) node[midway, above, font=\small, popdark] {Sans délai};
\draw[fleche] (m2.east) -- (m3.west) node[midway, above, font=\small, popdark] {Dès dispo};
\draw[fleche] (m3.east) -- (m4.west) node[midway, above, font=\small, popdark] {Relais};

% Texte sous les maillons
\node[below=0.3cm of m1, font=\small\itshape, text=popmuted, align=center] {Témoin};
\node[below=0.3cm of m2, font=\small\itshape, text=popmuted, align=center] {Témoin / Secouriste};
\node[below=0.3cm of m3, font=\small\itshape, text=popmuted, align=center] {Témoin / Secouriste};
\node[below=0.3cm of m4, font=\small\itshape, text=popmuted, align=center] {Professionnels};
\end{tikzpicture}
\legende{Schéma 1 : La chaîne de survie (version grand public). Les 4 maillons sont solidaires : la rupture d'un seul annule les efforts des autres.}
\end{popfigure}

\begin{aretenir}[Les 4 maillons de la chaîne de survie]
\begin{enumerate}
    \item \textbf{Alerter (Appeler le 112 ou 119 au Cameroun)} : Dès la reconnaissance de l'arrêt cardiaque (inconscient, ne respire pas). Donner : \emph{Où ? Quoi ? Combien ? État ? Ne pas raccrocher le premier.}
    \item \textbf{Masser (RCP précoce)} : Massage cardiaque externe (MCE) immédiat, \textbf{même imparfait}, maintient un débit cardiaque résiduel ($\sim$ 20--30 \% du normal) perfusant le cerveau et le cœur lui-même. \textbf{Ne pas s'arrêter} sauf arrivée DAE, secours, ou épuisement.
    \item \textbf{Défibriller (DAE précoce)} : Le DAE (Défibrillateur Automatisé Externe) analyse le rythme et délivre un choc si fibrillation ventriculaire (FV) ou tachycardie ventriculaire (TV) sans pouls. \textbf{Chaque minute de retard au choc réduit la survie de 10 \%}.
    \item \textbf{Secours médicalisés (Soins avancés)} : Intubation, médicaments (adrénaline, amiodarone), prise en charge post-réanimation (hypothermie thérapeutique, coronarographie).
\end{enumerate}
\end{aretenir}

\courssub{Le facteur temps : la courbe implacable de la survie}

La relation entre le délai de début de la RCP/défibrillation et la probabilité de survie à la sortie d'hôpital est l'un des faits les mieux établis en médecine d'urgence.

\begin{popfigure}
\begin{tikzpicture}
\begin{axis}[
    width=\linewidth,
    height=9cm,
    domain=0:15,
    samples=200,
    axis lines=middle,
    xlabel={Délai avant début RCP / DAE (min)},
    ylabel={Probabilité de survie (\% à la sortie d'hôpital)},
    xmin=0, xmax=15.5,
    ymin=0, ymax=105,
    xtick={0,1,2,3,4,5,6,7,8,9,10,11,12,13,14,15},
    ytick={0,10,20,30,40,50,60,70,80,90,100},
    grid=major,
    grid style={popline},
    tick label style={font=\small},
    label style={font=\small},
    clip=false
]
% Courbe théorique : survie = 100 * exp(-k*t) avec k ~ 0.15 pour RCP seul, k ~ 0.1 pour RCP+DAE précoce
% Sans RCP : chute très raide. Avec RCP immédiat : départ plus haut, pente moins raide.
% On trace deux courbes : "Sans RCP" et "RCP immédiat + DAE rapide"
\addplot[popcurve, very thick, popred, domain=0:12] {90*exp(-0.22*x)}; % Sans RCP
\addplot[popcurve, very thick, popgreen, domain=0:15] {75*exp(-0.10*x)}; % RCP + DAE

% Annotations
\node[popred, font=\small\bfseries, anchor=west] at (axis cs:1, {90*exp(-0.22*1)}) {Sans RCP};
\node[popgreen, font=\small\bfseries, anchor=west] at (axis cs:1, {75*exp(-0.10*1)}) {RCP immédiate + DAE rapide};

% Zone critique 3-5 min
\draw[poporange, dashed, thick] (axis cs:3,0) -- (axis cs:3,100);
\draw[poporange, dashed, thick] (axis cs:5,0) -- (axis cs:5,100);
\node[poporange, font=\tiny, rotate=90, anchor=south] at (axis cs:4, 50) {Fenêtre critique lésions cérébrales};

% Point 10%/min
\node[align=center, font=\footnotesize, popdark] at (axis cs:12, 20) {Perte $\approx$ 10 \% / min\\ sans RCP};
\end{axis}
\end{tikzpicture}
\legende{Figure 1 : Probabilité de survie à la sortie d'hôpital en fonction du délai d'intervention. Courbe rouge : absence de RCP (chute brutale, survie $\approx$ 0 \% à 10 min). Courbe verte : RCP immédiate par témoin + DAE précoce (survie doublée à triplée). La zone ombrée 3--5 min marque le seuil des lésions cérébrales irréversibles. Données synthétiques issues d'études registre (ex. Registre CARES, études nordiques).}
\end{popfigure}

\begin{exemplebox}[Exemple chiffré : l'impact du témoin]
Dans une ville comme Douala ou Yaoundé, le délai moyen d'arrivée des secours (SAMU, pompiers) est souvent \textbf{supérieur à 15--20 minutes} (embouteillages, distance, un seul numéro d'appel saturé).
\begin{itemize}
    \item \textbf{Scénario A (Témoin passif)} : Appel secours seulement. RCP débutée à $t = 18$ min. Survie $\approx$ \textbf{0 -- 1 \%}.
    \item \textbf{Scénario B (Témoin formé)} : Appel + RCP immédiate ($t = 1$ min) + DAE à $t = 4$ min (borne publique marché/mairie). Survie $\approx$ \textbf{30 -- 40 \%}.
\end{itemize}
\textbf{Le témoin est le premier maillon, et le plus critique.} Une RCP immédiate \textbf{double à triple} les chances de survie.
\end{exemplebox}

\courssub{Le rôle du témoin-citoyen : premier maillon, acteur majeur}

\begin{propriete}[Réalité de l'arrêt cardiaque extra-hospitalier (ACRH)]
\begin{itemize}
    \item \textbf{70 à 80 \% des arrêts cardiaques surviennent à domicile ou lieu public} (hors hôpital), \textbf{devant des témoins} (famille, collègues, passants).
    \item Au Cameroun, le taux de RCP par témoin (bystander CPR) reste \textbf{très faible ($<$ 10 \%)} par méconnaissance, peur de mal faire, peur juridique, ou tabous culturels (toucher un inconnu, femme/homme).
    \item La \textbf{Loi camerounaise (Code pénal, art. 273 ; Loi sur la protection civile)} protège le secouriste de bonne foi : \emph{« N'est pas responsable celui qui, face à un danger imminent, porte secours à une personne en péril, sauf faute intentionnelle ou négligence caractérisée. »}
\end{itemize}
\end{propriete}

\begin{methode}[Devenir un témoin actif : les 3 « A »]
\begin{enumerate}
    \item \textbf{Apprendre} : Suivre une formation PSC1 (Prévention et Secours Civiques de niveau 1) ou équivalent (Croix-Rouge Cameroun, Protection Civile, associations agréées). C'est accessible à tous, souvent gratuit ou peu coûteux.
    \item \textbf{Agir} : Surmonter la sidération. \textbf{Agir imparfaitement vaut mieux que ne rien faire.} Un massage « trop fort » ou « trop vite » perfuse encore le cerveau ; ne rien faire garantit la mort.
    \item \textbf{Alentir} : Repérer les DAE (Défibrillateurs Automatisés Externes) dans votre environnement (mairies, aéroports, grands magasins, stades, certaines gares). Applications mobiles : \emph{Staying Alive}, \emph{DAE RespondER} (cartographie participative).
\end{enumerate}
\end{methode}

\courssub{Prévention primaire : agir en amont pour ne pas avoir à secourir}

Le secourisme est la \textbf{dernière barrière} (prévention tertiaire). La meilleure lutte contre l'arrêt cardiaque reste la \textbf{prévention primaire} : réduire les facteurs de risque.

\begin{aretenir}[Prévention primaire cardiovasculaire : les 5 piliers (OMS / Société Camerounaise de Cardiologie)]
\begin{enumerate}
    \item \textbf{Activité physique} : $\ge$ 150 min/semaine d'endurance modérée (marche rapide, vélo, danse) + renforcement musculaire 2$\times$/sem. \emph{Ex. : Marcher 30 min aller-retour au travail/marché.}
    \item \textbf{Alimentation équilibrée} : Réduire sel ($<$ 5 g/j = 1 c. à café), sucres ajoutés, graisses saturées/trans. Augmenter fruits, légumes, légumineuses, poisson, céréales complètes. \emph{Attention au « ndolé » trop salé/huileux, aux bouillons cubes (glutamate + sel), aux fritures quotidiennes.}
    \item \textbf{Zéro tabac / Alcool modéré} : Arrêt total tabac (sevrage aidé possible). Alcool : $\le$ 2 verres/jour (homme), $\le$ 1 (femme), pas tous les jours.
    \item \textbf{Dépistage HTA / Diabète / Dyslipidémie} : \textbf{Contrôler sa TA au moins 1 fois/an} (pharmacie, centre de santé, mutuelle). HTA $\ge$ 140/90 mmHg $\rightarrow$ prise en charge. Glycémie à jeun $\ge$ 1,26 g/L $\rightarrow$ diabète.
    \item \textbf{Gestion du stress / Sommeil} : 7--8 h nuit, techniques relaxation, soutien social.
\end{enumerate}
\end{aretenir}

\begin{exoresolu}[Cas concret : M. K., 52 ans, commerçant au Marché Central de Douala]
\textbf{Profil :} M. K. fume 10 cigarettes/jour depuis 30 ans, boit 2--3 bières chaque soir, mange gras/salé (poissons braisés, plantains frits), ne fait pas de sport, ne s'est jamais fait mesurer la TA. Son père est mort « brutalement » à 55 ans.
\textbf{Question :} Évaluez son risque cardiovasculaire global et proposez un plan d'action prioritaire.
\textbf{Solution :}
\begin{enumerate}
    \item \textbf{Facteurs de risque majeurs présents :} Âge $>$ 50 ans (H), Tabagisme actif, Alcool excessif, Sédentarité, Alimentation atherogène, Antécédent familial de mort subite (père $<$ 60 ans). \textbf{Risque TRÈS ÉLEVÉ} (Score SCORE ou Framingham estimé $>$ 20 \% à 10 ans).
    \item \textbf{Actions prioritaires (ordres de grandeur d'impact) :}
    \begin{enumerate}
        \item \textbf{Arrêt tabac} (gain espérance de vie $\sim$ 10 ans, risque IDM divisé par 2 à 1 an).
        \item \textbf{Mesure TA urgente} (probable HTA non diagnostiquée). Si confirmée : traitement à vie + hygiéno-diététique.
        \item \textbf{Réduction alcool} $\le$ 1 verre/jour.
        \item \textbf{Marche 30 min/jour} (intégré au travail : stationnement plus loin, escaliers).
        \item \textbf{Bilan lipidique / glycémie} (dépistage diabète/dyslipidémie).
    \end{enumerate}
    \item \textbf{Message secourisme :} M. K. et son entourage (famille, collègues de stand) \textbf{doivent connaître la RCP et l'emplacement du DAE le plus proche}. Son risque d'ACR est réel.
\end{enumerate}
\end{exoresolu}

\begin{attention}[Erreurs fréquentes en prévention / dépistage]
\begin{itemize}
    \item « Je me sens bien, donc ma tension est bonne. » $\rightarrow$ \textbf{FAUX.} HTA = asymptomatique dans 90 \% des cas au stade initial.
    \item « Je suis mince, je ne peux pas faire d'hypertension / diabète. » $\rightarrow$ \textbf{FAUX.} Facteurs génétiques, sel, stress, alcool agissent indépendamment du poids.
    \item « Le sel de cuisine est le seul problème. » $\rightarrow$ \textbf{FAUX.} Au Cameroun, \textbf{bouillons cubes, sauces soja, conserves, plats préparés, pain} apportent 70--80 \% du sel caché.
    \item « Le DAE est dangereux, il peut tuer. » $\rightarrow$ \textbf{FAUX.} Le DAE \textbf{n'analyse et ne choque QUE si fibrillation ventriculaire}. Sur un rythme normal ou asystolie, il dit « Pas de choc conseillé ». Risque nul pour la victime ou l'utilisateur si consignes respectées.
\end{itemize}
\end{attention}

\courssub{Synthèse de la section 1}

\begin{itemize}
    \item Les \textbf{accidents cardiovasculaires} (arrêt cardiaque, infarctus, AVC) sont la première cause de mortalité au Cameroun et dans le monde, portés par une épidémie de facteurs de risque (HTA, tabac, diabète, sédentarité, alimentation).
    \item L'\textbf{arrêt cardiaque} se distingue de l'infarctus par l'\textbf{absence de conscience et de respiration} (cœur arrêté) $\rightarrow$ \textbf{RCP immédiate}. L'infarctus : victime consciente, douleur thoracique $\rightarrow$ \textbf{Appel 112, repos, aspirine}.
    \item La \textbf{chaîne de survie} (Alerter -- Masser -- Défibriller -- Secours médicalisés) structure la réponse. \textbf{Le témoin est le maillon 1 et 2.}
    \item \textbf{Chaque minute sans RCP réduit la survie de $\sim$ 10 \%}. Une RCP immédiate double à triple les chances de survie.
    \item La \textbf{prévention primaire} (hygiène de vie, dépistage HTA/diabète) reste l'arme la plus puissante pour réduire l'incidence de ces accidents.
\end{itemize}

\begin{exercice}[Entraînement : Reconnaissance et premiers réflexes]
Pour chaque situation, indiquez : (1) Pathologie probable, (2) Conduite à tenir immédiate (ordre des actions), (3) Erreur à ne pas commettre.
\begin{enumerate}
    \item M. B., 65 ans, hypertendu connu, se plaint brutalement d'une douleur thoracique en étau irradiant au bras gauche, transpire, est anxieux mais conscient et respire normalement.
    \item Mme K., 40 ans, s'effondre au bureau. Elle ne répond pas, ne bouge pas, son thorax ne se soulève pas, vous entendez un bruit de ronflement/gaspillage (respiration agonique).
    \item Un enfant de 3 ans joue avec une pièce de monnaie, se met à tousser violemment, ne peut plus parler, devient cyanotique, panique.
\end{enumerate}
\end{exercice}

\begin{corrige}[Exercice : Reconnaissance et premiers réflexes]
\begin{enumerate}
    \item \textbf{Infarctus du myocarde (IDM).} (1) Appeler 112/119 $\rightarrow$ (2) Installer en position demi-assise (confort respiratoire) $\rightarrow$ (3) Donner 300 mg aspirine à croquer (si pas d'allergie/ulcère) $\rightarrow$ (4) Surveiller, rassurer, préparer DAE au cas où dégradation en AC. \textbf{Erreur :} Faire marcher le patient, lui donner à boire/manger, masser le cœur (inutile et dangereux).
    \item \textbf{Arrêt cardiaque (AC).} La respiration agonique \textbf{NE EST PAS} une respiration normale. (1) Appeler 112/119 (ou faire appeler) $\rightarrow$ (2) \textbf{COMMENCER RCP IMMÉDIATEMENT} (30 compressions / 2 insufflations) $\rightarrow$ (3) Envoyer chercher DAE $\rightarrow$ (4) Continuer jusqu'à secours/DAE/réveil. \textbf{Erreur :} Attendre les secours sans masser, « vérifier le pouls » (perte de temps, peu fiable pour un laïc), arrêter pour « bouche-à-bouche » seul.
    \item \textbf{Étouffement total (corps étranger des voies aériennes supérieures).} (1) \textbf{5 claques dans le dos} (incliné en avant) $\rightarrow$ (2) Si échec : \textbf{5 compressions abdominales (Heimlich)} $\rightarrow$ (3) Alterner jusqu'à expulsion ou perte de connaissance. Si perte de connaissance $\rightarrow$ RCP (regarder dans la bouche avant insufflations). \textbf{Erreur :} Donner des tapes dans le dos \emph{debout} (risque de faire glisser l'objet plus bas), tenter d'extraire l'objet à l'aveugle avec les doigts (risque de le pousser), faire boire de l'eau.
\end{enumerate}
\end{corrige}

\coursec{L'arrêt cardiaque : bases physiologiques et reconnaissance}

Dans la section précédente, nous avons mesuré l'ampleur des accidents cardiovasculaires au Cameroun et dans le monde, et compris que la \cle{chaîne de survie} repose sur la rapidité d'exécution de quelques gestes simples par le premier témoin. Avant d'apprendre \emph{comment} réaliser ces gestes (réanimation cardio-pulmonaire, défibrillation), il est indispensable de comprendre \emph{pourquoi} ils fonctionnent et \emph{comment} reconnaître sans ambiguïté la situation qui les justifie. Cette section pose les bases physiologiques de l'arrêt cardiaque et la méthode rigoureuse de sa reconnaissance.

\courssub{Physiologie de l'arrêt cardiaque : pourquoi le temps compte}

\begin{experience}[Expériences historiques : l'intolérance du cerveau à l'anoxie]
Dès la fin du XIX\textsuperscript{e} siècle, des physiologistes (comme \textbf{François-Franck} en 1894, puis \textbf{Schneider} et \textbf{Kouwenhoven} dans les années 1950-60) ont réalisé des expériences de \textbf{ligature des gros vaisseaux} (aorte, veines caves) ou d'arrêt circulatoire contrôlé chez le chien et le singe.
\par \textbf{Protocole :} On stoppe brutalement la circulation sanguine (pinçage de l'aorte ou fibrillation ventriculaire induite par courant alternatif) chez l'animal anesthésié. On enregistre l'activité électrique cérébrale (EEG), la pression artérielle et on observe le comportement à la réanimation.
\par \textbf{Observations majeures :}
\begin{itemize}
    \item \textbf{0 à 10-15 s :} Perte de conscience (évanouissement). L'EEG s'aplatit (silence électrique cortical).
    \item \textbf{3 à 5 min :} Si la circulation n'est pas rétablie, les lésions neuronales deviennent \textbf{irréversibles} (mort neuronale, œdème cérébral). Au-delà de 5-6 min, la réanimation ne ramène qu'un « état végétatif » ou la mort.
    \item \textbf{Expérience clé (Kouwenhoven et al., 1960) :} La compression rythmée du sternum (massage cardiaque externe) maintient une pression artérielle carotidienne moyenne de \SI{30}{mmHg} (contre \SI{80}{mmHg} normal) et un débit cérébral d'environ \textbf{25 à 30 \% du débit normal}. Ce débit partiel suffit à retarder de plusieurs minutes le seuil des lésions irréversibles.
\end{itemize}
\par \textbf{Interprétation :} Le cerveau, qui ne représente que \SI{2}{\%} de la masse corporelle, consomme \SI{20}{\%} de l'oxygène total (\SI{50}{ml/min/100g} de tissu). Il ne possède \textbf{aucune réserve} en dioxygène (pas de myoglobine, glycogène très faible) ni en ATP. L'interruption de l'apport sanguin (\cle{ischémie}) entraîne en quelques secondes l'épuisement de l'ATP, la défaillance des pompes Na$^+$/K$^+$, l'entrée massive d'eau et de Ca$^{2+}$ dans le neurone, puis la nécrose.
\legende{Expériences fondatrices : l'arrêt circulatoire entraîne une perte de conscience en $\sim$10 s et des lésions cérébrales irréversibles dès 3-5 min. Le massage cardiaque externe (RCP) maintient un débit résiduel ($\sim$25-30 \%) qui repousse ce seuil critique.}
\end{experience}

\begin{popfigure}
\begin{tikzpicture}[scale=0.9, every node/.style={font=\small}]
    % Axes : x = temps (min), y = récupération (%) ; 100% = 5cm
    \draw[->] (0,0) -- (10,0) node[right] {Durée de l'arrêt circulatoire (min)};
    \draw[->] (0,0) -- (0,5.5) node[above] {Récupération neurologique (\%)};
    % Curve without CPR : 100 * exp(-x/1.5) scaled *0.05
    \draw[popred, thick, domain=0:8, samples=100] plot (\x, {100*exp(-\x/1.5)*0.05});
    \node[popred, anchor=west] at (8, {100*exp(-8/1.5)*0.05}) {Sans RCP : chute brutale};
    % Curve with CPR : 100 * exp(-x/4) scaled *0.05
    \draw[popblue, thick, domain=0:8, samples=100] plot (\x, {100*exp(-\x/4)*0.05});
    \node[popblue, anchor=west] at (8, {100*exp(-8/4)*0.05}) {Avec RCP immédiate : déclin ralenti};
    % Threshold lines
    \draw[dashed, popmuted] (0,1) -- (8,1) node[right] {Seuil lésions irréversibles ($\sim$20\%)};
    \draw[dashed, popmuted] (3,0) -- (3,5.5) node[above] {3 min};
    \draw[dashed, popmuted] (5,0) -- (5,5.5) node[above] {5 min};
    % Points at t=0 (100% = 5cm)
    \fill[popred] (0,5) circle (2pt) node[above left] {100\%};
    \fill[popblue] (0,5) circle (2pt);
\end{tikzpicture}
\legende{Évolution conceptuelle de la récupération neurologique en fonction du temps d'anoxie cérébrale complète (rouge) ou avec RCP immédiate (bleu). La RCP décale la courbe vers la droite, offrant une « fenêtre thérapeutique » élargie pour la défibrillation.}
\end{popfigure}

\begin{definition}[Arrêt cardiaque et anoxie cérébrale]
L'\textbf{arrêt cardiaque} (AC) est la cessation brutale et inattendue de l'activité mécanique efficace du cœur, entraînant l'arrêt de la circulation sanguine. L'\textbf{anoxie cérébrale} qui en résulte provoque une perte de conscience en \textbf{10 à 15 secondes} et des lésions irréversibles dès \textbf{3 à 5 minutes} sans réanimation. La \cle{réanimation cardio-pulmonaire (RCP)} précoce maintient un débit sanguin cérébral résiduel (\SIrange{25}{30}{\%} du normal) qui ralentit la nécrose neuronale et prolonge la durée pendant laquelle une défibrillation peut réussir.
\end{definition}

\courssub{Mécanismes de l'arrêt cardiaque : quand le cœur s'emballe ou s'arrête}

\begin{definition}[Automatisme cardiaque et fibrillation ventriculaire]
Le cœur possède un \textbf{automatisme} propre : le \textbf{nœud sinusal} (pacemaker physiologique) génère spontanément des potentiels d'action réguliers (\SI{60}{à 100}{bpm} au repos) qui se propagent aux oreillettes, au nœud atrio-ventriculaire, au faisceau de His et aux ventricules, assurant une contraction synchronisée et efficace.
\par La \textbf{fibrillation ventriculaire (FV)} est une activité électrique \textbf{anarchique, rapide (\SI{300}{à 500}{bpm}) et désorganisée} des ventricules. Les fibres musculaires se contractent de façon asynchrone : le cœur « tremble » comme un « sac de vers » et ne peut plus éjecter le sang. C'est la cause la plus fréquente (\textbf{environ \SIrange{70}{80}{\%}}) des arrêts cardiaques soudains de l'adulte.
\end{definition}

\begin{savaistu}[Naissance de la RCP moderne : Kouwenhoven, Knickerbocker et Jude (1960)]
En 1960, à l'Université Johns Hopkins (USA), l'ingénieur \textbf{William Kouwenhoven} et les chirurgiens \textbf{Guy Knickerbocker} et \textbf{James Jude} publient l'article fondateur : « \emph{Closed-chest cardiac massage} ». Ils démontrent sur le chien, puis valident sur l'homme, qu'une \textbf{compression rythmée du sternum} (60/min, enfoncement \SI{4}{à 5}{cm}) génère un débit sanguin suffisant pour maintenir la vie. Ils associent cette technique à la \textbf{ventilation bouche-à-bouche} (décrite par Safar et Elam en 1958) et à la \textbf{défibrillation externe} (Zoll, 1956). C'est la naissance du trio « \textbf{Masser - Ventiler - Défibriller} » qui sauve des millions de vies depuis.
\end{savaistu}

\begin{propriete}[Causes principales d'arrêt cardiaque chez l'adulte (Savoir admis)]
L'arrêt cardiaque résulte le plus souvent d'une pathologie coronaire sous-jacente, mais les circonstances varient :
\begin{itemize}
    \item \textbf{Fibrillation ventriculaire (FV) / Tachycardie ventriculaire sans pouls (TVsp) :} Cause n°1 (\SIrange{70}{80}{\%}) : infarctus du myocarde aigu, cardiopathie ischémique chronique.
    \item \textbf{Asystolie / Activité électrique sans pouls (AESP) :} Fin d'évolution de la FV non traitée, ou cause primaire : hypoxie sévère (noyade, étouffement, asthme), hypovolémie massive (hémorragie), tamponnade, embolie pulmonaire massive, hypo/hyperkaliémie, hypothermie, intoxication, traumatisme thoracique.
    \item \textbf{Circonstances spécifiques au Cameroun :} Électrocution (câbles haute tension tombés, branchements anarchiques), noyade (fleuves, barrages, piscines non surveillées), traumatismes routiers graves, complications obstétricales (hémorragie du post-partum), paludisme cérébral sévère chez l'enfant.
\end{itemize}
\end{propriete}

\begin{experience}[Hémodynamique pendant la RCP : la preuve du débit résiduel]
\par \textbf{Contexte :} Mesures invasives (cathétérisme artériel, sonde de Swan-Ganz) chez l'homme en arrêt cardiaque réanimé, ou modèle porcin instrumenté.
\par \textbf{Résultats types pendant une RCP de qualité (compressions \SI{5}{cm}, \SIrange{100}{120}{/min}, relâchement complet) :}
\begin{center}
\begin{tabularx}{\linewidth}{|l|X|X|}\hline
\textbf{Paramètre hémodynamique} & \textbf{Valeur normale} & \textbf{Pendant RCP de qualité} \\ \hline
Pression artérielle systolique & \SI{120}{mmHg} & \SIrange{60}{80}{mmHg} \\ \hline
Pression artérielle diastolique & \SI{80}{mmHg} & \SIrange{20}{30}{mmHg} (cruciale pour perfusion coronaire) \\ \hline
Pression artérielle moyenne (PAM) & \SI{93}{mmHg} & \SIrange{30}{40}{mmHg} \\ \hline
Débit cardiaque (DC) & \SI{5}{L/min} & \SIrange{1.5}{2.0}{L/min} (\textbf{\SIrange{25}{30}{\%}}) \\ \hline
Débit cérébral & \SI{750}{ml/min} & \SIrange{200}{250}{ml/min} \\ \hline
Pression de perfusion coronaire (PPC = PAD - PVC) & \SI{60}{mmHg} & \SIrange{15}{25}{mmHg} (cible $>$\SI{15}{mmHg} pour ROSC) \\ \hline
\end{tabularx}
\end{center}
\par \textbf{Interprétation :} La compression thoracique augmente la pression intrathoracique (mécanisme de la « pompe thoracique ») et comprime directement le cœur entre sternum et colonne (mécanisme de la « pompe cardiaque »). La \textbf{pression artérielle diastolique} (PAD) est le déterminant principal de la \textbf{perfusion coronaire} (le cœur se vascularise en diastole). Une PAD $>$\SI{20}{mmHg} (donc PPC $>$\SI{15}{mmHg}) est corrélée au retour de la circulation spontanée (ROSC).
\legende{Données hémodynamiques moyennes lors d'une RCP manuelle de haute qualité. Le maintien d'une pression diastolique suffisante est la clé de la revascularisation du myocarde et du cerveau. PVC = Pression Veineuse Centrale.}
\end{experience}

\begin{popfigure}
\begin{tikzpicture}[scale=0.85, every node/.style={font=\small}]
    % Heart simplified (Right heart left, Left heart right)
    % Right Heart
    \draw[popdark, thick, fill=popblue!10] (1,3) .. controls (0,4) and (0,1) .. (1,0) .. controls (2,-0.5) and (3,1.5) .. (2,4) .. controls (1.5,4.5) and (1,4.5) .. (1,3);
    % Vena Cava / Right Atrium inlet
    \draw[popblue, thick, ->] (0.2,3.5) -- (0.8,3.2) node[midway, left, popblue] {Sang désoxygéné};
    % Lungs
    \draw[poppurple, thick, fill=poppurple!10] (3.5,4) ellipse (1.2 and 0.8);
    \draw[poppurple, thick, ->] (2.2,3.5) -- (3,3.8) node[midway, above] {Artère pulmonaire};
    \draw[popred, thick, ->] (4.2,3.8) -- (5,3.5) node[midway, above] {Veines pulmonaires};
    % Left Heart
    \draw[popred, thick, fill=popred!20] (5,3) .. controls (5.5,3.5) and (6.5,3.5) .. (7,3) .. controls (7.5,2.5) and (7.5,1.5) .. (7,1) .. controls (6.5,0.5) and (5.5,0.5) .. (5,1) .. controls (4.5,1.5) and (4.5,2.5) .. (5,3);
    % Aorta to brain
    \draw[popred, very thick, ->] (7,2.5) -- (9,2.5) -- (9,5) node[above, popdark] {Cerveau};
    \draw[popred, very thick] (9,2.5) -- (9,0) node[below, popdark] {Corps};
    % Coronaries
    \draw[poporange, thick, ->] (7,2.8) .. controls (7.5,2.5) and (7.5,3.5) .. (7,3.2) node[right, poporange] {Coronaires};
    % Crossed out arrow for Cardiac Arrest
    \draw[popred, very thick] (7,2.5) -- (9,2.5);
    \node[popred, rotate=45, scale=1.5, font=\bfseries] at (8,2.5) {\texttimes};
    \node[popred, align=center, font=\bfseries\small] at (10.5,2.5) {ARRÊT\\CARDIAQUE\\= Flux = 0};
    % CPR Arrow
    \draw[popblue, thick, <->] (3.5,0.5) -- (6.5,0.5) node[midway, below] {Compressions thoraciques (RCP)};
    \node[popblue, align=center, font=\small] at (5,-0.5) {Génère flux rétrograde\\$\approx$ 25-30\% normal};
\end{tikzpicture}
\legende{Schéma de la circulation normale (flèches rouges continues) et effet de l'arrêt cardiaque (flèche barrée). En bas, les compressions thoraciques (flèche bleue double) créent un flux sanguin rétrograde partiel vers le cerveau et les coronaires.}
\end{popfigure}

\courssub{Reconnaissance de l'arrêt cardiaque : les 3 questions qui sauvent}

La reconnaissance précoce est le maillon le plus fragile de la chaîne de survie. Au Cameroun comme ailleurs, le retard à démarrer la RCP s'explique souvent par l'hésitation du témoin : « Est-ce vraiment un arrêt cardiaque ? » La méthode internationale (ERC/ILCOR/AHA) et le programme MINESEC reposent sur une évaluation \textbf{rapide (max 10 s)}, \textbf{binaire} et \textbf{sans matériel}.

\begin{methode}[Reconnaître un arrêt cardiaque en 3 questions (Algorithme « Vérifier - Appeler - Masser »)]
\emph{Devant une victime inerte (au sol, immobile), réaliser dans l'ordre :}
\begin{enumerate}
    \item \textbf{La victime répond-elle ?} \\
    \emph{Action :} Secouer doucement les épaules, parler fort (« Madame, Monsieur, m'entendez-vous ? »). \\
    \emph{Réponse attendue :} \textbf{Non} $\rightarrow$ \textbf{Inconsciente}. (Si oui : surveiller, alerter si besoin, ne pas faire de RCP).
    \item \textbf{La victime respire-t-elle \emph{normalement} ?} \\
    \emph{Action :} \textbf{Libérer les voies aériennes} (basculer la tête en arrière, soulever le menton). Puis \textbf{Regarder - Écouter - Sentir} pendant \textbf{10 secondes maximum} : regarder le soulèvement thoracique, écouter le bruit respiratoire, sentir l'air sur la joue. \\
    \emph{Réponse attendue :} \textbf{Non} (absence totale de respiration \textbf{OU} \textbf{gasps} / respirations agoniques) $\rightarrow$ \textbf{Arrêt cardiaque}.
    \item \textbf{Conclusion :} \textbf{Inconsciente + Ne respire pas normalement = ARRÊT CARDIAQUE.} \\
    \emph{Action immédiate :} \textbf{Alerter (ou faire alerter) le \SI{112}{/119} (SAMU/Pompiers Cameroun) + Demander un DAE} $\rightarrow$ \textbf{Commencer la RCP immédiatement (30 compressions / 2 insufflations)}.
\end{enumerate}
\end{methode}

\begin{attention}[Piège mortel : les gasps (respirations agoniques) ne sont PAS une respiration normale !]
C'est l'erreur \textbf{la plus fréquente et la plus grave} chez le secouriste débutant (et parfois chez des professionnels non entraînés).
\begin{itemize}
    \item \textbf{Aspect :} Mouvements respiratoires \textbf{bruyants} (ronflement, râle, grognement), \textbf{irréguliers} (pauses longues), \textbf{peu efficaces} (volume courant très faible, pas de soulèvement thoracique franc). Cela ressemble à un « hoquet » ou un « dernier souffle » répété.
    \item \textbf{Origine :} Activité résiduelle du centre respiratoire bulbaire hypoxique. Signe de \textbf{très mauvais pronostic} mais \textbf{témoin d'un arrêt cardiaque récent} (donc RCP urgente et potentiellement efficace).
    \item \textbf{Règle d'or :} \textbf{En cas de doute sur la normalité de la respiration $\rightarrow$ Considérer que la victime NE RESPIRE PAS et DÉMARRER LA RCP.} Il vaut mieux masser un vivant qui respire mal (risque faible de fracture costale) que de ne pas masser un mort qui ne respire plus (mort certaine).
\end{itemize}
\end{attention}

\begin{popfigure}
\begin{tikzpicture}[node distance=1.2cm and 1.5cm, every node/.style={font=\small, align=center},
    decision/.style={diamond, draw=popdark, thick, fill=popgoldL, aspect=2},
    process/.style={rectangle, draw=popdark, thick, fill=popblueL, rounded corners, minimum width=2.5cm, minimum height=1cm},
    startstop/.style={ellipse, draw=popdark, thick, fill=popgreenL, minimum width=2cm},
    arrow/.style={-Stealth, thick, popdark},
    yn/.style={font=\footnotesize, pos=0.5, sloped, anchor=center}]
    
    \node[startstop, align=center] (start){Victime découverte\\ inerte / au sol};
    \node[process, below of=start, align=center] (protect){1. \textbf{PROTÉGER} :\\ Sécuriser les lieux};
    \node[decision, below of=protect, align=center] (resp){2. \textbf{RÉPOND ?}\\ (Secouer, parler fort)};
    \node[process, below left=1.5cm and 0.5cm of resp, align=center] (conscious){Victime CONSCIENTE\\ $\rightarrow$ Surveiller, Alerter si besoin,\\ Position d'attente};
    \node[process, below of=resp, yshift=-1cm, align=center] (airway){3. \textbf{LIBÉRER VOIES AÉRIENNES}\\ (Basculer tête, soulever menton)};
    \node[decision, below of=airway, align=center] (breath){4. \textbf{RESPIRE NORMALMENT ?}\\ (Regarder-Écouter-Sentir $\le$ 10 s)};
    \node[process, below left=1.5cm and 0.5cm of breath, align=center] (breathing){Victime RESPIRE\\ $\rightarrow$ \textbf{PLS} + Alerter + Surveiller};
    \node[process, below of=breath, yshift=-1cm, align=center] (alert){5. \textbf{ALERTER \SI{112}{/119}}\\ (Ou faire alerter + DAE)};
    \node[process, below of=alert, fill=popred!20, draw=popred, align=center] (cpr){6. \textbf{DÉMARRER RCP}\\ 30 Compressions / 2 Insufflations\\ \textbf{SANS INTERRUPTION}};
    
    \draw[arrow] (start) -- (protect);
    \draw[arrow] (protect) -- (resp);
    \draw[arrow] (resp) -- node[yn, left] {NON} (airway);
    \draw[arrow] (resp) -- node[yn, above] {OUI} (conscious);
    \draw[arrow] (airway) -- (breath);
    \draw[arrow] (breath) -- node[yn, left] {NON / GASPS} (alert);
    \draw[arrow] (breath) -- node[yn, above] {OUI} (breathing);
    \draw[arrow] (alert) -- (cpr);
\end{tikzpicture}
\legende{Organigramme de la conduite à tenir face à une victime inerte (Algorithme ERC/ILCOR 2021 adapté MINESEC). La décision binaire « Respire normalement ? » est le point critique : tout doute = RCP.}
\end{popfigure}

\courssub{Conduite immédiate : de la reconnaissance à l'action}

Une fois l'arrêt cardiaque reconnu (Inconscient + Respiration absente ou agonique), la conduite à tenir est immuable et doit devenir un réflexe conditionné.

\begin{exoresolu}[Reconnaissance d'arrêt cardiaque : cas pratiques]
\textbf{Énoncé :} Pour chacune des situations suivantes, dites si vous démarrez la RCP (OUI/NON) et justifiez en une phrase.
\begin{enumerate}
    \item Un homme de 55 ans s'effondre au marché de Mfoundi. Il ne bouge pas, ne répond pas. Son thorax ne se soulève pas. Vous ne sentez pas d'air sur sa joue.
    \item Une femme de 30 ans, électrocutée par un câble sectionné, est au sol. Elle ne répond pas. Elle émet un bruit de ronflement rauque toutes les 5-6 secondes, la poitrine se soulève très peu.
    \item Un enfant de 8 ans retrouvé au fond d'une piscine. Il est bleu, ne répond pas, ne respire pas.
    \item Un homme de 70 ans, connu insuffisant cardiaque, assis sur un banc, dort la bouche ouverte. Il ronfle régulièrement. Il se réveille quand vous lui parlez.
\end{enumerate}
\tcblower
\textbf{Solution détaillée :}
\begin{enumerate}
    \item \textbf{OUI.} Inconscient (ne répond pas) + Apnée (absence totale de respiration) = Arrêt cardiaque. $\rightarrow$ Alerter \SI{112}{/119} + RCP immédiate.
    \item \textbf{OUI.} Inconsciente + \textbf{Gasps} (respiration agonique : bruit rauque, irrégulier, amplitude faible) = Arrêt cardiaque. $\rightarrow$ Alerter + RCP immédiate. \emph{Ne pas confondre avec une respiration normale.}
    \item \textbf{OUI.} Inconscient + Apnée (noyade = arrêt cardiaque par hypoxie primaire). $\rightarrow$ \textbf{5 insufflations initiales} (spécificité noyade) puis 30/2. Alerter.
    \item \textbf{NON.} La victime \textbf{répond} (se réveille à la voix) $\rightarrow$ Elle est \textbf{consciente}. Pas d'arrêt cardiaque. Surveiller, installer confortablement, alerter si symptômes persistants (dyspnée, douleur thoracique).
\end{enumerate}
\end{exoresolu}

\begin{aretenir}[Synthèse : Les 3 piliers de la reconnaissance]
\begin{enumerate}
    \item \textbf{Physiologie :} Cerveau = « organe roi » de l'oxygène. \SI{10}{s} sans flux = inconscience. \SI{3}{-5}{min} sans flux = mort neuronale. RCP = \SIrange{25}{30}{\%} débit = gain de temps cérébral.
    \item \textbf{Mécanisme :} FV (cœur « sac de vers ») = cause n°1. Asystolie/AESP = fin d'évolution ou causes non cardiaques (hypoxie, hypovolémie, toxiques).
    \item \textbf{Reconnaissance (Algorithme) :} \textbf{Ne répond pas} + \textbf{Ne respire pas normalement (apnée OU gasps)} = \textbf{ARRÊT CARDIAQUE} $\rightarrow$ \textbf{Alerter \SI{112}{/119} + RCP IMMÉDIATE}.
\end{enumerate}
\end{aretenir}

\begin{exercice}[Entraînement à la décision]
Vous êtes témoin des scènes suivantes. Pour chacune, cochez la bonne action : (A) Surveiller / PLS, (B) Alerter + RCP, (C) Alerter seulement.
\begin{enumerate}
    \item Joggeur au sol, ne bouge pas, ne répond pas, thorax immobile, pas de bruit respiratoire.
    \item Personne âgée dans un fauteuil, ronfle fort et régulièrement, ouvre les yeux quand on lui prend la main.
    \item Ouvrier tombé d'un échafaudage, inerte, respiration bruyante, saccadée, toutes les 10 secondes.
    \item Nourrisson de 6 mois dans son berceau, pâle, ne réagit pas, ventre immobile.
\end{enumerate}
\emph{(Corrigé en fin de chapitre / section 6).}
\end{exercice}

La maîtrise de cette reconnaissance — simple sur le papier, difficile sous le stress — est la compétence fondatrice du secouriste. La section suivante détaillera pas à pas la technique de la \textbf{Réanimation Cardio-Pulmonaire (RCP)} : compressions thoraciques de haute qualité, insufflations, utilisation du DAE et gestion des relais.

\coursec{Les gestes qui sauvent en cas d'arrêt cardiaque : la RCP}

Dans la section précédente, nous avons appris à \textbf{reconnaître} l'arrêt cardiaque : une victime inconsciente, qui ne respire pas ou présente une respiration agonique (bruyante, irrégulière, inefficace). Ce diagnostic clinique, posé en moins de 10 secondes, déclenche immédiatement l'action. Car en cas d'arrêt cardiaque, \textbf{chaque minute de retard à commencer la réanimation cardio-pulmonaire (RCP) réduit les chances de survie d'environ 10 \%}. La RCP n'est pas un geste médical réservé aux professionnels : c'est un acte citoyen, à la portée de tout témoin formé, qui maintient une circulation artificielle minimale en attendant le défibrillateur et les secours organisés. Comprenons \emph{pourquoi} et \emph{comment} ces gestes sauvent des vies.

\courssub{La chaîne de survie et l'alerte immédiate : le premier maillon}

L'efficacité de la prise en charge repose sur la \textbf{chaîne de survie}, dont les quatre maillons sont indissociables :
\begin{enumerate}
    \item \textbf{Alerte précoce} et reconnaissance de l'arrêt cardiaque.
    \item \textbf{RCP précoce} par un témoin (massage cardiaque, avec ou sans ventilation).
    \item \textbf{Défibrillation précoce} par DAE (Défibrillateur Automatisé Externe).
    \item \textbf{Prise en charge médicale avancée} par les secours (SAMU, pompiers) puis hospitalière.
\end{enumerate}
Le témoin est l'acteur des deux premiers maillons. \textbf{Alerter ou faire alerter} est la première action vitale, dès la constatation de l'inconscience et de l'absence de respiration normale.

\begin{methode}[L'alerte aux secours : le message structuré]
\begin{enumerate}
    \item \textbf{Composer le numéro d'urgence} : \textbf{119} (Sapeurs-pompiers au Cameroun) ou \textbf{112} (numéro d'urgence européen, fonctionne sur GSM même sans crédit). Si un SAMU structuré existe localement (Yaoundé, Douala), son numéro direct peut être utilisé.
    \item \textbf{Ne pas raccrocher avant l'opérateur}. Lui transmettre calmement le \textbf{message d'alerte} structuré :
    \begin{itemize}
        \item \textbf{Où ?} Adresse précise, repères visibles (quartier, carrefour, étage), coordonnées GPS si possible.
        \item \textbf{Qui ?} Votre nom, votre numéro de téléphone (pour rappel en cas de coupure).
        \item \textbf{Quoi ?} Nature de l'urgence : « Arrêt cardiaque, victime inconsciente, ne respire pas ».
        \item \textbf{Combien ?} Nombre de victimes (ici 1).
        \item \textbf{État ?} Préciser : « RCP en cours par un témoin » ou « En attente de consignes ».
    \end{itemize}
    \item \textbf{Exécuter les consignes} de l'opérateur (guidage pour la RCP, localisation du DAE le plus proche).
    \item \textbf{Désigner un témoin} pour aller chercher le DAE le plus proche et guider les secours à leur arrivée.
\end{enumerate}
\end{methode}

\begin{attention}[Piège fréquent : « J'ai appelé, j'attends »]
Ne \textbf{jamais} attendre les secours pour commencer le massage cardiaque. L'alerte et le massage sont \textbf{simultanés} : si vous êtes seul, alertez en mains libres (haut-parleur) \emph{pendant} que vous commencez les compressions. Si vous êtes deux, l'un alerte et cherche le DAE, l'autre commence immédiatement le massage. \textbf{Zéro minute sans compression thoracique.}
\end{attention}

\courssub{Le massage cardiaque externe (MCE) : principes physiologiques et technique}

\emph{Pourquoi comprimer le thorax ?} Le cœur est logé dans le médiastin, entre le sternum (en avant) et la colonne vertébrale (en arrière). Lors d'une compression thoracique efficace, l'augmentation de la pression intrathoracique et la compression directe des ventricules entre le sternum et la colonne vertébrale expulsent le sang vers l'aorte et l'artère pulmonaire (mécanismes dits de la \textbf{pompe cardiaque} et de la \textbf{pompe thoracique}). Lors du relâchement complet, la décompression thoracique crée une dépression qui favorise le retour veineux et le remplissage cardiaque. \textbf{Sans relâchement complet, pas de remplissage, donc pas de débit au cycle suivant.}

\begin{popfigure}
\begin{tikzpicture}[scale=1.0, every node/.style={font=\small}]
    % Colonne vertébrale (arrière)
    \draw[fill=popmuted!30, draw=popdark, thick] (0,0) rectangle (1,5.5);
    \node[rotate=90] at (0.5, 2.75) {Colonne vertébrale};
    % Cage thoracique (parois latérales)
    \draw[fill=popblueL!40, draw=popblue, thick] (1,0) .. controls (1.5,0.5) and (1.5,5) .. (1,5.5);
    \draw[fill=popblueL!40, draw=popblue, thick] (6,0) .. controls (5.5,0.5) and (5.5,5) .. (6,5.5);
    % Sternum (avant)
    \draw[fill=poporangeL!60, draw=poporange, thick] (1,5.0) rectangle (6,5.5);
    \node[poporange] at (3.5, 5.25) {Sternum};
    % Cœur (entre sternum et colonne)
    \draw[fill=popink!35, draw=popink, thick] (1.6, 1.5) .. controls (1.6, 3.6) and (4.4, 4.0) .. (5.0, 2.6) .. controls (5.4, 1.5) and (3.4, 1.0) .. (1.6, 1.5);
    \node[popink, font=\bfseries] at (3.3, 2.5) {C\oe ur};
    % Aorte et artère pulmonaire (schématique)
    \draw[popink, thick, ->] (3.3, 3.8) -- (3.3, 4.6) node[above, font=\footnotesize] {Aorte};
    \draw[popblue, thick, ->] (3.9, 3.5) -- (4.6, 4.3) node[right, font=\footnotesize] {Artère pulmonaire};
    % Mains (position)
    \draw[fill=poppinkL!60, draw=poppink, thick] (2.5, 5.5) rectangle (4.5, 6.2);
    \draw[fill=poppinkL!60, draw=poppink, thick] (2.7, 6.2) rectangle (4.3, 6.9);
    \node[poppink, font=\bfseries\footnotesize] at (3.5, 6.55) {Talons des mains};
    % Flèche de compression
    \draw[popdark, very thick, ->] (3.5, 5.5) -- (3.5, 4.1) node[midway, right, font=\footnotesize] {Compression 5--6 cm};
    \node[popdark, align=center, font=\footnotesize] at (3.5, 7.3) {Force verticale, bras tendus};
    \draw[popdark, thick, ->] (3.5, 7.1) -- (3.5, 6.95);
    % Légende repères
    \node[anchor=west, font=\footnotesize] at (6.4, 4.5) {1. Sternum (point d'appui)};
    \node[anchor=west, font=\footnotesize] at (6.4, 4.0) {2. C\oe ur (ventricules comprimés)};
    \node[anchor=west, font=\footnotesize] at (6.4, 3.5) {3. Colonne vertébrale (plan dur)};
    \node[anchor=west, font=\footnotesize] at (6.4, 3.0) {4. Talon de la main 1};
    \node[anchor=west, font=\footnotesize] at (6.4, 2.5) {5. Main 2 par-dessus, doigts écartés};
\end{tikzpicture}
\legende{Coupe schématique du thorax au niveau du cœur (vue antéro-postérieure). La compression verticale du sternum écrase les ventricules contre la colonne vertébrale, éjectant le sang vers l'aorte et l'artère pulmonaire. Le point d'appui est le centre de la poitrine (moitié inférieure du sternum).}
\end{popfigure}

\begin{definition}[Technique du Massage Cardiaque Externe (MCE) chez l'adulte]
\begin{enumerate}
    \item \textbf{Position de la victime} : sur le dos (décubitus dorsal), sur un \textbf{plan dur} (sol, planche). Un lit mou annule l'efficacité des compressions.
    \item \textbf{Position du secouriste} : à genoux, verticalement au-dessus de la poitrine de la victime, épaules alignées au-dessus des mains, bras \textbf{tendus} (coudes verrouillés) pour utiliser le poids du corps.
    \item \textbf{Point d'appui} : le \textbf{talon d'une main} placé au \textbf{centre de la poitrine} (moitié inférieure du sternum). L'autre main par-dessus, doigts entrelacés ou écartés, \textbf{ne touchant pas le thorax} (seul le talon compresse).
    \item \textbf{Paramètres techniques (cibles de qualité)} :
    \begin{itemize}
        \item \textbf{Fréquence} : \textbf{100 à 120 compressions par minute}.
        \item \textbf{Profondeur} : \textbf{5 à 6 cm} chez l'adulte (environ un tiers de l'épaisseur antéro-postérieure du thorax).
        \item \textbf{Relâchement} : \textbf{complet} à chaque cycle (la poitrine reprend sa forme initiale), sans décoller les mains.
        \item \textbf{Durée du cycle} : temps de compression égal au temps de relâchement.
    \end{itemize}
\end{enumerate}
\end{definition}

\begin{exemplebox}[Repère mnémotechnique : le rythme « Stayin' Alive »]
La fréquence cible de 100--120 compressions/min correspond au tempo de la célèbre chanson des \textbf{Bee Gees « Stayin' Alive »} (environ 103 battements/min). Dans les formations aux premiers secours, on utilise ce rythme musical pour guider le secouriste. L'essentiel : \textbf{un rythme soutenu, régulier, sans ralentir}.
\end{exemplebox}

\begin{attention}[Erreur fréquente : « J'ai peur de casser les côtes »]
C'est la \textbf{première réticence} des témoins. \textbf{Vérité scientifique} : des craquements costaux ou sternaux peuvent survenir (fractures de côtes ou du sternum) chez l'adulte, surtout chez le sujet âgé ou ostéoporotique. \textbf{Cependant} :
\begin{itemize}
    \item Une compression \emph{trop faible} (profondeur inférieure à 5 cm) ne génère \textbf{aucun débit sanguin cérébral efficace} : la victime décède avec des côtes intactes.
    \item Une compression \emph{efficace} (5--6 cm) maintient une perfusion cérébrale et coronaire minimale ; d'éventuelles fractures costales sont une \textbf{complication acceptable et secondaire} face à l'enjeu vital.
    \item \textbf{Règle d'or} : « Mieux vaut des côtes cassées qu'une vie brisée. » On compresse \textbf{fort et vite}, sans hésitation.
\end{itemize}
\end{attention}

\courssub{La ventilation artificielle et le ratio 30/2 : l'apport d'oxygène}

Le massage seul assure une circulation, mais le sang qui circule doit être oxygéné. Chez l'adulte victime d'un arrêt cardiaque \textbf{soudain} (origine cardiaque), le sang artériel est initialement bien saturé en O$_2$ ; pendant les premières minutes, les compressions seules (RCP « mains seules ») suffisent souvent. Mais au-delà de 4--5 min, ou si l'arrêt est \textbf{secondaire à une cause hypoxique} (noyade, étouffement, asthme, intoxication, enfant ou nourrisson), l'apport d'oxygène par ventilation devient \textbf{critique}.

\begin{methode}[Ventilation bouche-à-bouche : technique]
\begin{enumerate}
    \item Après 30 compressions, \textbf{libérer les voies aériennes} : bascule prudente de la tête en arrière (une main sur le front) + élévation du menton (deux doigts).
    \item \textbf{Pincer le nez} de la victime (pouce et index).
    \item \textbf{Inspirer normalement} (pas à fond), poser votre bouche \textbf{hermétiquement} sur la sienne (ou utiliser un masque de poche ou une protection si disponible).
    \item \textbf{Souffler doucement} pendant \textbf{1 seconde environ}, en surveillant le \textbf{soulèvement de la poitrine} (signe d'efficacité).
    \item \textbf{Relâcher}, laisser la poitrine retomber (expiration passive).
    \item \textbf{Deuxième insufflation} (même procédure).
    \item \textbf{Reprendre immédiatement les compressions} (viser moins de 10 s d'interruption pour les 2 insufflations).
\end{enumerate}
\textbf{Ratio universel adulte (secouriste unique) : 30 compressions / 2 insufflations} (30:2).
\end{methode}

\begin{popfigure}
\begin{tikzpicture}[scale=0.95, every node/.style={font=\small}]
    % Ligne de temps
    \draw[popdark, thick, ->] (0,0) -- (13.2,0) node[right] {Temps};
    % Bloc compressions
    \draw[fill=popink!30, draw=popink, thick] (0.2,0) rectangle (8.6,1.0);
    \node[popink, align=center] at (4.4,0.5) {\textbf{30 compressions} (environ 15--18 s à 100--120/min)};
    % Bloc insufflations
    \draw[fill=popblueL!60, draw=popblue, thick] (8.8,0) rectangle (9.9,1.0);
    \node[popblue, align=center, font=\footnotesize] at (9.35,0.5) {Insuffl. 1};
    \draw[fill=popblueL!60, draw=popblue, thick] (10.1,0) rectangle (11.2,1.0);
    \node[popblue, align=center, font=\footnotesize] at (10.65,0.5) {Insuffl. 2};
    \node[popblue, align=center, font=\footnotesize] at (10.0,1.35) {2 insufflations : moins de 10 s au total};
    % Flèche retour
    \draw[popdark, thick, ->] (11.4,0.5) .. controls (12.4,1.8) and (1.0,1.8) .. (0.4,1.1);
    \node[popdark, align=center, font=\footnotesize] at (6.4,2.0) {Reprise immédiate des compressions : nouveau cycle 30:2};
    % Légende couleurs
    \node[anchor=west, font=\footnotesize] at (0.2,-0.9) {\textcolor{popink}{\rule{5mm}{2mm}} Compression thoracique (pousser fort et vite)};
    \node[anchor=west, font=\footnotesize] at (7.0,-0.9) {\textcolor{popblue}{\rule{5mm}{2mm}} Insufflation (souffler 1 s, voir la poitrine monter)};
\end{tikzpicture}
\legende{Frise chronologique d'un cycle de RCP standardisé (30:2). Les 30 compressions durent environ 15 à 18 secondes (à 100--120/min). Les 2 insufflations ne doivent pas excéder 10 secondes au total. La régularité et la minimisation des pauses sont les garants de l'efficacité hémodynamique.}
\end{popfigure}

\courssub{La RCP « mains seules » (Compression-Only CPR) : une alternative validée}

De nombreuses études épidémiologiques (registres japonais et suédois notamment) ont montré que chez l'adulte victime d'un \textbf{arrêt cardiaque soudain devant témoin} (origine cardiaque présumée), la RCP par \textbf{compressions thoraciques continues sans ventilation} (« hands-only CPR ») offre une survie \textbf{au moins équivalente} à la RCP conventionnelle (30:2) réalisée par des témoins non professionnels.

\begin{savaistu}[Pourquoi la RCP « mains seules » fonctionne-t-elle ?]
\begin{enumerate}
    \item \textbf{Réserve en O$_2$} : au moment de l'arrêt soudain, le sang artériel et les poumons contiennent encore une réserve d'oxygène suffisante pour les premières minutes de circulation artificielle.
    \item \textbf{Éviter les interruptions} : la ventilation bouche-à-bouche par un non-expert entraîne souvent des pauses longues (plus de 10 s), mal réalisées, qui font chuter la pression de perfusion coronaire.
    \item \textbf{Levée de la réticence} : la peur du contact buccal (maladies transmissibles, vomissements) paralyse de nombreux témoins. La RCP « mains seules » lève cet obstacle : \textbf{« Alerter, Masser, Ne pas arrêter »} est un message simple qui augmente le taux de RCP par les témoins.
    \item \textbf{Exceptions où la ventilation reste indispensable} : arrêt \textbf{non vu}, \textbf{enfant ou nourrisson}, \textbf{noyade}, \textbf{asphyxie}, \textbf{intoxication} : dans ces cas, l'arrêt est d'origine hypoxique et le sang est déjà pauvre en O$_2$ ; il faut ventiler (ratio 30:2).
\end{enumerate}
\end{savaistu}

\courssub{Le Défibrillateur Automatisé Externe (DAE) : le maillon technologique}

Le massage cardiaque maintient une perfusion, mais \textbf{ne défibrille pas}. La cause majeure de l'arrêt cardiaque soudain de l'adulte est la \textbf{fibrillation ventriculaire (FV)} : activité électrique chaotique et inefficace des ventricules, qui ne produit aucun débit sanguin. Seul un choc électrique externe (défibrillation) peut « réinitialiser » simultanément les cellules myocardiques et permettre le retour d'un rythme organisé. Le DAE rend ce geste accessible à tout citoyen.

\begin{definition}[Défibrillateur Automatisé Externe (DAE)]
Appareil portable, alimenté par batterie, qui \textbf{analyse automatiquement le rythme cardiaque} via des électrodes adhésives posées sur le thorax nu de la victime. Si l'algorithme détecte un rythme \textbf{chocable} (fibrillation ventriculaire ou tachycardie ventriculaire sans pouls), il \textbf{charge le condensateur} et \textbf{ordonne (messages vocaux et lumineux) de délivrer le choc} (en appuyant sur un bouton clignotant pour le DAE semi-automatique, ou automatiquement pour le DAE entièrement automatique). Si le rythme n'est pas chocable (asystolie, activité électrique sans pouls), il n'autorise pas le choc et invite à reprendre la RCP.
\end{definition}

\begin{methode}[Utilisation du DAE : procédure « Allumer -- Coller -- Laisser faire »]
\begin{enumerate}
    \item \textbf{Allumer} le DAE (bouton marche ou ouverture du capot) \textbf{dès qu'il arrive} : il guide l'utilisateur par des messages vocaux clairs.
    \item \textbf{Dénuder le thorax} (couper les vêtements, essuyer la peau si elle est mouillée ou transpirante, raser si très poilue -- kit fourni avec l'appareil).
    \item \textbf{Coller les électrodes} selon le schéma dessiné dessus :
    \begin{itemize}
        \item une sous la clavicule droite (bord droit du sternum) ;
        \item une sur le flanc gauche (au niveau de la pointe du cœur, sous l'aisselle).
    \end{itemize}
    \item \textbf{Ne plus toucher la victime} : le DAE analyse (« Ne touchez pas la victime »).
    \item \textbf{Si « Choc recommandé »} :
    \begin{itemize}
        \item vérifier que \textbf{personne ne touche la victime} (crier « Éloignez-vous ! ») ;
        \item appuyer sur le bouton \textbf{CHOC} (clignotant) si l'appareil est semi-automatique ; ne rien faire s'il est entièrement automatique.
    \end{itemize}
    \item \textbf{Immédiatement après le choc (ou si « Pas de choc conseillé ») : reprendre la RCP} (30:2) pendant 2 minutes, sans retirer les électrodes. Le DAE réanalysera automatiquement le rythme toutes les 2 minutes.
\end{enumerate}
\end{methode}

\begin{attention}[Pièges DAE au Bac et en pratique]
\begin{itemize}
    \item \textbf{Ne pas retarder le massage} pour installer le DAE : un secouriste continue le MCE pendant qu'un autre pose les électrodes.
    \item \textbf{Médicaments transdermiques} (patchs de nicotine, de nitroglycérine) : les retirer (avec des gants) et essuyer la peau avant de coller l'électrode à cet endroit (risque de brûlure ou d'étincelle).
    \item \textbf{Dispositifs implantés} (pacemaker, défibrillateur interne) : boîtier visible sous la peau (environ 5 cm). \textbf{Ne pas placer l'électrode dessus} (la décaler de quelques centimètres).
    \item \textbf{Environnement humide ou métallique} : sortir la victime de l'eau, la sécher ; éviter le contact direct avec une surface métallique conductrice.
    \item \textbf{Enfant de moins de 8 ans (moins de 25 kg)} : utiliser si possible des \textbf{électrodes pédiatriques} (énergie atténuée). À défaut, les électrodes adultes sont acceptées (une sur le devant du thorax, une dans le dos si le thorax est trop petit pour que les électrodes ne se touchent pas).
\end{itemize}
\end{attention}

\courssub{Conduite à tenir globale : la boucle « Alerter -- Masser -- Défibriller »}

La synthèse opérationnelle pour le citoyen-secouriste face à un adulte en arrêt cardiaque :

\begin{experience}[Scénario type : arrêt cardiaque soudain devant témoin -- conduite à tenir pas à pas]
\textbf{Matériel} : mannequin de RCP adulte, DAE d'entraînement, téléphone (simulation d'appel au 119), chronomètre.

\textbf{Protocole} :
\begin{enumerate}
    \item \textbf{Sécuriser} la zone (circulation, incendie, risque électrique) : protéger la victime, les témoins et soi-même.
    \item \textbf{Évaluer la conscience} : « Monsieur/Madame, m'entendez-vous ? » + tapotement des épaules. \textit{Résultat : inconscient.}
    \item \textbf{Libérer les voies aériennes} (bascule de la tête, élévation du menton) + \textbf{regarder, écouter, sentir} la respiration (10 s maximum). \textit{Résultat : absente ou agonique.}
    \item \textbf{Alerter ou faire alerter} (119/112) + \textbf{envoyer chercher le DAE} (si disponible). \textit{Top départ du chronomètre.}
    \item \textbf{Commencer le MCE} : 30 compressions (centre de la poitrine, 5--6 cm, 100--120/min, relâchement total).
    \item \textbf{Si formé et consentant} : 2 insufflations efficaces (la poitrine se soulève). \textbf{Sinon} : compressions continues sans pause.
    \item \textbf{Arrivée du DAE} : allumer, coller les électrodes (sans arrêter le MCE si possible), laisser analyser.
    \item \textbf{Choc si indiqué} : éloigner tout le monde, délivrer le choc, \textbf{reprendre immédiatement la RCP 30:2 pendant 2 min}.
    \item \textbf{Boucle} : RCP 2 min $\rightarrow$ analyse DAE $\rightarrow$ choc ou non $\rightarrow$ RCP... \textbf{jusqu'à} :
    \begin{itemize}
        \item l'arrivée des secours médicalisés (ils prennent le relais) ;
        \item la reprise d'une \textbf{respiration normale} (régulière, efficace) $\rightarrow$ \textbf{PLS} (Position Latérale de Sécurité, voir section 5) ;
        \item l'épuisement physique du ou des secouristes (relais toutes les 2 min si possible).
    \end{itemize}
\end{enumerate}
\textbf{Interprétation} : la qualité de la RCP (profondeur, fréquence, relâchement complet, interruptions minimales) et la précocité du choc (objectif : moins de 3--5 min) sont les deux déterminants pronostiques majeurs sur lesquels le témoin peut agir.
\end{experience}

\begin{aretenir}[Synthèse des critères de qualité de la RCP]
\begin{center}
\begin{tabularx}{\linewidth}{|l|X|X|}\hline
\rowcolor{popblueL} \textbf{Paramètre} & \textbf{Cible} & \textbf{Erreur fréquente à corriger} \\ \hline
Fréquence & 100 -- 120 / min & Trop lent (moins de 100) ou trop rapide (plus de 120 : remplissage incomplet) \\ \hline
Profondeur (adulte) & 5 -- 6 cm & Trop superficiel (moins de 5 cm) : perfusion inefficace \\ \hline
Relâchement & Complet (la poitrine remonte) & Appui résiduel entre deux compressions : remplissage cardiaque incomplet \\ \hline
Interruptions & Les plus courtes possibles (moins de 10 s) & Pauses longues pour ventiler, changer de secouriste ou installer le DAE \\ \hline
Position des mains & Centre de la poitrine (sternum) & Trop bas (appendice xiphoïde : risque de lésion du foie), trop haut, ou sur les côtes \\ \hline
\end{tabularx}
\end{center}
\end{aretenir}

\begin{exoresolu}[Cas pratique : gestion d'un arrêt cardiaque au marché]
\textbf{Énoncé} : Monsieur K., 58 ans, commerçant au Marché Central de Douala, s'effondre brutalement devant son étal. Vous êtes témoin direct. Il ne répond pas et ne respire pas normalement (râles irréguliers). Un DAE est installé à l'entrée du marché (200 m). Vous êtes seul pendant les 2 premières minutes, puis un collègue arrive. Détaillez votre conduite à tenir minute par minute.
\tcblower
\textbf{Solution détaillée} :
\begin{enumerate}
    \item \textbf{T = 0 (immédiat)} : sécuriser (écarter la foule, vérifier l'absence de danger). Évaluer la conscience : « Monsieur ! » + tapotement des épaules : inconscient. Ouvrir les voies aériennes (front-menton). Regarder-écouter-sentir 10 s : respiration agonique = \textbf{arrêt cardiaque}.
    \item \textbf{T = 0 min 30 s} : \textbf{alerter} : téléphone en mains libres, composer le \textbf{119} (pompiers de Douala). Message : « Arrêt cardiaque, homme d'environ 60 ans, Marché Central, entrée Nord, étal n°45, RCP en cours ». \textbf{Ne pas raccrocher} et suivre les consignes.
    \item \textbf{T = 1 min} : \textbf{démarrer le MCE} : victime au sol (plan dur). À genoux à côté d'elle. Talon d'une main au centre de la poitrine, bras tendus. \textbf{30 compressions} (compter à voix haute : « 1, 2, 3... 30 », environ 15--18 s), sur le rythme mental de « Stayin' Alive ».
    \item \textbf{T = 1 min 20 s} : \textbf{le collègue arrive}. Consignes claires : « \textbf{Va chercher le DAE à l'entrée Nord, cours, reviens vite !} » et « \textbf{Guide les pompiers jusqu'ici !} ».
    \item \textbf{T = 1 min 25 s} : \textbf{si formé} : 2 insufflations (bascule de la tête, pincer le nez, souffler 1 s, voir la poitrine monter). \textbf{Si non formé ou non consentant} : \textbf{continuer les compressions seules} sans pause.
    \item \textbf{T = 1 min 45 s} : reprise des cycles 30:2 (ou compressions continues). \textbf{Minimiser les interruptions.}
    \item \textbf{T = 3 à 4 min} : \textbf{le collègue revient avec le DAE}. \textbf{Ne pas arrêter le MCE} pendant qu'il allume l'appareil et colle les électrodes (thorax nu et sec). Ne s'écarter qu'à l'annonce « Ne touchez pas la victime ».
    \item \textbf{T = 4 min 10 s} : analyse du DAE. \textbf{Si « Choc recommandé »} : vérifier que personne ne touche la victime (crier « Éloignez-vous ! »), appuyer sur le bouton CHOC. \textbf{Immédiatement après} : reprendre la RCP 30:2 (se relayer avec le collègue toutes les 2 min pour éviter la fatigue, qui dégrade la qualité des compressions).
    \item \textbf{T = 6 min 10 s} : deuxième analyse automatique du DAE. Poursuivre la boucle selon le rythme annoncé.
    \item \textbf{Arrivée des pompiers/SAMU} : transmission orale structurée : « Monsieur K., 58 ans, arrêt cardiaque vu il y a environ 8 min, RCP par témoins continue, 2 chocs délivrés pour fibrillation ventriculaire, pas de reprise de la respiration ». \textbf{Ne pas arrêter} avant la prise en charge effective par les secours.
\end{enumerate}
\textbf{Points clés évaluables} : alerte précoce + MCE immédiat + envoi d'un relais chercher le DAE + qualité des compressions (profondeur, fréquence, relâchement) + interruptions inférieures à 10 s + transmission claire aux secours.
\end{exoresolu}

\begin{exemplebox}[Attitude empathique et communication : le « savoir-être » du secouriste]
Au-delà de la technique, le programme exige d'\textbf{assister la victime et de faire preuve d'empathie}.
\begin{itemize}
    \item \textbf{Parler à la victime} : même inconsciente, elle peut percevoir des sons. « Je suis là, les secours arrivent, je te masse le cœur, tiens bon. »
    \item \textbf{Gérer l'entourage} : désigner des rôles précis (« Toi, appelle le 119 ; toi, va chercher le DAE ; vous autres, écartez la foule et faites de la place »). Cela canalise la panique et fait de chacun un acteur utile.
    \item \textbf{Respecter la dignité} : couvrir la victime (surtout si ses vêtements ont été coupés pour poser les électrodes), écarter les regards indiscrets, ne pas filmer ni photographier (atteinte à la dignité et à la vie privée).
    \item \textbf{Après l'événement} : en cas de décès, ne pas cacher la vérité aux proches et les orienter vers un soutien psychologique ; en cas de survie, transmettre les informations aux secours et rester disponible pour témoigner.
\end{itemize}
\end{exemplebox}

\courssub{Exercices d'entraînement}

\begin{exercice}[QCM -- Paramètres de la RCP adulte]
Parmi les propositions suivantes, laquelle correspond aux recommandations internationales en vigueur pour la RCP adulte par un secouriste unique formé ?
\begin{enumerate}
    \item Fréquence 80--100/min, profondeur 4--5 cm, ratio 15:2.
    \item Fréquence 100--120/min, profondeur 5--6 cm, ratio 30:2.
    \item Fréquence 120--140/min, profondeur 6--7 cm, ratio 30:2.
    \item Fréquence 100--120/min, profondeur 3--4 cm, compressions seules continues.
\end{enumerate}
\end{exercice}

\begin{exercice}[Cas concret -- DAE et particularités]
Une victime de 65 ans est en arrêt cardiaque. Vous posez les électrodes du DAE. Vous remarquez : (1) un patch de nitroglycérine sur le côté droit du thorax, (2) un boîtier de pacemaker visible sous la clavicule gauche, (3) un thorax très poilu. Quelle est la conduite à tenir correcte pour chaque point ?
\end{exercice}

\begin{corrige}[Exercice 1]
\textbf{Bonne réponse : 2.} Fréquence 100--120/min, profondeur 5--6 cm, ratio 30:2. \textbf{Justification} : ce sont les valeurs des recommandations internationales (ERC/ILCOR) en vigueur. Une fréquence de 80--100/min est trop lente (anciennes recommandations) ; au-delà de 120/min, le remplissage diastolique des ventricules devient insuffisant ; une profondeur de 3--4 cm est inefficace (perfusion coronaire quasi nulle) ; le ratio 15:2 est réservé à l'enfant lorsque deux secouristes formés sont présents.
\end{corrige}

\begin{corrige}[Exercice 2]
\begin{enumerate}
    \item \textbf{Patch de nitroglycérine} : retirer le patch (avec des gants si possible), \textbf{essuyer soigneusement la peau} pour éliminer tout résidu de médicament (risque de brûlure ou d'étincelle sous l'électrode lors du choc), puis coller l'électrode.
    \item \textbf{Pacemaker (boîtier sous la clavicule gauche)} : \textbf{ne pas placer l'électrode gauche directement sur le boîtier} ; la décaler de quelques centimètres (plus bas et plus latéralement, vers le flanc). Le choc reste possible et nécessaire.
    \item \textbf{Thorax très poilu} : si les poils empêchent l'adhérence des électrodes, \textbf{raser rapidement} la zone (rasoir jetable fourni dans le kit du DAE) ou appuyer fermement les électrodes. Ne pas retarder le choc pour une épilation parfaite si l'adhérence est suffisante.
\end{enumerate}
\end{corrige}

\coursec{L'étouffement total : mécanisme et gestes selon la victime}

Après avoir maîtrisé la réanimation cardio-pulmonaire (RCP) face à un arrêt cardiaque, nous abordons maintenant une urgence vitale fréquente, souvent soudaine et dramatiquement évitable : \textbf{l'étouffement total par corps étranger}. Au Cameroun, comme partout, les repas de fête, les collations d'arachides ou de maïs grillé chez l'enfant, ou la mastication difficile chez la personne âgée sont autant de situations à risque. Comprendre le mécanisme physiopathologique permet de choisir le geste adapté à la morphologie de la victime, car \emph{un geste inadapté peut aggraver l'obstruction ou provoquer des lésions viscérales graves}.

\courssub{Mécanisme physiopathologique : l'obstruction des voies aériennes supérieures}

\paragraph{Rappel du trajet de l'air (classe de Première).} L'air inspiré traverse les fosses nasales (ou la bouche) $\to$ le pharynx (carrefour aéro-digestif) $\to$ le larynx (gardien des voies aériennes, siège des cordes vocales) $\to$ la trachée $\to$ les bronches $\to$ les alvéoles pulmonaires. L'épiglotte, clapet cartilagineux, bascule lors de la déglutition pour fermer l'entrée du larynx et diriger le bol alimentaire vers l'œsophage, situé en arrière de la trachée.

\paragraph{L'accident.} L'étouffement total survient lorsqu'un corps étranger (morceau de viande, arachide, arête de poisson, petit jouet, dentier) franchit l'épiglotte et se coince \textbf{au niveau du larynx} (seuil glottique) ou de la \textbf{trachée proximale}. L'obstruction est \emph{complète} : aucun flux d'air ne peut passer vers les poumons, ni en inspiration (impossible), ni en expiration (résiduelle seulement). La victime ne peut plus ventiler ses poumons $\to$ l'hypoxémie sévère et l'hypercapnie s'installent en \textbf{moins d'une minute} $\to$ perte de connaissance $\to$ arrêt cardiaque hypoxique.

\begin{definition}[Étouffement total]
Obstruction complète de la lumière des voies aériennes supérieures (larynx ou trachée proximale) par un corps étranger, rendant \textbf{impossible la parole, la toux efficace et la respiration}. C'est une \cle{urgence vitale absolue} : sans intervention dans les 1 à 3 minutes, le pronostic vital est engagé par asphyxie puis arrêt cardiaque.
\end{definition}

\paragraph{Distinction cruciale : étouffement partiel vs total.} Savoir différencier les deux tableaux cliniques conditionne la conduite à tenir.

\begin{center}
\begin{tabularx}{\linewidth}{|l|X|X|}
\hline
\rowcolor{popblueL}
\textbf{Signe clinique} & \textbf{Étouffement partiel (obstruction incomplète)} & \textbf{Étouffement total (obstruction complète)} \\
\hline
\textbf{Toux} & Présente, forte, efficace, sonore & \textbf{Absente ou silencieuse (gestuelle sans son)} \\
\hline
\textbf{Parole} & Possible (phrases courtes, voix étranglée) & \textbf{Impossible} \\
\hline
\textbf{Respiration} & Difficile, sifflante (stridor inspiratoire), angoissée & \textbf{Impossible (pas d'entrée d'air, tirage sous-claviculaire vain)} \\
\hline
\textbf{Cyanose} & Absente ou discrète (lèvres) & \textbf{Apparition rapide (lèvres, ongles, visage)} \\
\hline
\textbf{Attitude} & Victime consciente, anxieuse, tousse spontanément & \textbf{Victime paniquée, porte les mains à la gorge (signe universel), ne peut crier} \\
\hline
\textbf{Conduite à tenir} & \textbf{NE RIEN FAIRE} d'invasif. Encourager à \textbf{tousser} (seul moyen physiologique d'expulser). Surveiller. Ne pas frapper dans le dos (risque de basculer en obstruction totale). & \textbf{AGIR IMMÉDIATEMENT} : gestes d'expulsion (claques + compressions). Alerter les secours (112/119) \emph{pendant} les gestes si témoin unique, ou faire alerter par un tiers. \\
\hline
\end{tabularx}
\end{center}

\begin{attention}[Erreur fréquente : frapper dans le dos une victime qui tousse]
Si la victime \textbf{tousse efficacement}, elle a un étouffement \emph{partiel}. Lui taper dans le dos ou tenter une manœuvre de Heimlich risque de faire basculer le corps étranger en position plus distale et de transformer l'obstruction partielle en \textbf{obstruction totale}. \textbf{Conduite : « Laissez-la tousser, encouragez-la, ne la touchez pas ».} N'intervenez physiquement que si la toux devient inefficace ou silencieuse (basculement en étouffement total).
\end{attention}

\begin{popfigure}
\begin{tikzpicture}[scale=0.9, every node/.style={font=\small}]
% Anatomie sagittale simplifiée
% Tête
\draw[popdark, thick] (0,5) .. controls (1,5.5) and (2,5.5) .. (2.5,4.5) .. controls (3,3.5) and (3,2.5) .. (2.5,1.5); % Profil face
\draw[popdark, thick] (0,5) -- (0,1.5); % Dos
% Langue
\draw[poppink, thick] (0.5,4) .. controls (1,3.8) and (1.5,3.5) .. (1.8,3) .. controls (1.5,2.8) and (1,2.8) .. (0.5,3);
% Pharynx / Œsophage (arrière)
\draw[popmuted, thick] (0.2,4.5) -- (0.2,0.5);
% Trachée (avant)
\draw[popblue, thick] (1.8,3.5) -- (1.8,0.5);
% Anneaux trachéaux
\foreach \y in {3.2, 2.8, 2.4, 2.0, 1.6, 1.2, 0.8} {
    \draw[popblue, thick] (1.6,\y) -- (2.0,\y);
}
% Larynx (zone élargie)
\draw[poppurple, thick] (1.6,3.5) -- (1.4,3.8) -- (1.8,3.8) -- (2.0,3.5);
% Corps étranger (bolus) bloqué au larynx
\fill[poporange!80!popdark] (1.5,3.6) ellipse (0.25 and 0.15);
\node[poporange!80!popdark, anchor=west] at (2.1,3.6) {\textbf{Corps étranger}};
% Flèches : Claque dans le dos -> Pression -> Expulsion
\draw[popgreen, very thick, ->] (-1.5, 2.5) -- (-0.3, 2.5) node[midway, above] {\textbf{Claque dans le dos}};
\draw[popgreen, very thick, ->] (-0.2, 2.5) .. controls (0.5, 2.5) and (0.8, 2.8) .. (1.2, 3.2);
\draw[popgreen, dashed, thick, ->] (1.2, 3.2) .. controls (1.5, 3.5) and (1.7, 3.6) .. (1.8, 3.7) node[right] {Onde de pression};
% Flèche air expiratoire forcé (Heimlich)
\draw[popblue, very thick, ->] (1.8, 0.2) -- (1.8, 1.5) node[midway, right] {Air résiduel};
\draw[popblue, very thick, ->] (1.8, 1.5) -- (1.8, 3.0);
\draw[popblue, dashed, thick, ->] (1.8, 3.0) .. controls (1.8, 3.4) and (1.7, 3.6) .. (1.6, 3.7) node[left] {Surpression};
% Légendes anatomiques
\node[popdark, anchor=east] at (-0.2, 4.8) {Crâne};
\node[popdark, anchor=east] at (-0.2, 4.0) {Pharynx};
\node[poppurple, anchor=west] at (2.1, 3.9) {Larynx (glotte)};
\node[popblue, anchor=west] at (2.1, 2.0) {Trachée};
\node[popmuted, anchor=east] at (-0.3, 1.0) {Œsophage};
\node[poppink, anchor=west] at (2.1, 3.2) {Langue};
\end{tikzpicture}
\legende{Schéma sagittal des voies aéro-digestives. Le corps étranger (orange) obstrue l'entrée du larynx. \textbf{Flèches vertes :} claques dans le dos (victime penchée en avant) créant une onde de pression rétrograde et une vibration pour déloger le corps étranger. \textbf{Flèches bleues :} manœuvre de Heimlich (compressions abdominales) augmentant brutalement la pression intra-thoracique et intra-alvéolaire, expulsant l'air résiduel comme un piston pour chasser le corps étranger (principe du « bouchon de champagne »).}
\end{popfigure}

\courssub{Prise en charge de l'étouffement total chez l'adulte et l'enfant de plus d'un an}

\paragraph{Reconnaissance immédiate.} La victime est consciente, ne parle pas, ne tousse pas, ne respire pas. Elle porte souvent les deux mains à la gorge (signe universel de l'étouffement), le regard angoissé, la cyanose apparaît. \textbf{Interrogez : « Vous étouffez ? » $\to$ Hochement de tête ou impossibilité de répondre = confirmation.}

\paragraph{Algorithme standardisé (Recommandations ERC/ILCOR 2021, adaptées MINESEC).}
\begin{enumerate}
    \item \textbf{Appeler à l'aide / Alerter} (ou faire alerter le 112/119 par un témoin) \textbf{sans quitter la victime}.
    \item \textbf{Jusqu'à 5 claques dans le dos :} Victime debout ou assise, \textbf{penchée en avant} (tronc quasi horizontal) pour que la gravité aide l'expulsion. Une main soutient le haut du thorax (région sternale). L'autre main (paume) frappe \textbf{fermement entre les deux omoplates} (région inter-scapulaire). Observer après chaque claque : expulsion ?
    \item \textbf{Si échec : jusqu'à 5 compressions abdominales (Manœuvre de Heimlich) :} Se placer \textbf{derrière la victime}, passer les bras sous les siens. \textbf{Poing fermé} (pouce vers l'intérieur) placé \textbf{au-dessus du nombril, bien en dessous du processus xiphoïde du sternum} (éviter le sternum et les côtes flottantes). Saisir le poing avec l'autre main. Effectuer \textbf{5 tractions vigoureuses, rapides, vers soi et vers le haut} (mouvement en « J » ou en cuillère). Observer après chaque compression.
    \item \textbf{Alterner} les séries de 5 claques et 5 compressions jusqu'à expulsion du corps étranger, ou jusqu'à ce que la victime perde connaissance.
\end{enumerate}

\begin{methode}[Manœuvre de Heimlich (Compressions abdominales) — Adulte / Enfant > 1 an]
\textbf{Position :} Derrière la victime, bras sous ses aisselles.\\
\textbf{Repérage :} Poing fermé, face pouciaire vers l'abdomen, placé \textbf{au-dessus du nombril et au-dessous du sternum (processus xiphoïde)}.\\
\textbf{Prise :} Saisir le poing avec l'autre main.\\
\textbf{Action :} \textbf{5 compressions rapides, vigoureuses, dirigées vers le haut et vers l'arrière (vers soi)}.\\
\textbf{Vérification :} Après chaque compression, regarder dans la bouche si le corps étranger est visible et accessible (ne pas faire de « balayage digital » à l'aveugle).\\
\textbf{Répétition :} Alterner avec 5 claques dans le dos jusqu'à succès ou perte de connaissance.
\end{methode}

\begin{savaistu}[Le principe physique de la manœuvre de Heimlich (1974)]
Le Dr Henry Heimlich a démontré que la compression brutale de l'abdomen (piston abdominal) transmet une \textbf{surpression subdiaphragmatique} aux poumons via le diaphragme. Cette surpression (pouvant atteindre \SI{100}{mmHg} soit \SI{13,3}{kPa}) force l'air résiduel des alvéoles à remonter violemment dans la trachée. Ce flux d'air expiratoire forcé, bien que la glotte soit fermée par le corps étranger, agit comme un \textbf{« coup de bélier » pneumatique} qui déloge le corps étranger vers le haut et l'extérieur. C'est l'équivalent physiologique d'éternuer ou de tousser avec une force décuplée par l'effort musculaire du secouriste.
\end{savaistu}

\begin{attention}[Erreur fréquente et dangereuse : le « balayage digital » à l'aveugle]
Ne \textbf{JAMAIS} introduire ses doigts dans la bouche de la victime pour « chercher » ou « tirer » le corps étranger si vous ne le voyez \textbf{PAS distinctement}. Risque majeur : \textbf{enfoncer le corps étranger plus profondément} (transformant une obstruction sus-glottique en obstruction sous-glottique, impossible à déloger par Heimlich) ou léser le palais mou / l'épiglotte. \textbf{Règle :} On ne retire que ce qui est \textbf{visiblement accessible} au fond de la bouche (ex: dentier, gros morceau visible) avec un doigt en crochet (index), idéalement sous contrôle visuel direct.
\end{attention}

\courssub{Adaptation des gestes pour le nourrisson (moins de 1 an)}

\paragraph{Pourquoi pas de Heimlich ?} Le nourrisson a un foie volumineux, fragile, occupant une grande partie de l'hypochondre droit, juste sous le diaphragme. Les compressions abdominales violentes risquent \textbf{rupture hépatique, hémorragie interne, lésion rate/rate}. \textbf{Interdit formel.}

\paragraph{Technique adaptée (alternance dos / thorax) :}
\begin{enumerate}
    \item \textbf{5 claques dans le dos :} Asseyez-vous ou genou à terre. Placez le nourrisson \textbf{à califourchon sur votre avant-bras}, tête plus basse que le tronc (inclinaison \SI{30}{\degree}--\SI{45}{\degree}), en maintenant la mâchoire (pouce/index en « V » sur les maxillaires) pour garder la bouche ouverte et protéger la tête. Frappez \textbf{fermement entre les omoplates} (paume de l'autre main), 5 fois. Vérifiez.
    \item \textbf{Si échec : 5 compressions thoraciques :} Retournez le nourrisson sur l'autre avant-bras (ou sur votre cuisse), toujours \textbf{tête basse}. \textbf{2 doigts (index + majeur) sur le tiers inférieur du sternum} (un doigt de large sous la ligne intermammaire). Comprimez le sternum de \textbf{\SI{1,5}{cm} à \SI{2}{cm}} (environ 1/3 du diamètre antéro-postérieur), 5 fois, rythme \SI{100}{--}\SI{120}{/min}. Vérifiez.
    \item \textbf{Alterner} jusqu'à expulsion ou perte de connaissance.
\end{enumerate}

\begin{attention}[Danger absolu : Manœuvre de Heimlich sur nourrisson]
Pratiquer des compressions abdominales (Heimlich) sur un nourrisson de moins d'un an expose à un risque \textbf{très élevé de rupture du foie} (hémorragie interne mortelle) et de lésion de la rate. \textbf{Seules les claques dans le dos (tête basse) et les compressions thoraciques (2 doigts) sont autorisées.}
\end{attention}

\courssub{Adaptation des gestes pour la femme enceinte (et la personne obèse)}

\paragraph{Contre-indication aux compressions abdominales.} Chez la femme enceinte (utérus gravide volumineux, surtout 2\textsuperscript{e} et 3\textsuperscript{e} trimestres) et chez la personne très obèse (paroi abdominale épaisse, foie souvent hypertrophié/stéatosique), les compressions abdominales sont \textbf{inefficaces} (transmission de pression atténuée) et \textbf{dangereuses} (risque de décollement placentaire, rupture utérine, traumatisme viscéral).

\paragraph{Technique : Compressions thoraciques (sternales) après les claques.}
\begin{enumerate}
    \item \textbf{5 claques dans le dos} (identiques : victime penchée en avant, entre omoplates).
    \item \textbf{Si échec : 5 compressions thoraciques :} Secouriste derrière la victime. \textbf{Poing fermé sur la moitié inférieure du sternum} (éviter le processus xiphoïde). Saisir le poing avec l'autre main. \textbf{Tractions vers l'arrière (vers soi), sans composante vers le haut marquée} (pour ne pas comprimer l'abdomen). Rythme \SI{100}{--}\SI{120}{/min}.
    \item Alterner claques / compressions thoraciques.
\end{enumerate}

\courssub{Évolution vers l'inconscience et prévention}

\paragraph{Si la victime perd connaissance pendant les manœuvres :}
\begin{enumerate}
    \item La \textbf{poser délicatement au sol} (protéger la tête).
    \item \textbf{Alerter / Faire alerter les secours (112/119)} si pas encore fait.
    \item \textbf{Commencer immédiatement la RCP} (30 compressions thoraciques / 2 insufflations). \textbf{Les compressions thoraciques de la RCP génèrent elles-mêmes une surpression intrathoracique capable d'expulser le corps étranger.} Avant chaque tentative d'insufflation, regarder dans la bouche : si le corps étranger est visible et accessible, le retirer (doigt en crochet). Ne pas pratiquer de « balayage digital » à l'aveugle.
\end{enumerate}

\paragraph{Prévention primaire (éducation à la santé, contexte camerounais).}
\begin{itemize}
    \item \textbf{Enfants :} Couper les aliments ronds/durs (arachides, maïs grillé, raisins, saucisses, bonbons durs) en \textbf{petits morceaux} (pas de morceaux ronds de taille trachéale). Ne pas donner d'arachides/maïs grillé aux \textbf{moins de 4--5 ans}. Surveiller les repas (pas de course/jeu la bouche pleine).
    \item \textbf{Adultes :} Bien mâcher (prothèses dentaires adaptées). Éviter de parler/rire avec la bouche pleine. Modérer l'alcool (diminue le réflexe de toux et la vigilance).
    \item \textbf{Personnes âgées / dysphagiques :} Adapter la texture (mixé, épaissi), position assise droite pendant et après le repas.
    \item \textbf{Environnement :} Éloigner les petits objets (boutons, perles, piles boutons, pièces de monnaie) de la portée des nourrissons (exploration buccale physiologique).
\end{itemize}

\begin{popfigure}
\begin{tikzpicture}[font=\small, scale=0.95]
% Tableau comparatif 3 colonnes
% En-têtes
\draw[popblue, thick] (0,0) rectangle (14,1.2);
\node[popblue, font=\bfseries] at (2,0.6) {Adulte / Enfant > 1 an};
\node[poppurple, font=\bfseries] at (7,0.6) {Nourrisson < 1 an};
\node[poppink, font=\bfseries] at (12,0.6) {Femme enceinte / Obèse};

% Lignes
\draw[popline, thick] (0,0) -- (14,0);
\draw[popline, thick] (0,1.2) -- (14,1.2);
\draw[popline, thick] (0,2.4) -- (14,2.4);
\draw[popline, thick] (0,3.6) -- (14,3.6);
\draw[popline, thick] (0,4.8) -- (14,4.8);
\draw[popline, thick] (0,6.0) -- (14,6.0);
% Colonnes
\draw[popline, thick] (4.67,0) -- (4.67,6);
\draw[popline, thick] (9.33,0) -- (9.33,6);

% Contenu
% Ligne 1 : Claques dans le dos
\node[anchor=west] at (0.2,5.4) {\textbf{Claques dans le dos}};
\node[popgreen, align=center] at (2.3,5.4) {\checkmark\ \textbf{AUTORISÉES}\\(5 max, inter-scapulaires,\\victime penchée en avant)};
\node[popgreen, align=center] at (7,5.4) {\checkmark\ \textbf{AUTORISÉES}\\(5 max, inter-scapulaires,\\nourrisson tête basse sur avant-bras)};
\node[popgreen, align=center] at (11.7,5.4) {\checkmark\ \textbf{AUTORISÉES}\\(5 max, inter-scapulaires,\\victime penchée en avant)};

% Ligne 2 : Compressions abdominales (Heimlich)
\node[anchor=west, align=center] at (0.2,4.2){\textbf{Compressions abdominales (Heimlich)}};
\node[popgreen, align=center] at (2.3,4.2) {\checkmark\ \textbf{AUTORISÉES}\\(Poing au-dessus nombril,\\sous sternum, traction haut/arrière)};
\node[poporange!80!popdark, align=center] at (7,4.2) {\textbf{INTERDITES}\\(Risque rupture hépatique,\\lésions viscérales graves)};
\node[poporange!80!popdark, align=center] at (11.7,4.2) {\textbf{INTERDITES}\\(Risque décollement placentaire,\\rupture utérine, traumatisme viscéral)};

% Ligne 3 : Compressions thoraciques
\node[anchor=west, align=center] at (0.2,3.0){\textbf{Compressions thoraciques (Sternum)}};
\node[poporange!80!popdark, align=center] at (2.3,3.0) {Second choix\\(si Heimlich échoue/impossible)};
\node[popgreen, align=center] at (7,3.0) {\checkmark\ \textbf{AUTORISÉES}\\(2 doigts, tiers inf. sternum,\\prof. 1,5--2 cm, 5 max)};
\node[popgreen, align=center] at (11.7,3.0) {\checkmark\ \textbf{RECOMMANDÉES (1\textsuperscript{er} choix après claques)}\\(Poing sur sternum, traction arrière)};

% Ligne 4 : Position spécifique
\node[anchor=west] at (0.2,1.8) {\textbf{Position victime}};
\node[align=center] at (2.3,1.8) {Debout/Assis,\\penché en avant};
\node[align=center] at (7,1.8) {Allongé ventre sur\\avant-bras, tête basse};
\node[align=center] at (11.7,1.8) {Debout/Assis,\\penché en avant};

% Ligne 5 : Si inconscience
\node[anchor=west] at (0.2,0.6) {\textbf{Si inconscience}};
\node[align=center, popblue] at (2.3,0.6) {RCP standard\\(30 comp. / 2 ins.)};
\node[align=center, poppurple] at (7,0.6) {RCP nourrisson\\(30 comp. 2 doigts / 2 ins.)};
\node[align=center, poppink] at (11.7,0.6) {RCP standard\\(30 comp. / 2 ins.)};
\end{tikzpicture}
\legende{Tableau comparatif synthétique des gestes de premiers secours face à un étouffement total selon le profil de la victime. \textbf{Vert = geste recommandé / autorisé. Orange = geste interdit ou non prioritaire.} La règle d'or : alterner 5 claques dans le dos et 5 compressions (abdominales, thoraciques 2 doigts, ou thoraciques poing) jusqu'à expulsion ou perte de connaissance.}
\end{popfigure}

\begin{exoresolu}[Cas pratique : Étouffement lors d'un repas de fête à Douala]
\textbf{Énoncé :} Lors d'un repas de fin d'année, M. K., 52 ans, se lève brusquement de table, les mains au cou, le visage rouge puis cyanosé. Il ne peut ni parler, ni tousser, ni respirer. Sa femme paniquée tente de lui donner de l'eau à boire. Un neveu, élève en Terminale D, intervient.

\textbf{Questions :}
\begin{enumerate}
    \item Qualifiez le type d'étouffement (partiel ou total) et justifiez par trois signes cliniques.
    \item Expliquez pourquoi donner de l'eau est une erreur grave.
    \item Décrivez précisément la séquence des gestes que le neveu doit réaliser jusqu'à l'arrivée des secours ou la résolution.
    \item Si M. K. perd connaissance avant l'expulsion, quelle conduite tenir ?
\end{enumerate}

\tcblower
\textbf{Solution détaillée :}
\begin{enumerate}
    \item \textbf{Étouffement TOTAL (obstruction complète).} Justification : (1) \textbf{Impossibilité de parler} (aphonie totale) ; (2) \textbf{Impossibilité de tousser} (absence de son, gestuelle vaine) ; (3) \textbf{Impossibilité de respirer} (pas d'entrée d'air, cyanose rapide, tirage sous-claviculaire paradoxal). Le signe des mains au cou (signe universel) confirme la détresse respiratoire aiguë.
    \item \textbf{Erreur grave :} Donner de l'eau à une victime d'étouffement total est \textbf{formellement contre-indiqué}. L'eau ne peut pas passer (obstruction complète) et risque de : (a) provoquer une fausse route supplémentaire (inhalation d'eau dans les poumons $\to$ œdème aigu du poumon) ; (b) faire glisser le corps étranger plus bas (trachée) ; (c) retarder les gestes efficaces (claques/Heimlich) ; (d) compliquer la RCP ultérieure (liquide dans voies aériennes).
    \item \textbf{Séquence gestes (Adulte) :}
    \begin{enumerate}
        \item Cri d'alerte : « Quelqu'un appelle le 112 (ou 119) ! » (Ne pas quitter la victime).
        \item Se placer derrière M. K., le pencher en avant (tronc horizontal), une main sur le sternum pour soutenir.
        \item \textbf{5 claques fermes} paume de la main entre les omoplates. Vérifier après chaque claque.
        \item Si échec : \textbf{Manœuvre de Heimlich}. Poing fermé (pouce ventral) au-dessus du nombril, sous le sternum. Saisir le poing de l'autre main. \textbf{5 tractions vigoureuses vers le haut et vers l'arrière (vers soi)}. Vérifier.
        \item \textbf{Alterner} séries de 5 claques / 5 Heimlich jusqu'à expulsion du corps étranger (souvent bruyant, suivi d'une grande inspiration sifflante puis toux grasse) ou perte de connaissance.
    \end{enumerate}
    \item \textbf{Si perte de connaissance :}
    \begin{enumerate}
        \item Allonger M. K. délicatement au sol (protéger la tête).
        \item Confirmer l'alerte (112/119) si pas fait.
        \item \textbf{Commencer RCP immédiatement :} 30 compressions thoraciques (centre du sternum, \SI{5}{--}\SI{6}{cm}, \SI{100}{--}\SI{120}{/min}) $\to$ Ouvrir la voie aérienne (relevé du menton) $\to$ Regarder dans la bouche : si corps étranger \textbf{visible et accessible} $\to$ retirer (doigt en crochet). Sinon $\to$ 2 insufflations (bouche-à-bouche ou masque). $\to$ Reprendre 30 compressions.
        \item \textbf{Ne pas pratiquer de Heimlich sur victime inconsciente au sol.} Les compressions thoraciques de la RCP remplacent la manœuvre de Heimlich (même principe de surpression intrathoracique).
    \end{enumerate}
\end{enumerate}
\end{exoresolu}

\begin{aretenir}[Synthèse : Conduite à tenir unifiée face à l'étouffement total]
\begin{itemize}
    \item \textbf{Reconnaître :} Victime consciente, mains au cou, \textbf{ne parle pas, ne tousse pas, ne respire pas} (cyanose).
    \item \textbf{Alerter :} 112 / 119 (ou faire alerter) \textbf{avant ou pendant} les gestes.
    \item \textbf{Agir selon le profil :}
    \begin{tabular}{ll}
    \textbf{Adulte / Enfant > 1 an} & 5 claques dos (penché avant) $\leftrightarrow$ 5 Heimlich (poing nombril/sternum) \\
    \textbf{Nourrisson < 1 an} & 5 claques dos (tête basse sur avant-bras) $\leftrightarrow$ 5 compressions thoraciques (2 doigts sternum) \\
    \textbf{Femme enceinte / Obèse} & 5 claques dos (penché avant) $\leftrightarrow$ 5 compressions thoraciques (poing sternum)
    \end{tabular}
    \item \textbf{Interdits absolus :} Heimlich sur nourrisson / femme enceinte ; Balayage digital à l'aveugle ; Boire / manger / parler pendant l'étouffement ; Frapper dans le dos si la victime tousse encore (étouffement partiel).
    \item \textbf{Évolution :} Si inconscience $\to$ Sol $\to$ Alerte $\to$ RCP standard (les compressions expulsent le corps étranger).
\end{itemize}
\end{aretenir}

\coursec{La perte de connaissance : reconnaissance et position latérale de sécurité}

\courssub{Transition : de l'obstruction aiguë à la surveillance de l'inconscient qui respire}

Dans la section précédente, nous avons vu comment intervenir face à une \cle{obstruction totale des voies aériennes} par un corps étranger : la victime est consciente mais ne peut plus respirer, parler ni tousser. Nous allons maintenant aborder une situation radicalement différente : la victime a \emph{perdu connaissance} (elle ne répond plus), \textbf{mais elle respire encore normalement}. C'est une urgence vitale d'un autre type : le danger ne vient pas d'un objet coincé, mais de l'inconscience elle-même qui fait perdre les réflexes de protection des voies aériennes. Savoir reconnaître cet état et réaliser la \cle{Position Latérale de Sécurité (PLS)} est la compétence clé pour prévenir l'étouffement secondaire par la langue ou les vomissements.

\courssub{Définition, causes et mécanisme physiopathologique}

\begin{definition}[Perte de connaissance (malaise, syncope, coma)]
La \textbf{perte de connaissance} est l'\textbf{abolition temporaire ou prolongée de la conscience} avec \textbf{perte de la réponse aux stimulations verbales et tactiles} (la victime ne répond pas quand on lui parle, ne réagit pas quand on lui secoue doucement les épaules), \textbf{la ventilation spontanée étant conservée} (la victime respire normalement). Elle se distingue de l'arrêt cardiaque où la respiration est absente ou agonique (gasps).
\end{definition}

\begin{savaistu}[Pourquoi le cerveau « s'éteint-il » ?]
Le cerveau est l'organe le plus gourmand en oxygène : il ne représente que \SI{2}{\%} de la masse corporelle mais consomme \SI{20}{\%} de l'oxygène et du glucose. Toute chute brutale de la \cle{perfusion cérébrale} (pression artérielle trop basse, arythmie, hémorragie), de la \cle{glycémie} (hypoglycémie sévère chez le diabétique), ou toute perturbation électrique brutale (crise d'épilepsie) ou toxique (alcool, médicaments, monoxyde de carbone) entraîne une défaillance du \cle{système activateur réticulé ascendant} (tronc cérébral), centre de la vigilance. Au Cameroun, les causes fréquentes en milieu scolaire ou familial sont le \cle{paludisme cérébral}, l'\cle{hypoglycémie} (jeûne prolongé, effort intense), les \cle{traumatismes crâniens} (accidents de moto, chutes), les \cle{crises d'épilepsie} et les \cle{accidents vasculaires cérébraux (AVC)} chez l'adulte âgé.
\end{savaistu}

\begin{experience}[Observation clinique : le danger de la position dorsale chez l'inconscient]
\textbf{Contexte :} Une victime est retrouvée inconsciente, allongée sur le dos (décubitus dorsal), respirant bruyamment (ronflement).
\textbf{Observation :} On entend un bruit de \emph{ronflement} (stridor inspiratoire) à chaque inspiration. La poitrine se soulève difficilement. Si on bascule la tête en arrière (libération des voies aériennes), le bruit disparaît et la respiration devient silencieuse et ample.
\textbf{Interprétation :} En décubitus dorsal, l'\cle{hypotonie musculaire} totale (relaxation de tous les muscles) fait chuter la \cle{langue} en arrière par gravité : sa base obstrue le \cle{pharynx} (voie aérienne supérieure). C'est une \textbf{obstruction mécanique par la langue}. De plus, le sphincter œsophagien inférieur étant relâché, tout contenu gastrique (vomissements, salive, sang) peut refluer passivement dans le pharynx et être aspiré vers la trachée (\cle{fausse route}, \cle{pneumonie d'inhalation}).
\textbf{Conclusion :} La position dorsale est \textbf{interdite} chez tout inconscient qui respire. Il faut utiliser la gravité pour écarter la langue et laisser s'écouler les liquides : c'est le principe de la PLS.
\end{experience}

\begin{popfigure}
\begin{tikzpicture}[scale=0.9, every node/.style={font=\small}]
% Couleurs
\colorlet{airway}{popblue!80!black}
\colorlet{tongue}{poppink}
\colorlet{spine}{popdark}
\colorlet{skin}{popcream}
\colorlet{arrowcol}{popgreen!70!black}

% --- GAUCHE : Décubitus dorsal (Obstruction) ---
\begin{scope}[xshift=0cm]
% Tête - profil sagittal
\draw[skin, line width=2pt, fill=skin!30] (0,0) .. controls (0,3) and (2,4) .. (3,4) .. controls (4,4) and (4.5,3.5) .. (5,2.5) .. controls (5.5,1.5) and (5,0.5) .. (4,0) .. controls (3,-0.5) and (1,-0.5) .. (0,0);
% Crâne
\draw[popdark, line width=1.5pt] (0,3) .. controls (0.5,3.8) and (2.5,4.2) .. (3.5,4) .. controls (4.2,3.8) .. (4.5,3.2);
% Colonne cervicale
\draw[spine, line width=2pt] (1.5,0) -- (1.5,-1.5) node[below, font=\tiny] {Colonne vertébrale};
% Langue (obstruante)
\draw[tongue, fill=tongue!40, line width=1pt] (2.2,1.5) .. controls (2.8,1.8) and (3.2,1.5) .. (3.5,1.0) .. controls (3.2,0.5) and (2.5,0.5) .. (2.2,0.8) .. controls (2.0,1.1) .. cycle;
\node[tongue, align=center, font=\tiny\bfseries] at (3.0, 1.3) {Langue\\(chute postérieure)};
% Voies aériennes (obstruées)
\draw[airway, line width=1.5pt, dashed] (3.5,2.0) -- (3.5,1.5); % Pharynx
\draw[airway, line width=1.5pt] (3.5,1.5) -- (3.5,0.5); % Trachée (partie visible)
\node[airway, align=center, font=\tiny] at (4.5, 1.2) {Voies aériennes\\ \small\textbf{OBSTRUÉES}};
% Flèches air bloqué
\draw[poporange, ->, thick, decoration={snake, amplitude=1pt, segment length=4pt}, decorate] (4.5, 2.5) -- (3.8, 2.0) node[above right, font=\tiny, poporange] {Air (difficile)};
\draw[poporange, ->, thick] (3.8, 0.8) -- (3.8, 0.3) node[right, font=\tiny, poporange] {Vers poumons};

% Légende schéma gauche
\node[align=center, font=\small\bfseries, popdark] at (2.5, -2.2) {A. Décubitus dorsal : Obstruction par la langue};
\end{scope}

% --- DROITE : PLS (Voies aériennes libres) ---
\begin{scope}[xshift=7.5cm]
% Tête - profil sagittal (tournée sur le côté, menton relevé)
\draw[skin, line width=2pt, fill=skin!30] (0,0) .. controls (0,3) and (2,4) .. (3,4) .. controls (4,4) and (4.5,3.5) .. (5,2.5) .. controls (5.5,1.5) and (5,0.5) .. (4,0) .. controls (3,-0.5) and (1,-0.5) .. (0,0);
% Crâne
\draw[popdark, line width=1.5pt] (0,3) .. controls (0.5,3.8) and (2.5,4.2) .. (3.5,4) .. controls (4.2,3.8) .. (4.5,3.2);
% Colonne cervicale (alignée)
\draw[spine, line width=2pt] (1.5,0) -- (1.5,-1.5) node[below, font=\tiny] {Colonne vertébrale};
% Langue (tombée en avant par gravité)
\draw[tongue, fill=tongue!40, line width=1pt] (1.8,2.5) .. controls (1.5,2.8) and (1.2,2.5) .. (1.0,2.0) .. controls (1.1,1.5) and (1.5,1.5) .. (1.8,1.8) .. controls (2.0,2.1) .. cycle;
\node[tongue, align=center, font=\tiny\bfseries] at (1.3, 2.2) {Langue\\(antérieure)};
% Voies aériennes (libres)
\draw[airway, line width=2pt] (2.5,2.5) -- (2.5,0.5); % Trachée ouverte
\node[airway, align=center, font=\tiny] at (3.8, 1.5) {Voies aériennes\\ \small\textbf{LIBRES}};
% Écoulement liquides
\draw[popblue, ->, thick, decoration={snake, amplitude=1.5pt, segment length=3pt}, decorate] (2.5, 2.2) -- (1.5, 1.5) node[left, font=\tiny, popblue, align=center]{Écoulement\\ sécrétions/vomiss.};
% Flèches air libre
\draw[arrowcol, ->, thick] (4.0, 2.8) -- (3.0, 2.5) node[above left, font=\tiny, arrowcol] {Air (libre)};
\draw[arrowcol, ->, thick] (2.5, 0.8) -- (2.5, 0.3) node[right, font=\tiny, arrowcol] {Vers poumons};

% Sol (représentation gravité)
\draw[popmuted, line width=3pt] (-0.5, -0.2) -- (5.5, -0.2);
\draw[->, popmuted, thick] (5.8, 0) -- (5.8, -1) node[right, font=\tiny, popmuted] {Gravité};

% Légende schéma droit
\node[align=center, font=\small\bfseries, popdark] at (2.5, -2.2) {B. Position Latérale de Sécurité : Voies aériennes dégagées};
\end{scope}
\end{tikzpicture}
\legende{Comparaison sagittale de la perméabilité des voies aériennes. \textbf{A.} En décubitus dorsal, l'hypotonie linguale fait obstruction (flèches orange : passage d'air turbulent). \textbf{B.} En PLS, la gravité tire la langue en avant (vers le menton) et ouvre le pharynx ; les liquides s'écoulent par la bouche (flèches bleues ondulées). La bascule de la tête en arrière (extension cervicale) aligne les axes bucco-pharyngo-trachéaux.}
\end{popfigure}

\courssub{Reconnaissance : distinguer l'inconscient qui respire de l'arrêt cardiaque}

C'est l'étape \textbf{critique} qui conditionne tout le reste. Une erreur de diagnostic (confondre respiration agonique et respiration normale) retarde la RCP et réduit les chances de survie.

\begin{methode}[Bilan initial : « Regarder, Écouter, Sentir » (\SI{10}{\second} max)]
\begin{enumerate}
    \item \textbf{Assurer la protection} : sécuriser les lieux (circulation, électricité, feu, foule).
    \item \textbf{Évaluer l'état de conscience} : s'agenouiller à côté de la victime. Lui parler fort (\emph{« Entendez-vous ? Ouvrez les yeux ! Serrez ma main ! »}) et \textbf{secouer vigoureusement les deux épaules} (stimulation tactile douloureuse modérée).
    \item \textbf{Si pas de réponse} : la victime est \textbf{inconsciente}.
    \item \textbf{Libérer les voies aériennes} : une main sur le front pour basculer la tête en arrière (\cle{bascule prudente de la tête}), l'index et le majeur de l'autre main sous le menton pour \textbf{relever le menton} (élévation du menton). \emph{Ne pas appuyer sur les parties molles du cou.}
    \item \textbf{Vérifier la respiration} (pendant \textbf{\SI{10}{\second} maximum}) :
    \begin{itemize}
        \item \textbf{Regarder} : le thorax et l'abdomen se soulèvent-ils régulièrement ?
        \item \textbf{Écouter \&\ Sentir} : poser sa joue au-dessus de la bouche/nez de la victime pour sentir le flux d'air expiré et entendre les bruits respiratoires.
    \end{itemize}
    \item \textbf{Interprétation} :
    \begin{itemize}
        \item \textbf{Respiration NORMALE} (régulière, ample, silencieuse ou peu bruyante, $\ge$ \SI{2}{mouvements} en \SI{10}{\second} soit \SIrange{10}{12}{cycles/min}) $\rightarrow$ \textbf{PLS + Alerte + Surveillance}.
        \item \textbf{PAS de respiration} OU \textbf{Respiration AGONIQUE} (mouvements respiratoires lents, irréguliers, bruyants, « gasps », halètements, \textbf{$<$ \SI{2}{mouvements} en \SI{10}{\second}}) $\rightarrow$ \textbf{ARRÊT CARDIAQUE} $\rightarrow$ \textbf{Alerte + RCP immédiate}.
    \end{itemize}
\end{enumerate}
\end{methode}

\begin{attention}[Piège mortel : la respiration agonique (gasping)]
Chez \SIrange{30}{40}{\%} des arrêts cardiaques en début d'évolution, le tronc cérébral génère des \textbf{mouvements respiratoires anormaux} : la bouche s'ouvre, les épaules se soulèvent brutalement, un bruit de râle ou de hoquet se produit. \textbf{CE N'EST PAS UNE RESPIRATION NORMALE.} C'est un signe de mort encéphalique imminente. \textbf{Ne pas faire de PLS !} Il faut alerter et commencer les compressions thoraciques immédiatement. \emph{Règle d'or : « En cas de doute sur la respiration, considérez qu'elle est absente et commencez la RCP. »}
\end{attention}

\begin{exoresolu}[Diagnostic différentiel : Syncope vasovagale vs Arrêt cardiaque]
\textbf{Énoncé :} Lors d'une assemblée générale dans une salle des fêtes à Yaoundé, M. K., 62 ans, diabétique, s'effondre soudainement. Il est au sol, yeux fermés. Vous êtes le premier secouriste. Vous lui parlez, secouez ses épaules : pas de réponse. Vous libérez les voies aériennes (tête en arrière, menton relevé). Vous observez pendant \SI{10}{\second} : son abdomen se soulève régulièrement toutes les \SI{5}{\second} environ, vous entendez un flux d'air calme sur votre joue. Que faites-vous ?
\textbf{Solution détaillée :}
\begin{enumerate}
    \item \textbf{Analyse :} Victime \textbf{inconsciente} (pas de réponse aux stimulations) \textbf{MAIS respiration NORMALE} (régulière, fréquence ~\SI{12}{min^{-1}}, flux d'air perçu). Ce n'est \textbf{PAS} un arrêt cardiaque. C'est une perte de connaissance (probable syncope vasovagale, hypoglycémie ou AVC).
    \item \textbf{Action immédiate :} Mettre en \textbf{Position Latérale de Sécurité (PLS)} pour protéger les voies aériennes (langue, vomissements possibles).
    \item \textbf{Alerte :} Appeler le \textbf{112} (numéro d'urgence européen, fonctionne au Cameroun depuis mobile) ou le \textbf{119} (SAMU Cameroun) ou les \textbf{Sapeurs-Pompiers (\SI{118}{})} ou demander à un témoin d'appeler en précisant : « Inconscient, respire, mis en PLS ». Ne pas raccrocher avant l'opérateur.
    \item \textbf{Surveillance :} Rester à côté, lui parler, vérifier la respiration \textbf{toutes les minutes} jusqu'à l'arrivée des secours. Si la respiration s'arrête ou devient agonique : remettre sur le dos dur (plan dur) et démarrer la RCP (\SI{30}{ compressions} / \SI{2}{ insufflations}).
    \item \textbf{Recherche cause :} Si un témoin signale qu'il est diabétique et n'a pas mangé $\rightarrow$ hypoglycémie possible. \textbf{Ne rien donner par la bouche} (risque fausse route). Attendre les secours (injection de glucagon ou perfusion de sérum glucosé).
\end{enumerate}
\end{exoresolu}

\courssub{La Position Latérale de Sécurité (PLS) : technique normalisée en 5 étapes}

La PLS est une position \textbf{latérale, stable, au sol}, qui maintient les voies aériennes pérennes grâce à la gravité. Elle doit être réalisable par un seul secouriste, sans matériel, en protégeant la colonne vertébrale (autant que possible) et en permettant une surveillance continue.

\begin{methode}[Mise en PLS : 5 étapes codifiées (Recommandations ERC/ILCOR 2021 adaptées)]
\textbf{Position de départ :} Victime allongée sur le dos (décubitus dorsal), jambes tendues, bras le long du corps. Secouriste agenouillé à côté d'elle (côté choisi arbitrairement, idéalement côté où il y a de la place).

\begin{enumerate}
    \item \textbf{Bras proche en angle droit :} Saisir le bras de la victime le plus proche de vous. Le placer \textbf{perpendiculaire au corps} (angle droit à l'épaule), coude plié à \SI{90}{\degree} (main vers la tête, paume vers le haut). \emph{Cet bras servira d'appui antérieur pour stabiliser la position finale.}
    \item \textbf{Bras lointain contre la joue :} Saisir l'autre bras (lointain) au niveau du poignet. Le porter \textbf{contre la joue opposée} (côté secouriste), paume de la victime contre sa propre joue. \emph{Votre main maintient cette position (votre paume contre le dos de sa main). Cela protège la tête lors du basculement et maintient la bouche basse.}
    \item \textbf{Plier la jambe lointaine :} De votre main libre (celle qui était sur le front), saisir la \textbf{jambe opposée} (lointaine) juste au-dessus du genou. La \textbf{plier} en posant le pied à plat au sol (genou relevé). \emph{Ce genou sera le levier de rotation.}
    \item \textbf{Basculer (Rouler vers soi) :} En gardant une main sur la joue/bras de la victime (pour guider la tête) et l'autre sur le genou plié, \textbf{tirer le genou vers vous} pour faire rouler la victime \textbf{sur le côté vers vous}. La victime vient se coucher sur son flanc, le genou plié au sol servant d'ancrage. La tête repose sur le dos de sa propre main (menton relevé).
    \item \textbf{Ajuster et Dégager les voies aériennes :}
    \begin{itemize}
        \item Ajuster la \textbf{jambe supérieure} (celle qui est pliée) : hanche et genou fléchis à angle droit (position stable « en chien de fusil »).
        \item \textbf{Re-basculer la tête en arrière} (une main front, deux doigts menton) pour \textbf{ouvrir la bouche} (pouce sur le menton) et laisser l'écoulement libre.
        \item Vérifier que la respiration est toujours normale.
    \end{itemize}
\end{enumerate}
\end{methode}

\begin{popfigure}
\begin{tikzpicture}[scale=0.75, every node/.style={font=\small}]
% Styles
\tikzset{
    victim/.style={draw=popdark, line width=1.5pt, fill=popcream!50, rounded corners=3pt},
    ground/.style={draw=popmuted, line width=3pt},
    arrow/.style={->, thick, popblue},
    stepnum/.style={circle, draw=popblue, fill=popblue, text=white, font=\small\bfseries, minimum size=0.6cm},
    labelbox/.style={rectangle, draw=popdark, fill=white, rounded corners=2pt, inner sep=2pt, font=\tiny, align=center}
}

% --- ÉTAPE 1 : Position initiale + Bras proche ---
\begin{scope}[xshift=0cm, yshift=0cm]
\draw[ground] (-1, -0.5) -- (5, -0.5);
% Victime dos
\draw[victim] (0.5, 0) ellipse (1.5 and 0.6); % Tronc
\draw[victim] (0.5, 1.2) circle (0.6); % Tête
% Bras proche (gauche) en angle droit
\draw[victim, line width=4pt] (-1.0, 0.8) -- (-1.0, 1.8) -- (-0.2, 1.8);
% Bras lointain (droit) le long du corps
\draw[victim, line width=4pt] (2.0, 0.2) -- (2.8, -0.5);
% Jambes
\draw[victim, line width=4pt] (0.5, -0.5) -- (0.5, -1.8);
\draw[victim, line width=4pt] (1.5, -0.5) -- (1.5, -1.8);
% Secouriste (agenouillé à gauche)
\draw[popblue, line width=2pt] (-2.5, -0.2) -- (-2.5, -1.2) -- (-1.8, -1.2);
\draw[popblue, fill=popblue!20] (-2.5, 0.2) circle (0.3);
\node[stepnum] at (-3.5, 1.5) {1};
\node[labelbox, popblue, align=center] at (-3.5, 0.8){Bras proche\\$\approx$ 90°};
\draw[arrow] (-3.0, 1.2) .. controls (-2.0, 1.5) .. (-1.0, 1.5);
\node[align=center, font=\small\bfseries, popdark] at (0.5, -2.5) {Étape 1 : Position initiale\\\& Bras proche en angle droit};
\end{scope}

% --- ÉTAPE 2 : Bras lointain joue (Ajout pour clarté) ---
\begin{scope}[xshift=7.5cm, yshift=0cm]
\draw[ground] (-1, -0.5) -- (5, -0.5);
% Victime dos, bras proche en place
\draw[victim] (0.5, 0) ellipse (1.5 and 0.6);
\draw[victim] (0.5, 1.2) circle (0.6);
\draw[victim, line width=4pt] (-1.0, 0.8) -- (-1.0, 1.8) -- (-0.2, 1.8); % Bras proche
% Bras lointain amené à la joue
\draw[victim, line width=4pt, poporange] (2.0, 0.2) -- (1.2, 1.0) -- (0.8, 1.4); % Trajet main vers joue
\draw[victim, line width=4pt] (0.8, 1.4) -- (0.5, 1.6); % Main sur joue
% Secouriste maintient
\draw[popblue, line width=2pt] (-2.5, -0.2) -- (-2.5, -1.2) -- (-1.8, -1.2);
\draw[popblue, fill=popblue!20] (-2.5, 0.2) circle (0.3);
\node[stepnum] at (-3.5, 1.5) {2};
\node[labelbox, poporange, align=center] at (-3.5, 0.6){Bras lointain\\\textgreater\ Joue};
\draw[arrow, poporange] (-3.0, 1.0) .. controls (-1.5, 1.2) .. (0.8, 1.4);
\node[align=center, font=\small\bfseries, popdark] at (0.5, -2.5) {Étape 2 : Bras lointain\\contre la joue (protection)};
\end{scope}

% --- ÉTAPE 3-4 : Rotation (Levier genou) ---
\begin{scope}[xshift=15cm, yshift=0cm]
\draw[ground] (-1, -0.5) -- (5, -0.5);
% Victime en rotation (sur le flanc, pas tout à fait fini)
\draw[victim] (1.0, 0.2) ellipse (1.2 and 0.5); % Tronc de profil
\draw[victim] (1.0, 1.1) circle (0.5); % Tête
% Bras proche (au sol, antérieur)
\draw[victim, line width=4pt] (-0.2, 0.5) -- (0.5, 1.0); % Coude au sol
% Bras lointain (sous la joue)
\draw[victim, line width=4pt] (1.5, 1.0) -- (1.8, 1.4) -- (1.2, 1.6); % Main sous joue
% Jambe proche (tendue)
\draw[victim, line width=4pt] (0.5, -0.3) -- (0.5, -1.6);
% Jambe lointaine (pliiée, levier)
\draw[victim, line width=4pt] (1.8, -0.2) -- (2.5, 0.5) -- (2.0, 1.2); % Genou haut, pied au sol
% Secouriste tire le genou
\draw[popblue, line width=2pt] (3.0, 0.5) -- (3.0, -1.0) -- (3.7, -1.0);
\draw[popblue, fill=popblue!20] (3.0, 0.9) circle (0.3);
\node[stepnum] at (4.0, 1.5) {3-4};
\node[labelbox, popblue, align=center] at (4.0, 0.6){Tirer genou\\\textgreater\ Rotation};
\draw[arrow, decoration={snake, amplitude=2pt, segment length=5pt}, decorate] (2.5, 0.8) -- (1.5, 1.0) node[above left, font=\tiny] {Roulement};
\node[align=center, font=\small\bfseries, popdark] at (1.0, -2.5) {Étapes 3-4 : Plier jambe loin \& Basculer vers soi};
\end{scope}

% --- ÉTAPE 5 : PLS Finale ---
\begin{scope}[xshift=22.5cm, yshift=0cm]
\draw[ground] (-1, -0.5) -- (5, -0.5);
% Victime stable sur le côté
\draw[victim] (1.0, 0.2) ellipse (1.3 and 0.5); % Tronc
\draw[victim] (1.0, 1.1) circle (0.5); % Tête
% Bras antérieur (proche) stabilisateur
\draw[victim, line width=4pt] (-0.3, 0.4) -- (0.4, 0.9) -- (0.8, 0.5); % Coude au sol, main vers tête
% Bras postérieur (lointain) sous la joue
\draw[victim, line width=4pt] (1.5, 1.0) -- (1.8, 1.3) -- (1.3, 1.5); % Main joue
% Jambe antérieure (proche) tendue
\draw[victim, line width=4pt] (0.4, -0.3) -- (0.4, -1.6);
% Jambe postérieure (lointaine) pliée 90/90 (stable)
\draw[victim, line width=4pt] (1.8, -0.2) -- (2.4, 0.4) -- (2.0, 1.1); % Genou hanche 90
% Tête : bascule arrière, bouche ouverte
\draw[poporange, thick] (1.3, 1.3) -- (1.6, 1.1); % Ligne menton relevé
\draw[popgreen, thick, ->] (1.6, 1.0) -- (2.0, 0.8) node[right, font=\tiny, popgreen] {Écoulement};
% Secouriste surveille
\draw[popblue, line width=2pt] (3.0, 0.8) -- (3.0, -0.5) -- (3.7, -0.5);
\draw[popblue, fill=popblue!20] (3.0, 1.2) circle (0.3);
\node[stepnum] at (4.0, 2.0) {5};
\node[labelbox, popgreen, align=center] at (4.0, 1.2){Tête en arrière\\\textgreater\ Voies aériennes\\ouvertes + Bouche basse};
\node[align=center, font=\small\bfseries, popdark] at (1.0, -2.5) {Étape 5 : PLS Finale\\\& Surveillance continue};
\end{scope}
\end{tikzpicture}
\legende{Les 5 étapes clés de la mise en Position Latérale de Sécurité (PLS). \textbf{1.} Bras du côté du secouriste mis en angle droit (stabilisateur antérieur). \textbf{2.} Bras opposé amené contre la joue (protection tête). \textbf{3.} Jambe opposée pliée, pied au sol (levier). \textbf{4.} Basculement vers le secouriste en tirant le genou (la tête glisse sur la main). \textbf{5.} Ajustement : jambe supérieure fléchie hanche/genou \SI{90}{\degree} (stabilité), tête basculée en arrière, bouche ouverte vers le bas (gravité = drainage).}
\end{popfigure}

\courssub{Surveillance continue, alerte et conduite à tenir jusqu'à l'arrivée des secours}

Une fois la victime en PLS, le travail du secouriste n'est pas fini. Il entre dans une phase de \textbf{surveillance active}.

\begin{propriete}[Conduite à tenir post-PLS : « Protéger, Alerter, Surveiller »]
\begin{enumerate}
    \item \textbf{Protéger / Conforter :} Couvrir la victime avec une couverture de survie (côté argenté vers la victime par temps froid pour réfléchir la chaleur corporelle, côté doré vers la victime par temps très chaud pour réfléchir le soleil) ou tout vêtement/cover disponible. \textbf{Ne pas découvrir la victime} pour « voir si elle respire » : observer le thorax/abdomen à travers le tissu ou glisser une main discrètement.
    \item \textbf{Alerter (si pas déjà fait) :} Appeler le \textbf{112} (numéro d'urgence européen, fonctionne au Cameroun depuis mobile) ou le \textbf{119} (SAMU Cameroun) ou les \textbf{Sapeurs-Pompiers (\SI{118}{})}. Donner : \textbf{Lieu précis} (quartier, repères, étage, code porte), \textbf{Nombre de victimes}, \textbf{État} (« Inconsciente, respire, en PLS »), \textbf{Circonstances} (chute, malaise, maladie connue). \textbf{Ne pas raccrocher le premier.}
    \item \textbf{Surveiller la respiration \textbf{en continu} (toutes les \SI{1}{ à \SI{2}{minute}}) :} Vérifier visuellement (mouvements thoraco-abdominaux) et tactilement (souffle sur la joue/arrière de la main près de la bouche). \textbf{Ne pas quitter la victime des yeux.}
    \item \textbf{Communiquer / Empathie :} Lui parler calmement (\emph{« Les secours arrivent, je suis avec vous, tout va bien »}). Même inconsciente, l'audition est souvent le dernier sens à partir. Rassurer les témoins, écarter la foule.
    \item \textbf{Préparer l'arrivée des secours :} Dégager l'accès, allumer les lumières, désigner quelqu'un pour guider l'ambulance. Rassembler les ordonnances, carte d'identité, carte de groupe sanguin si disponibles.
\end{enumerate}
\end{propriete}

\begin{attention}[Erreurs fréquentes et interdictions formelles]
\begin{itemize}
    \item \textbf{Ne JAMAIS donner à boire ni à manger} à une victime inconsciente ou qui « revient à elle » (confusion, risque majeur de \cle{fausse route} $\rightarrow$ étouffement / pneumonie d'inhalation). Même si elle dit « j'ai soif ». Attendre l'avis médical (perfusion si besoin).
    \item \textbf{Ne JAMAIS laisser seule} une victime en PLS. La respiration peut s'arrêter à tout moment (aggravation AVC, arythmie, décompensation).
    \item \textbf{Ne PAS mettre en PLS} si suspicion de \textbf{traumatisme rachidien grave} (chute de hauteur, accident moto/vélo rapide, plongeon, AVP) \textbf{SAUF} si les voies aériennes sont menacées (vomissements, sang dans la bouche, obstruction par la langue) et qu'on ne peut pas les libérer autrement (aspiration, bascule tête/menton en maintenant l'alignement tête-cou-tronc). Dans ce cas : \cle{PLS modifiée} (rotation en bloc « comme un billot de bois » avec 2-3 secouristes) ou maintien en décubitus dorsal avec aspiration continue et bascule tête/menton prudente. \emph{Au Bac : la PLS standard est la règle pour l'inconscient qui respire ; la PLS modifiée est une connaissance de culture générale, pas une technique exigible en examen pratique seul.}
    \item \textbf{Ne pas mettre en PLS une femme enceinte de plus de \SI{20}{SA} (semaines d'aménorrhée) sur le côté droit :} L'utérus gravide comprime la \cle{veine cave inférieure} (retour veineux $\downarrow$ $\rightarrow$ choc hypotensif maternel $\rightarrow$ détresse fœtale). \textbf{Privilégier le côté GAUCHE} (PLS gauche) pour décomprimer la veine cave. Si côté gauche impossible (blessure), mettre en PLS droite mais avec un coussin/sous-vêtement sous le flanc droit pour basculer l'utérus vers la gauche.
\end{itemize}
\end{attention}

\begin{exoresolu}[Cas pratique : Jeune fille de 16 ans, malaise à l'école]
\textbf{Énoncé :} Au lycée de Nkolbisson, une élève de 16 ans s'effondre dans la cour de récréation. Elle est au sol, yeux fermés. Ses camarades l'entourent. Vous arrivez en premier. Elle ne répond pas à votre voix ni aux secousses d'épaules. Vous libérez les voies aériennes (tête en arrière, menton relevé). Vous observez : respiration présente, régulière, \SI{14}{cycles/min}, sans bruit anormal. Une camarade dit : « Elle n'a pas mangé ce matin, elle a eu ses règles abondantes hier. »
\textbf{Questions :}
\begin{enumerate}
    \item Quel est votre diagnostic immédiat ? Justifiez.
    \item Quelle est la première action à réaliser sur la victime ?
    \item Quelle alerte passez-vous et que dites-vous ?
    \item La victime reprend connaissance 5 min plus tard, elle est confuse, demande de l'eau. Que faites-vous ?
\end{enumerate}
\textbf{Solution détaillée :}
\begin{enumerate}
    \item \textbf{Perte de connaissance (malaise) avec respiration conservée.} Justification : Inconscience (pas de réponse stimulations) + Respiration normale (fréquence, rythme, amplitude normaux). Causes probables : \cle{syncope vasovagale} (déclenchée par la station debout, la chaleur, la faim) ou \cle{hypoglycémie} (jeûne) ou \cle{anémie} (règles abondantes). \textbf{Pas d'arrêt cardiaque.}
    \item \textbf{Mettre en Position Latérale de Sécurité (PLS)} immédiatement pour protéger les voies aériennes (langue, vomissements possibles).
    \item \textbf{Appeler le 112 (ou 119/118).} Message : « Lycée de Nkolbisson, cour principale. Une élève de 16 ans, inconsciente mais qui respire, mise en PLS. Probable malaise vagal / hypoglycémie (jeûne + règles). Pas de traumatisme. J'attends les secours. »
    \item \textbf{Ne pas lui donner d'eau.} La maintenir en PLS ou la rassurer en position assise \textbf{seulement si elle est parfaitement lucide, orientée (temps, lieu, identité), sans nausées ni vertiges}, et sous surveillance. Si confusion persiste : garder en PLS. Surveiller la respiration. Informer les secours de la reprise de conscience à leur arrivée.
\end{enumerate}
\end{exoresolu}

\courssub{Synthèse : la conduite à tenir du citoyen-secouriste face à l'inconscient}

\begin{aretenir}[Arbre décisionnel : Inconscient = Respiration ?]
\begin{center}
\begin{tikzpicture}[node distance=1.2cm and 1.5cm, every node/.style={font=\small, align=center},
    decision/.style={diamond, draw=popblue, thick, fill=popblue!10, aspect=2},
    action/.style={rectangle, draw=popgreen, thick, fill=popgreen!10, rounded corners},
    start/.style={ellipse, draw=popdark, thick, fill=popcream},
    stop/.style={rectangle, draw=poporange, thick, fill=poporange!10, rounded corners}]
\node[start, align=center] (start){Victime Inconsciente\\ (Ne répond pas)};
\node[action, below of=start, align=center] (airway){Libérer voies aériennes\\\text{Bascule tête / Relève menton}};
\node[decision, below of=airway, align=center] (resp){Respiration\\ NORMALE ?};
\node[action, left of=resp, xshift=-2.5cm, align=center] (pls){PLS \\+ Alerte 112/119/118 \\+ Surveiller respiration};
\node[action, right of=resp, xshift=2.5cm, align=center] (rcp){ARRÊT CARDIAQUE\\ Alerte 112/119/118 \\+ RCP (30/2) + DAE};
\node[stop, below of=pls, align=center] (surv){Surveillance continue\\ jusqu'aux secours\\ Si arrêt resp. $\rightarrow$ RCP};
\node[stop, below of=rcp, align=center] (cprcont){RCP sans interruption\\ jusqu'aux secours\\ ou reprise resp.};

\draw[->, thick] (start) -- (airway);
\draw[->, thick] (airway) -- (resp);
\draw[->, thick] (resp) -- node[above, font=\tiny, popgreen] {OUI} (pls);
\draw[->, thick] (resp) -- node[above, font=\tiny, poporange] {NON / AGONIQUE} (rcp);
\draw[->, thick] (pls) -- (surv);
\draw[->, thick] (rcp) -- (cprcont);
\end{tikzpicture}
\end{center}
\textbf{Règle d'or :} \textbf{« Inconscient qui respire = PLS. Inconscient qui ne respire pas = RCP. »} La vérification de la respiration (\SI{10}{\second}) est le \textbf{seul} critère de bascule. L'empathie, la parole, la couverture et la surveillance continue font partie intégrante du geste qui sauve.
\end{aretenir}

\begin{savaistu}[Le saviez-vous ? : La PLS, une invention française devenue standard mondial]
C'est le médecin militaire français \textbf{Dr Frédéric-Vincent \textsc{Larrey}} (1766-1842), chirurgien en chef de la Garde impériale de Napoléon, qui a le premier codifié le « \emph{tournement sur le côté} » pour les blessés inconscients sur les champs de bataille, afin d'éviter l'étouffement par le sang ou les vomissements. La technique a été affinée tout au long du XX\textsuperscript{e} siècle (Croix-Rouge, Protection Civile, SAMU) pour devenir la \textbf{Position Latérale de Sécurité (PLS)} enseignée aujourd'hui par le \emph{European Resuscitation Council (ERC)} et l'\emph{International Liaison Committee on Resuscitation (ILCOR)}. Au Cameroun, elle est enseignée dès le collège (Éducation Physique et Sportive, Éducation à la Santé) et est un prérequis pour le Brevet National de Sécurité et de Sauvetage Aquatique (BNSSA) et les formations de secourisme (PSC1, PSE1/2). Un geste simple, gratuit, qui ne demande que 30 secondes d'apprentissage et sauve des vies chaque jour dans nos rues, nos maisons et nos écoles.
\end{savaistu}

\begin{exercice}[Entraînement : Scénarios de décision]
\textbf{Pour chaque situation, cochez la conduite à tenir : (A) PLS + Alerte + Surveillance \quad (B) RCP + Alerte + DAE \quad (C) Surveillance simple (assise) + Alerte si aggravation.}
\begin{enumerate}
    \item Victime en état d'ébriété avancée, ronfle fort, ne répond pas, ventre qui monte/descend régulièrement.
    \item Victime au sol après choc électrique, ne répond pas, \textbf{pas de mouvement respiratoire}, peau bleutée (cyanose).
    \item Victime diabétique connue, trouvée confuse, assise, parle difficilement, transpire, tension artérielle basse. Répond à la voix.
    \item Victime épileptique, crise convulsive terminée depuis 2 min, yeux fermés, respiration bruyante (ronflement), ne répond pas.
    \item Victime AVC suspecté (bouche déviée, paralysie hémicorps), consciente, assise, respire normalement.
\end{enumerate}
\end{exercice}

\begin{corrige}[Exercice : Scénarios de décision]
\begin{enumerate}
    \item \textbf{(A) PLS.} Inconsciente + Respiration normale (même bruyante/ronflement = obstruction partielle par langue). PLS libère voies aériennes. Alerte 112/119/118. Surveiller.
    \item \textbf{(B) RCP.} Inconsciente + \textbf{Absence de respiration} = Arrêt cardiaque. RCP immédiate (\SI{30}{ compressions} / \SI{2}{ insufflations}). Alerte + DAE.
    \item \textbf{(C) Surveillance simple + Alerte.} \textbf{Consciente} (répond, orientée même si confuse). Pas de PLS (réservée à l'inconscient). Position demi-assise confortable. Alerte 112/119/118 (suspicions hypoglycémie / AVC). Surveiller niveau conscience : si devient inconsciente $\rightarrow$ PLS.
    \item \textbf{(A) PLS.} Post-critique épileptique : victime \textbf{inconsciente} (ne répond pas) + \textbf{Respiration présente} (bruyante = langue en arrière). PLS impérative pour drainer salive/sang et écarter langue. Alerte 112/119/118 si 1\textsuperscript{re} crise, crise $>$ \SI{5}{min}, traumatisme, grossesse, ou pas de reprise conscience en \SI{10}{min}.
    \item \textbf{(C) Surveillance simple + Alerte.} \textbf{Consciente}. Position demi-assise (tête surélevée \SI{30}{\degree}) pour réduire pression intracrânienne si AVC hémorragique. Alerte 112/119/118 urgente (unité neurovasculaire). \textbf{Pas de PLS.}
\end{enumerate}
\end{corrige}

\coursec{Synthèse : la conduite à tenir du citoyen-secouriste}

Après avoir détaillé, dans les sections précédentes, la prise en charge spécifique de la \cle{perte de connaissance} (reconnaissance, \cle{position latérale de sécurité}), de l'\cle{étouffement total} (nourrisson, adulte, femme enceinte) et de l'\cle{arrêt cardiaque} (\cle{RCP}, \cle{DAE}), il est indispensable de rassembler ces savoir-faire dans une \cle{conduite à tenir unique, logique et mémorisable}. En situation d'urgence réelle, le stress, le bruit et l'émotion perturbent le raisonnement. Seul un \cle{schéma décisionnel} solidement ancré, fruit d'un entraînement régulier, permet d'agir vite et juste. Cette synthèse vise à construire ce \emph{réflexe citoyen} en articulant protection, alerte et secours selon l'état de la victime.

\courssub{Le triptyque universel : Protéger — Alerter — Secourir}

Toute intervention de secourisme, quels que soient le contexte (domicile, voie publique, lieu de travail, établissement scolaire) et la pathologie, obéit à une \cle{chronologie immuable} : \textbf{Protéger}, puis \textbf{Alerter}, puis \textbf{Secourir}. Inverser cet ordre expose le secouriste à devenir une victime supplémentaire et retarde la prise en charge spécialisée.

\begin{popfigure}
\begin{tikzpicture}[
    node distance=2.5cm and 3cm,
    main/.style={draw=popblue, thick, fill=popblueL, rounded corners, minimum width=3.5cm, minimum height=1.5cm, align=center, font=\bfseries},
    arrow/.style={-Stealth, thick, popdark},
    labelarrow/.style={midway, fill=white, inner sep=2pt, font=\small, text=popdark}
]
\node[main, align=center] (proteger){\textbf{1. PROTÉGER}\\\textit{Sécuriser la zone}\\\textit{Éloigner le danger}\\\textit{Ne pas s'exposer}};
\node[main, right=of proteger, align=center] (alerter){\textbf{2. ALERTER}\\\textit{Appeler le 112 / 119}\\\textit{Message structuré}\\\textit{Ne pas raccrocher}};
\node[main, below=of alerter, align=center] (secourir){\textbf{3. SECOURIR}\\\textit{Gestes adaptés}\\\textit{RCP, PLS, Heimlich}\\\textit{Jusqu'à l'arrivée des secours}};

\draw[arrow] (proteger) -- node[labelarrow] {Si zone sûre} (alerter);
\draw[arrow] (alerter) -- node[labelarrow, right] {Alerte lancée} (secourir);
\draw[arrow, dashed, poporange] (secourir.west) to[out=180,in=-90] node[labelarrow, left=0.5cm] {Réévaluer en permanence} (proteger.south);
\end{tikzpicture}
\legende{Le triptyque Protéger — Alerter — Secourir : séquence logique et itérative. La flèche pointillée rappelle que la sécurité de la scène doit être réévaluée en permanence (incendie qui reprend, circulation qui reprend, victime qui se déplace).}
\end{popfigure}

\begin{methode}[Protéger : la priorité absolue]
\begin{enumerate}
    \item \textbf{Analyser la scène} : identifier les dangers (circulation, électricité, incendie, produits chimiques, foule, agressivité).
    \item \textbf{Supprimer le danger} si possible (couper le courant, éteindre un début d'incendie avec un extincteur adapté) \textbf{SANS SE METTRE EN DANGER}.
    \item \textbf{Éloigner la victime} seulement si le danger est imminent et impossible à supprimer (traîner par les vêtements en protégeant la tête, technique du \emph{traînage en abri}).
    \item \textbf{Baliser} la zone (triangle de pré-signalisation, lampes torches, témoins) pour protéger les autres usagers et les secours à venir.
\end{enumerate}
\end{methode}

\begin{attention}[Piège mortel : le « réflexe héros »]
\textbf{Se précipiter vers la victime sans avoir sécurisé les lieux est l'erreur n°1.} Un secouriste électrocuté, percuté par un véhicule ou intoxiqué par des fumées \textbf{ne peut plus secourir personne} et double la charge des secours. \textbf{Exemple} : sur la route Yaoundé-Douala, un témoin traverse la chaussée pour porter secours à un accidenté sans poser de triangle ; il est percuté par un poids lourd. Bilan : deux victimes graves au lieu d'une. \textit{Protéger, c'est déjà secourir.}
\end{attention}

\begin{methode}[Alerter : le message type « PASSE »]
Dès la zone sécurisée, alerter \textbf{avant} de secourir (règle générale). \textbf{Exceptions} : arrêt cardiaque non témoin (noyade, enfant, asphyxie) ou témoin seul sans téléphone → 1 min de RCP puis alerter. Composer le \textbf{112} (numéro d'urgence européen, gratuit, fonctionne sans crédit, même clavier verrouillé) ou le \textbf{119} (Cameroun, pompiers/SAMU). Suivre la structure \cle{PASSE} :
\begin{enumerate}
    \item \textbf{P} — \textbf{P}our qui ? (Votre identité, numéro de rappel \textbf{impératif}).
    \item \textbf{A} — \textbf{A}dresse exacte (rue, n°, étage, porte, code d'accès, \textbf{repères visibles} : « en face de la station Total, au carrefour de l'Amitié »).
    \item \textbf{S} — \textbf{S}ituation / \textbf{S}igne apparent (nature : arrêt cardiaque, étouffement, AVC, traumatisme ; nombre de victimes ; état : conscient/inconscient, respire/ne respire pas, saigne).
    \item \textbf{S} — \textbf{S}oins déjà prodigués (RCP en cours, PLS, compression hémorragie, DAE posé).
    \item \textbf{E} — \textbf{E}ntretenir la ligne : \textbf{NE JAMAIS RACROCHER LE PREMIER}. Le régulateur peut guider la RCP, localiser l'appelant, déclencher les secours adaptés (SMUR, VSAV, hélicoptère).
\end{enumerate}
\end{methode}

\begin{exemplebox}[Message d'alerte efficace — Exemple réel]
\textbf{Appelant :} « Allô, 112 ? Bonjour, je suis M. \textsc{Fouda}, 6 99 12 34 56. Je suis au quartier \textsc{Essos}, carrefour de l'\textsc{Amitié}, devant l'entrée du lycée \textsc{Leclerc}. Un homme d'environ 50 ans vient de s'effondrer en marchant. Il est \textbf{inconscient}, \textbf{ne respire pas} (pas de mouvement thoracique, pas de souffle sur la joue). Je commence la \textbf{RCP} (30 compressions / 2 insufflations). Il n'y a pas de DAE visible. Je reste en ligne. »
\textbf{Régulateur :} « Reçu, l'ambulance SMUR part de l'Hôpital Central dans 2 min. Continuez la RCP, je vous guide pour le rythme. Ne raccrochez pas. »
\end{exemplebox}

\courssub{L'arbre décisionnel de synthèse : de l'observation au geste adapté}

Face à une victime, le secouriste doit réaliser un \cle{bilan rapide} (conscience, respiration) en \textbf{moins de 10 secondes} pour orienter sa conduite. L'arbre ci-dessous synthétise les quatre scénarios majeurs.

\begin{popfigure}
\begin{tikzpicture}[
    node distance=1.8cm and 2.2cm,
    decision/.style={diamond, draw=poppurple, thick, fill=poppurpleL, aspect=2, inner sep=4pt, align=center, font=\small\bfseries},
    action/.style={rectangle, draw=popgreen, thick, fill=popgreenL, rounded corners, minimum width=3.2cm, minimum height=1.2cm, align=center, font=\small\bfseries},
    start/.style={ellipse, draw=popblue, thick, fill=popblueL, minimum width=2.5cm, minimum height=1.2cm, align=center, font=\bfseries},
    arrow/.style={-Stealth, thick, popdark},
    label/.style={midway, fill=white, inner sep=1.5pt, font=\scriptsize, text=popdark, align=center}
]
\node[start, align=center] (start){VICTIME\\ DÉCOUVERTE};
\node[decision, below=of start, align=center] (conscience){CONSCIENTE ?\\(Réagit à la voix,\\au toucher)};
\node[decision, below=of conscience, align=center] (respiration){RESPIRE ?\\(Mouvements thoraciques,\\souffle sur joue\\< 10 s)};
\node[decision, right=of conscience, xshift=1.5cm, align=center] (etouffe){ÉTOUFFEMENT ?\\(Ne peut parler,\\tient la gorge,\\sifflement)};
\node[action, left=of respiration, xshift=-1.5cm, align=center] (pls){\textbf{PLS}\\\textit{Position Latérale  de Sécurité}\\\textit{Surveiller respiration  Alerter 112}};
\node[action, right=of respiration, xshift=1.5cm, align=center] (rcp){\textbf{RCP + DAE}\\\textit{30 compressions / 2 ins.  100-120/min, 5-6 cm  DAE dès dispo}};
\node[action, below=of etouffe, align=center] (heimlich){\textbf{CLAQUES + HEIMLICH}\\\textit{Adulte/Enfant : 5 claques  + 5 compressions abdo  Femme enceinte/obèse :  compressions thoraciques  Nourrisson : 5 claques dos  + 5 compressions thoraciques}};
\node[action, right=of etouffe, xshift=1cm, align=center] (toux){\textbf{ENCOURAGER LA TOUX}\\\textit{Ne rien faire d'autre  Surveiller aggravation}};

\draw[arrow] (start) -- (conscience);
\draw[arrow] (conscience) -- node[label, left] {NON} (respiration);
\draw[arrow] (conscience) -- node[label, above] {OUI} (etouffe);
\draw[arrow] (respiration) -- node[label, above] {OUI} (pls);
\draw[arrow] (respiration) -- node[label, above] {NON} (rcp);
\draw[arrow] (etouffe) -- node[label, above, align=center]{OUI \\ (Ne parle plus)} (heimlich);
\draw[arrow] (etouffe) -- node[label, above, align=center]{NON \\ (Tousse efficacement)} (toux);
\end{tikzpicture}
\legende{Arbre décisionnel unique du citoyen-secouriste. Deux entrées : Conscience (OUI/NON) → si NON : Respiration (OUI/NON). Si Conscience = OUI : rechercher signe d'étouffement total (impossibilité de parler, toux inefficace, signe de la main à la gorge). Chaque feuille correspond à une conduite à tenir standardisée.}
\end{popfigure}

\begin{aretenir}[Règle d'or du bilan respiratoire]
La vérification de la respiration (regarder le thorax, écouter/ressentir le souffle) ne doit \textbf{pas dépasser 10 secondes}. En cas de doute \textbf{considérer que la victime NE RESPIRE PAS} et démarrer la RCP. Mieux vaut faire une RCP sur un cœur qui bat (risque faible de lésions costales) que de ne pas en faire sur un cœur arrêté (mort certaine).
\end{aretenir}

\begin{exoresolu}[Application de l'arbre décisionnel : trois cas concrets]
\textbf{Cas 1 :} Au marché Mokolo, une femme de 35 ans s'effondre brutalement. Vous vous approchez (zone sûre), lui tapez les épaules : « Madame, m'entendez-vous ? » $\rightarrow$ \textbf{Pas de réponse}. Vous penchez la tête, relevez le menton, regardez le thorax 10 s : \textbf{pas de mouvement, pas de souffle}. \textbf{Décision :} Alerter 112 (ou envoyer un témoin), demander un DAE, \textbf{démarrer RCP immédiatement} (30 compressions / 2 ins.).

\textbf{Cas 2 :} Au restaurant, un homme de 60 ans mangeant un morceau de viande se lève brutalement, porte les mains à la gorge (signe universel), ne peut pas parler, devient cyanotique. Il est \textbf{conscient} mais \textbf{ne peut pas tousser}. \textbf{Décision :} \textbf{5 claques inter-omoplatiques} (buste penché en avant) + \textbf{5 compressions abdominales (Heimlich)}. Répéter jusqu'à expulsion ou perte de conscience (alors $\rightarrow$ RCP).

\textbf{Cas 3 :} Dans un bus, un élève de Terminale s'endort, tête sur l'épaule du voisin. Vous le secouez doucement : il grogne, ouvre les yeux, vous répond. \textbf{Conscient}. Il respire normalement (thorax se soulève). \textbf{Décision :} \textbf{PLS} pour maintenir les voies aériennes libres (langue, vomissements) et surveiller jusqu'à l'arrivée ou la reprise complète.
\tcblower
\textbf{Analyse :} Dans chaque cas, le triptyque \textbf{Protéger-Alerter-Secourir} est respecté. L'arbre décisionnel évite l'hésitation : « Conscient ? $\rightarrow$ Non $\rightarrow$ Respire ? $\rightarrow$ Non $\rightarrow$ RCP ». La rapidité du bilan ($< 10$ s) est critique : chaque minute sans RCP réduit la survie de 10 \%.
\end{exoresolu}

\courssub{Protection juridique et morale : le devoir d'agir}

Au-delà de la technique, le secourisme engage la \cle{responsabilité citoyenne}. La loi (Code pénal camerounais, art. 277 ; lois similaires dans de nombreux pays) sanctionne la \cle{non-assistance à personne en danger}.

\begin{propriete}[Obligation d'assistance — Savoir admis au Baccalauréat]
\textbf{Toute personne qui, sans risque pour elle ni pour les tiers, peut empêcher par son action immédiate un crime ou un délit contre l'intégrité corporelle d'une personne, ou porter secours à une personne en péril, est tenue de le faire.} À défaut, elle encourt des sanctions pénales (emprisonnement, amende).
\textbf{Corollaire (Protection du « Bon Samaritain ») :} Le secouriste qui agit \textbf{conformément aux gestes enseignés} (formation validée, respect des protocoles) \textbf{n'engage pas sa responsabilité civile ni pénale} en cas d'aggravation ou d'échec, même s'il commet une maladresse (ex. : fracture costale lors d'une RCP bien réalisée). La loi protège l'intention de sauver.
\end{propriete}

\begin{attention}[Limites de la protection juridique]
La protection tombe si le secouriste :
\begin{itemize}
    \item Agit \textbf{avec imprudence caractérisée} (ex. : déplacer un traumatisé rachidien sans indication vitale, pratiquer une trachéotomie au cutter sans compétence).
    \item Ne \textbf{respecte pas les protocoles} enseignés (ex. : faire un massage cardiaque sur un lit mou au lieu de durcir le plan, ne pas appeler les secours).
    \item Met \textbf{volontairement} la victime en danger.
\end{itemize}
\textit{La formation régulière (recyclage annuel recommandé) est la meilleure garantie juridique et éthique.}
\end{attention}

\courssub{Dimension citoyenne et empathique : au-delà du geste technique}

Le secouriste n'est pas un robot exécutant un algorithme. Il est un \cle{être humain} face à un autre être humain en détresse. La \cle{qualité de la relation} influence l'issue (stress de la victime, confiance, acceptation des soins) et le vécu du secouriste lui-même.

\begin{definition}[Empathie en secourisme]
Capacité à \textbf{se mettre à la place de la victime} pour adapter sa communication (ton calme, mots simples, contact physique rassurant — main sur l'épaule, regard), \textbf{respecter sa dignité} (couvrir, écarter les curieux, ne pas juger) et \textbf{accompagner ses proches} (expliquer ce qui se passe, les orienter, les protéger de la vue de la scène si nécessaire).
\end{definition}

\begin{methode}[Conduite relationnelle : 4 piliers]
\begin{enumerate}
    \item \textbf{Se présenter} : « Bonjour, je m'appelle [Prénom], je suis formé aux premiers secours, je vais vous aider. »
    \item \textbf{Expliquer chaque geste} avant de le faire : « Je vais vous mettre sur le côté pour que vous respiriez mieux. »
    \item \textbf{Rassurer sans mentir} : « Les secours arrivent, je reste avec vous. » (Pas de « Tout va bien » si l'état est grave).
    \item \textbf{Transmettre aux secours} : À l'arrivée des pompiers/SAMU, faire un \textbf{compte-rendu oral structuré} (Identité, Circonstances, Bilan initial, Gestes faits, Évolution). Remplir la \cle{fiche de liaison} si existante.
\end{enumerate}
\end{methode}

\begin{savaistu}[L'impact sociétal de la formation grand public — Le modèle danois]
Au \textbf{Danemark}, une stratégie nationale (formation obligatoire à l'école, au permis de conduire, dans l'armée, DAE publics massifs, application smartphone alertant les secouristes volontaires à proximité) a fait passer le taux de RCP par témoin de \textbf{22 \% (2001) à 45 \% (2014)} et la survie à 30 jours après arrêt cardiaque extra-hospitalier de \textbf{3,5 \% à 11 \%} (multiplié par \textbf{3,1}). Au \textbf{Cameroun}, l'intégration du secourisme au programme de Terminale D (4 h) et la multiplication des formations Croix-Rouge, SAMU, ONG (ex. \textit{Heart Foundation Cameroon}) visent le même objectif : faire de \textbf{chaque citoyen le premier maillon de la chaîne de survie}. Un élève de Terminale D formé aujourd'hui peut sauver son père, sa mère, un camarade, un inconnu demain.
\end{savaistu}

\courssub{Prévention des accidents cardiovasculaires : rappel des facteurs de risque (complément Bac)}

Le secourisme est la \textbf{dernière barrière} (prévention tertiaire). La \textbf{meilleure prise en charge} reste la \cle{prévention primaire} : éviter l'accident. Le programme de SVT (Module I : Santé) et cette leçon convergent vers les mêmes messages de santé publique, fondés sur l'épidémiologie (Étude INTERHEART, OMS).

\begin{exemplebox}[Synthèse chiffrée des facteurs de risque modifiables — Cibles pour le Bac]
\begin{center}
\begin{tabularx}{\linewidth}{|l|X|X|}\hline
\textbf{Facteur de risque} & \textbf{Mécanisme pathophysiologique (rappel SVT)} & \textbf{Message de prévention concret (contexte Cameroun)} \\ \hline
\textbf{HTA non contrôlée} & Surcharge de post-charge $\rightarrow$ hypertrophie VG $\rightarrow$ insuffisance cardiaque, rupture anévrisme, AVC hémorragique. & \textbf{« Connaissez votre tension ! »} Dépistage gratuit dans les centres de santé (CS) et hôpitaux de district. Traitement observance à vie. Réduire le sel (cube Maggi, poisson séché, sauce arachide salée). \\ \hline
\textbf{Diabète type 2} & Hyperglycémie chronique $\rightarrow$ athérosclérose accélérée (macro/micro-angiopathie) $\rightarrow$ IDM silencieux, AVC, artérite membres inférieurs. & Dépistage glycémie à jeun (CS). Alimentation : réduire sucres rapides (boissons gazeuses, beignets), privilégier igname, plantain, légumes verts. Activité physique 30 min/jour. \\ \hline
\textbf{Tabagisme} & CO (hypoxie), nicotine (vasoconstriction, agrégation plaquettaire), goudrons (endothélium toxique) $\rightarrow$ thrombose coronaire/cérébrale. & \textbf{Arrêt total.} Sevrage aidé (CS, pneumologues). Interdiction fumer lieux publics (loi 2006). Chicha = danger égal cigarette. \\ \hline
\textbf{Dyslipidémie} & $LDL\text{-}c$ élevé $\rightarrow$ plaque d'athérome instable $\rightarrow$ rupture $\rightarrow$ thrombose occlusive (IDM, AVC ischémique). & Limiter graisses saturées (viande rouge, peau poulet, huile de palme non traitée, fritures). Privilégier poisson (sardine, maquereau), huiles végétales (arachide, coton), fibres (légumineuses, fruits). \\ \hline
\textbf{Sédentarité / Obésité} & Résistance à l'insuline, HTA, dyslipidémie, inflammation chronique. IMC $> 25$ kg/m$^2$ = risque majoré. & Marcher 30 min/jour (aller à pied au travail/marché, monter escaliers). Lutte contre obésité infantile (cantines scolaires équilibrées, éducation nutritionnelle). \\ \hline
\textbf{Stress / Alcool} & Cortisol/catécholamines chroniques $\rightarrow$ HTA, arythmies. Alcool excessif $\rightarrow$ cardiomyopathie dilatée, HTA, AVC. & Gestion du stress (sommeil 7-8 h, pause, soutien social). Alcool : max 2 verres/jour (homme), 1 (femme), jours d'abstinence. \\ \hline
\end{tabularx}
\end{center}
\legende{Principaux facteurs de risque cardiovasculaires modifiables. La maîtrise de ces facteurs (prévention primaire) réduit de 80 \% le risque d'IDM et d'AVC (OMS). L'élève de Terminale D doit pouvoir citer les mécanismes (liens avec physiologie cardiaque, vasculaire, métabolique) et proposer des actions de prévention adaptées au contexte local.}
\end{exemplebox}

\begin{exercice}[Entraînement à la synthèse — Sujet type Bac]
\textbf{Énoncé :} Un homme de 55 ans, connu hypertendu non suivi, fumeur, s'effondre au travail. Un collègue (formé secourisme) intervient.
\begin{enumerate}
    \item Citez, dans l'ordre, les trois temps de la conduite à tenir universelle.
    \item Le collègue vérifie la conscience (pas de réponse) et la respiration (absente). Quelle conduite tenir ? Justifiez par l'arbre décisionnel.
    \item Détaillez les paramètres chiffrés de la RCP (fréquence, profondeur, ratio, changement de secouriste).
    \item Le DAE arrive. Citez les 4 étapes de son utilisation.
    \item Citez 3 facteurs de risque cardiovasculaires présents chez cette victime et, pour chacun, le mécanisme physiopathologique principal.
\end{enumerate}
\end{exercice}

\begin{corrige}[Exercice n]
\begin{enumerate}
    \item \textbf{1. Protéger} (sécuriser zone, ne pas s'exposer), \textbf{2. Alerter} (112/119, message PASSE, ne pas raccrocher), \textbf{3. Secourir} (gestes adaptés selon bilan).
    \item \textbf{Victime Inconsciente + Ne respire pas} $\rightarrow$ \textbf{Arrêt cardiaque} $\rightarrow$ \textbf{RCP immédiate + DAE dès que possible}. (Arbre : Consciente ? NON $\rightarrow$ Respire ? NON $\rightarrow$ RCP).
    \item Fréquence : \textbf{100 à 120 compressions/min}. Profondeur : \textbf{5 à 6 cm} (adulte). Ratio : \textbf{30 compressions / 2 insufflations}. Changement de secouriste : \textbf{toutes les 2 min} (ou 5 cycles) pour éviter la fatigue $\rightarrow$ baisse qualité.
    \item 1. Allumer DAE (bouton ON/OUVRIR). 2. Coller électrodes (schéma sur boîtier : sternum droit / apex gauche). 3. Laisser analyser (\textbf{NE TOUCHER PLUS LA VICTIME}). 4. Si choc conseillé : \textbf{Éloigner tout le monde} $\rightarrow$ Appuyer bouton CHOC $\rightarrow$ Reprendre RCP immédiatement 2 min. Si pas de choc : Reprendre RCP 2 min.
    \item \begin{itemize}
        \item \textbf{HTA} : Surcharge de post-charge $\rightarrow$ hypertrophie VG $\rightarrow$ rupture plaque / anévrisme / IDM / AVC.
        \item \textbf{Tabagisme} : Endothélium dysfonctionnel + hypercoagulabilité + vasoconstriction $\rightarrow$ thrombose coronaire aiguë (IDM).
        \item \textbf{Non observance traitement} : Perte de contrôle tensionnel $\rightarrow$ pic tensionnel brutal $\rightarrow$ rupture plaque athéromateuse ou AVC hémorragique.
    \end{itemize}
\end{enumerate}
\end{corrige}

\begin{aretenir}[Check-list finale du citoyen-secouriste Terminale D]
\begin{itemize}
    \item \textbf{Je sais} : Protéger-Alerter-Secourir (ordre immuable).
    \item \textbf{Je sais} : Faire le bilan (Conscience, Respiration) en $< 10$ s.
    \item \textbf{Je sais} : L'arbre décisionnel $\rightarrow$ 4 issues : Toux encouragée / Claques+Heimlich / PLS / RCP+DAE.
    \item \textbf{Je sais} : RCP = 100-120/min, 5-6 cm, 30/2, changer ttes 2 min. DAE = Allumer, Coller, Analyser, Choc/Non-choc.
    \item \textbf{Je sais} : Message d'alerte PASSE (Numéro de rappel ! Ne pas raccrocher).
    \item \textbf{Je sais} : Devoir légal d'assistance + Protection si gestes conformes.
    \item \textbf{Je sais} : Empathie, dignité, transmission aux secours.
    \item \textbf{Je sais} : Facteurs de risque CV modifiables (HTA, Diabète, Tabac, Lipides, Sédentarité) + Mécanismes + Prévention.
\end{itemize}
\textbf{Cet ensemble de compétences, c'est le pouvoir de sauver une vie. Entraînez-vous, recylez-vous, restez citoyen-secouriste.}
\end{aretenir}

\newpage

\coursec{Activité d'intégration}

\courssub{Objectifs de l'activité}

À la fin de cette activité d'intégration, tu seras capable de :
\begin{itemize}
\item interpréter une courbe de survie en fonction du délai d'intervention et en tirer des conclusions chiffrées sur l'urgence d'agir ;
\item poser un \cle{diagnostic d'urgence} (arrêt cardiaque, étouffement total, perte de connaissance) à partir de signes cliniques simples ;
\item rédiger, pour chaque situation, la \cle{séquence complète des gestes qui sauvent} avec leurs paramètres chiffrés (fréquence, profondeur, nombre de cycles) ;
\item rédiger un \cle{message d'alerte} téléphonique complet et hiérarchisé ;
\item décrire et ordonner les étapes de la \cle{position latérale de sécurité} (PLS) ;
\item identifier les erreurs fréquentes à éviter et argumenter sur le rôle du citoyen-secouriste dans la chaîne de survie au Cameroun.
\end{itemize}

\courssub{Situation}

\textbf{Urgence au marché Mokolo de Yaoundé : devenez le maillon de la chaîne de survie.}

Il est 11 h 30 un samedi matin au marché Mokolo, l'un des plus grands marchés de Yaoundé. La foule est dense, la chaleur accablante (34 \textdegree C). Tu fais tes courses avec ta classe de Terminale D dans le cadre d'une sortie pédagogique sur l'hygiène alimentaire, encadrée par ton professeur de SVT et une infirmière du Centre de Santé Intégré du quartier. En l'espace de quarante minutes, trois événements se produisent.

\textbf{Document 1 — La courbe de survie.} L'infirmière vous montre un document portant sur la probabilité de survie sans séquelle après un arrêt cardiaque, en fonction du délai de début de la réanimation cardio-pulmonaire (RCP) par un témoin.

\begin{popfigure}
\begin{center}
\begin{tikzpicture}
\begin{axis}[
width=0.85\linewidth, height=6cm,
xlabel={Délai avant le début de la RCP (min)},
ylabel={Probabilité de survie (\%)},
xmin=0, xmax=12, ymin=0, ymax=80,
xtick={0,1,2,3,4,5,6,7,8,9,10,11,12},
ytick={0,10,20,30,40,50,60,70,80},
grid=both, grid style={popline!40},
legend style={at={(0.97,0.95)}, anchor=north east, font=\small}
]
\addplot[popcurve, domain=0:12, samples=100] {70*exp(-0.23*x)};
\addlegendentry{Survie sans RCP immédiate}
\addplot[poporange, very thick, dashed, domain=0:12, samples=100] {70*exp(-0.05*x)};
\addlegendentry{Survie si RCP débutée dans la minute}
\draw[popgreen, thick, dashed] (axis cs:1,0) -- (axis cs:1,56);
\draw[popgreen, thick, dashed] (axis cs:8,0) -- (axis cs:8,11);
\node[popgreen, font=\small, anchor=west] at (axis cs:1.1,60) {1 min : $\approx$ 56 \%};
\node[popgreen, font=\small, anchor=west] at (axis cs:8.1,18) {8 min : $\approx$ 11 \%};
\end{axis}
\end{tikzpicture}
\end{center}
\legende{Probabilité de survie après arrêt cardiaque en fonction du délai de début de la RCP. Sans geste immédiat, les chances de survie s'effondrent : chaque minute perdue en retranche environ 7 à 10 \%.}
\end{popfigure}

\textbf{Document 2 — Trois cas cliniques.}
\begin{itemize}
\item \textbf{Cas (a).} Un vendeur de tomates de 55 ans s'effondre brutalement devant son étal. Il ne réagit ni à la voix ni aux stimulations. Son thorax est immobile : il ne respire pas.
\item \textbf{Cas (b).} Une maman affolée arrive en courant avec son nourrisson de 8 mois : l'enfant, qui jouait avec des arachides grillées, est devenu bleu (cyanosé), silencieux, et ne tousse plus.
\item \textbf{Cas (c).} Au restaurant « Chez Maman Ngo Bell » attenant au marché, une femme enceinte de 7 mois se lève brusquement, les deux mains sur la gorge, incapable de parler, de tousser ou de respirer.
\end{itemize}

\textbf{Document 3.} En sortant du restaurant, vous découvrez un motard (un « bendskineur ») allongé sur le dos après une chute sans traumatisme apparent. Il ne répond pas quand on lui parle, mais son thorax se soulève régulièrement : il respire. Personne n'a encore rien fait et il commence à émettre des bruits de gargouillis.

\courssub{Consignes}

Travail en groupes de 4 élèves. Chaque groupe rédige ses réponses sur une fiche, puis la restitution se fait par simulation (sur mannequin ou camarade volontaire, \textbf{sans aucune pression réelle} sur le thorax ou l'abdomen).

\begin{enumerate}
\item \textbf{Document 1.} À l'aide de la courbe, calcule la perte de chances de survie entre une RCP débutée à 1 minute et une RCP débutée à 8 minutes. Exprime cette perte en points de pourcentage, puis en proportion relative. Quelle conclusion pratique en tires-tu sur le rôle du témoin ?
\item \textbf{Cas (a).} Pose le diagnostic d'urgence en justifiant par les signes observés. Rédige la séquence complète des gestes qui sauvent, avec tous les paramètres chiffrés (fréquence, profondeur, rythme compressions/ventilations, durée des cycles). Rédige le message d'alerte téléphonique que tu adresserais au SAMU Cameroun (119) ou aux pompiers (118).
\item \textbf{Cas (b).} Pose le diagnostic. Décris la conduite à tenir chez ce nourrisson : position, nombre et localisation des claques dans le dos, puis des compressions thoraciques. Cite deux erreurs à éviter absolument.
\item \textbf{Cas (c).} Pose le diagnostic. Explique pourquoi la manœuvre de Heimlich classique est contre-indiquée ici, et décris la technique adaptée à la femme enceinte.
\item \textbf{Document 3.} Quel est le diagnostic ? Décris et ordonne les étapes de la position latérale de sécurité (PLS), en expliquant son intérêt physiologique face aux bruits de gargouillis.
\item \textbf{Débat.} Au Cameroun, les délais d'arrivée des secours organisés peuvent dépasser 30 minutes en ville. Argumente, en t'appuyant sur le document 1, l'importance de former la population aux gestes qui sauvent, et propose deux actions concrètes pour ton lycée ou ton quartier.
\end{enumerate}

\courssub{Ressources à mobiliser}

\begin{aretenir}[Les réflexes du secouriste]
\begin{itemize}
\item \textbf{Chaîne de survie} : alerter $\rightarrow$ masser $\rightarrow$ défibriller $\rightarrow$ secours médicalisés. Chaque minute sans RCP fait perdre environ 7 à 10 \% de chances de survie.
\item \textbf{Arrêt cardiaque} : victime inconsciente + absence de respiration (ou respiration agonique : rares soupirs irréguliers) $\Rightarrow$ RCP : compressions thoraciques au centre de la poitrine, fréquence de 100 à 120 compressions par minute, profondeur de 5 à \SI{6}{cm} chez l'adulte, alternance 30 compressions / 2 ventilations, sans interruption jusqu'à l'arrivée des secours.
\item \textbf{Étouffement total} (victime silencieuse, incapable de tousser) : 5 claques dans le dos puis 5 compressions (abdominales = manœuvre de Heimlich chez l'adulte ; thoraciques chez la femme enceinte et le nourrisson), en alternance.
\item \textbf{Perte de connaissance avec respiration présente} : position latérale de sécurité, surveillance de la respiration, alerte.
\item \textbf{Message d'alerte} : qui appelle, quoi (nature de l'accident), où (adresse précise), combien de victimes, état de la victime, gestes déjà réalisés ; ne jamais raccrocher le premier.
\end{itemize}
\end{aretenir}

\courssub{Démarche et corrigé commenté}

\begin{exoresolu}[Corrigé commenté de l'activité]
\textbf{1. Exploitation de la courbe de survie.}

Lecture graphique sur la courbe « sans RCP immédiate » : à 1 minute, la probabilité de survie est d'environ 56 \% ; à 8 minutes, elle tombe à environ 11 \%.
\begin{align*}
\text{Perte absolue} &= 56\ \% - 11\ \% = 45\ \text{points de pourcentage}\\
\text{Perte relative} &= \frac{56 - 11}{56} \approx 0,80 \approx 80\ \%
\end{align*}
\tcblower
\textbf{Interprétation.} Attendre 8 minutes au lieu d'agir dans la première minute, c'est perdre 45 points de survie, soit environ 80 \% des chances initiales de la victime. La courbe décroît d'environ 7 à 10 \% par minute : le cerveau, privé d'oxygène, subit des lésions irréversibles au-delà de 3 à 5 minutes d'arrêt circulatoire. En revanche, si la RCP est débutée dans la minute (courbe en pointillés orange), la décroissance est beaucoup plus lente : le massage maintient un débit sanguin minimal vers le cerveau et le cœur. \textbf{Conclusion pratique :} le témoin est le premier maillon de la chaîne de survie ; c'est lui, et non l'hôpital, qui détermine le pronostic dans les premières minutes.

\textbf{2. Cas (a) — l'arrêt cardiaque.}

\emph{Diagnostic.} Victime inconsciente (ne réagit ni à la voix ni à la stimulation) ET sans respiration (thorax immobile) : c'est un \cle{arrêt cardiaque}. Le cœur ne pompe plus ; sans RCP, la mort survient en quelques minutes.

\emph{Séquence des gestes.}
\begin{enumerate}
\item \textbf{Protéger} : écarter la foule, sécuriser la zone (pas de danger électrique ni de circulation).
\item \textbf{Alerter} : désigner une personne précise (« Toi, appelle le 119 ! ») ; si l'on est seul, téléphoner en haut-parleur tout en commençant le massage.
\item \textbf{Masser} : placer le talon d'une main au centre de la poitrine (sur la moitié inférieure du sternum), l'autre main par-dessus, doigts entrelacés, bras tendus, épaules au-dessus des mains. Compressions de 5 à \SI{6}{cm} de profondeur, à une fréquence de 100 à 120 par minute, en laissant le thorax se regonfler complètement entre deux compressions.
\item \textbf{Ventiler} (si l'on est formé) : bascule de la tête en arrière, 30 compressions puis 2 insufflations d'environ 1 seconde chacune, en cycles ininterrompus. À défaut, poursuivre les compressions seules.
\item \textbf{Ne pas arrêter} jusqu'à l'arrivée des secours, la reprise d'une respiration normale, ou l'épuisement total du secouriste (relais possible toutes les 2 minutes).
\end{enumerate}

\emph{Message d'alerte.} « Allô SAMU ? Je suis témoin au marché Mokolo, entrée principale côté avenue, à Yaoundé. Un homme de 55 ans est inconscient et ne respire plus : arrêt cardiaque. Nous avons commencé le massage cardiaque. Envoyez une ambulance, je reste en ligne. »

\emph{Erreurs à éviter} : chercher le pouls (perte de temps, peu fiable pour un non-professionnel), mettre la victime en PLS (elle ne respire pas : il faut masser !), arrêter le massage pour téléphoner, donner à boire.

\textbf{3. Cas (b) — l'étouffement total du nourrisson.}

\emph{Diagnostic.} Nourrisson cyanosé (bleu), silencieux, sans toux après inhalation d'arachide : obstruction \textbf{totale} des voies aériennes. (Si l'enfant toussait encore, on le laisserait tousser sans intervenir, la toux étant le meilleur mécanisme d'expulsion.)

\emph{Conduite à tenir.} Alerter le 119. Allonger le nourrisson à plat ventre sur l'avant-bras, tête plus basse que le tronc, mâchoire maintenue ; porter \textbf{5 claques fermes dans le dos} entre les omoplates avec le talon de la main. Si échec, retourner l'enfant sur le dos (tête toujours déclive) et porter \textbf{5 compressions thoraciques} avec deux doigts au centre de la poitrine. Alterner 5 claques / 5 compressions jusqu'à l'expulsion du corps étranger ou l'arrivée des secours. Si l'enfant perd connaissance, débuter la RCP.

\emph{Erreurs à éviter} : introduire un doigt à l'aveugle dans la bouche (risque d'enfoncer l'arachide) ; secouer l'enfant par les pieds la tête en bas.

\textbf{4. Cas (c) — l'étouffement de la femme enceinte.}

\emph{Diagnostic.} Victime silencieuse, mains sur la gorge (signe universel d'étouffement), incapable de parler : obstruction totale des voies aériennes.

\emph{Conduite à tenir.} 5 claques dans le dos entre les omoplates, victime penchée en avant. Si échec, la manœuvre de Heimlich \textbf{abdominale} est contre-indiquée : l'utérus gravide (7 mois) occupe l'abdomen et une pression à ce niveau risquerait de léser le fœtus et l'utérus. On réalise donc des \textbf{compressions thoraciques} : se placer derrière la victime, passer les bras sous ses aisselles, poser le poing fermé au centre du sternum (et non au creux épigastrique), le saisir avec l'autre main, et tirer fermement vers soi, 5 fois. Alterner 5 claques / 5 compressions thoraciques.

\textbf{5. Document 3 — la perte de connaissance.}

\emph{Diagnostic.} Victime inconsciente (ne répond pas) mais qui \textbf{respire} (thorax mobile, régulier) : perte de connaissance avec respiration conservée. Les gargouillis signalent un risque d'inhalation de liquide (vomissements, salive) : la langue, relâchée, peut en outre obstruer les voies aériennes chez une victime allongée sur le dos.

\emph{Étapes de la PLS.}
\begin{enumerate}
\item S'agenouiller à côté de la victime ; libérer les voies aériennes par une bascule prudente de la tête en arrière.
\item Placer le bras le plus proche de soi à angle droit du corps, paume vers le haut.
\item Ramener le bras opposé contre la joue proche (dos de la main contre la joue) et le maintenir.
\item Saisir le genou opposé, le ramener en gardant le pied au sol.
\item Tout en maintenant la main contre la joue, faire rouler la victime vers soi d'un mouvement régulier.
\item Ajuster la jambe supérieure pour que la hanche et le genou soient à angle droit ; basculer la tête en arrière pour garder les voies aériennes libres.
\item Surveiller la respiration en permanence jusqu'à l'arrivée des secours ; alerter le 119 si ce n'est pas fait.
\end{enumerate}
\emph{Intérêt physiologique :} en décubitus latéral, la langue ne tombe plus en arrière et les liquides (vomissements, sang, salive) s'écoulent vers l'extérieur au lieu d'être inhalés dans les poumons : on prévient l'asphyxie.

\textbf{6. Débat.} Avec des délais de secours souvent supérieurs à 30 minutes (embouteillages de Yaoundé et Douala, distances en zone rurale), la courbe du document 1 montre que sans témoin formé, la survie est quasi nulle : à 10 minutes, il reste moins de 10 \% de chances. Former chaque citoyen, c'est transformer la population en premier maillon de la chaîne de survie. Actions proposées : créer un club de secourisme au lycée avec formation annuelle sur mannequin ; organiser des séances de sensibilisation au marché et dans les églises/mosquées du quartier en partenariat avec le Centre de Santé Intégré.
\end{exoresolu}

\begin{attention}[Pièges fréquents dans ce type d'exercice]
\begin{itemize}
\item Confondre perte de connaissance et arrêt cardiaque : la PLS est \textbf{interdite} si la victime ne respire pas — il faut masser.
\item Oublier d'alerter \emph{avant} ou \emph{pendant} les gestes : l'alerte fait partie de la chaîne de survie.
\item Donner des chiffres vagues : le correcteur attend 100--120 compressions/min, 5--6 cm, 30/2, 5 claques/5 compressions.
\item Pratiquer la manœuvre de Heimlich abdominale sur une femme enceinte ou un nourrisson : c'est une faute grave.
\end{itemize}
\end{attention}

\courssub{Bilan de l'activité}

Cette activité t'a permis de mobiliser l'ensemble des savoir-faire de la leçon dans des situations authentiques et différenciées. Elle a mis en évidence trois idées essentielles. D'abord, le \textbf{temps est le paramètre décisif} : la courbe de survie montre quantitativement que chaque minute compte et que le témoin, plus que l'hôpital, fait la différence. Ensuite, \textbf{le bon geste dépend du bon diagnostic} : masser une victime qui ne respire pas, désobstruer en adaptant la technique à la victime (nourrisson, adulte, femme enceinte), mettre en PLS une victime inconsciente qui respire — trois conduites distinctes qu'aucune improvisation ne peut remplacer. Enfin, le secourisme est un \textbf{acte citoyen} : au Cameroun, où les secours organisés peuvent tarder, la formation de tous aux gestes qui sauvent est un enjeu de santé publique. En maîtrisant ces gestes et en faisant preuve d'empathie envers la victime, tu deviens un maillon actif de la chaîne de survie.

\newpage

\coursec{Méthodes et exercices résolus}

\coursec{Les accidents cardiovasculaires : ampleur et enjeu de la chaîne de survie}

\begin{savaistu}[Épidémiologie et contexte camerounais]
Au Cameroun, les maladies cardiovasculaires (hypertension artérielle, insuffisance cardiaque, accident vasculaire cérébral) sont en forte augmentation du fait de la transition nutritionnelle, du vieillissement de la population et de l'urbanisation. L'\cle{arrêt cardiaque} extra-hospitalier touche plusieurs milliers de personnes chaque année ; la survie à la sortie d'hôpital reste faible (\textless 5 \%) faute de témoins formés et de défibrillateurs accessibles. Les numéros d'urgence à retenir sont le \textbf{119} (SAMU) et le \textbf{118} (Corps des Sapeurs-Pompiers). La \cle{chaîne de survie} — reconnaissance précoce et alerte, RCP précoce, défibrillation précoce, soins médicalisés — est le concept central : chaque minute gagnée sur les trois premiers maillons augmente significativement la probabilité de survie sans séquelle neurologique.
\end{savaistu}

\begin{popfigure}
\begin{tikzpicture}[scale=0.9, every node/.style={font=\small}]
% Nodes
\node[draw, rounded corners, fill=popblueL, thick] (1) at (0,0) {1. Reconnaissance précoce\& Alerte (119/118)};
\node[draw, rounded corners, fill=popgreenL, thick, align=center] (2) at (0,-1.5){2. RCP précoce par témoin\\(30 compressions / 2 insufflations)};
\node[draw, rounded corners, fill=poporangeL, thick, align=center] (3) at (0,-3){3. Défibrillation précoce\\(DAE public ou secours)};
\node[draw, rounded corners, fill=poppurpleL, thick, align=center] (4) at (0,-4.5){4. Soins médicalisés avancés\\(Réanimation, coronarographie, USIC)};
% Arrows
\draw[->, thick, popdark] (1.south) -- (2.north) node[midway, right] {Gain de temps : \SIrange{3}{5}{min} avant lésions cérébrales};
\draw[->, thick, popdark] (2.south) -- (3.north) node[midway, right] {Maintient fibrillation « fine » récupérable};
\draw[->, thick, popdark] (3.south) -- (4.north) node[midway, right] {Traitement étiologique + stabilisation};
% Survival curve sketch on right
\begin{scope}[xshift=7cm, yshift=-2.25cm, scale=0.6]
\draw[->] (0,0) -- (4,0) node[right] {Temps (min)};
\draw[->] (0,0) -- (0,3.5) node[above] {Survie (\%)\%};
\draw[popcurve, thick, domain=0:3.5, samples=50] plot (\x, {3.5*exp(-\x/1.5)});
\node[popdark, anchor=north west] at (0,3.5) {Sans RCP};
\draw[popcurve, thick, popgreen, domain=0:3.5, samples=50] plot (\x, {3.5*exp(-\x/3)});
\node[popgreen, anchor=south west] at (0,3.5) {Avec RCP immédiate};
\draw[dashed, thin] (1,0) -- (1,{3.5*exp(-1/1.5)});
\draw[dashed, thin] (1,0) -- (1,{3.5*exp(-1/3)});
\node[below] at (1,0) {1 min};
\node[below] at (3,0) {3 min};
\node[below] at (5,0) {5 min};
\end{scope}
\end{tikzpicture}
\legende{Schéma de la chaîne de survie (quatre maillons) et illustration de l'impact de la RCP précoce sur la courbe de survie. Chaque minute de retard sans RCP réduit la survie d'environ 10 \%.}
\end{popfigure}

\courssub{L'arrêt cardiaque : bases physiologiques et reconnaissance}

\begin{definition}[Arrêt cardiaque]
C'est la cessation brutale et inattendue de l'activité mécanique efficace du cœur. Le débit sanguin devient nul (\SI{0}{L.min^{-1}}), entraînant l'arrêt immédiat de la perfusion cérébrale. La conscience se perd en \SIrange{10}{15}{s}, la respiration s'arrête en \SIrange{30}{60}{s} (après quelques \cle{gasp} ou soupirs agoniques). Sans oxygénation, les lésions cérébrales irréversibles débutent à \SIrange{3}{5}{min} ; la mort biologique survient vers \SI{10}{min}.
\end{definition}

\begin{propriete}[Critères diagnostiques du secourisme]
Devant une victime inanimée, le diagnostic d'arrêt cardiaque repose sur l'association de deux signes \textbf{simultanés} :
\begin{enumerate}
\item \textbf{Inconscience} : absence de réponse à la stimulation verbale (interpellation forte) et tactile (secousses des épaules).
\item \textbf{Absence de respiration normale} : pas de mouvement thoracique, pas de souffle perçu, ou présence exclusive de \emph{gasp} (mouvements respiratoires saccadés, bruyants, inefficaces) lors de l'observation de \SI{10}{s} maximum.
\end{enumerate}
\emph{La recherche du pouls carotidien est exclue du protocole laïc : elle est peu fiable, source d'erreur et de perte de temps.}
\end{propriete}

\courssub{Fiche méthode 1 : Reconnaître un arrêt cardiaque et déclencher la chaîne de survie}

\begin{methode}[Conduite à tenir face à une victime inanimée : Protéger -- Alerter -- Secourir]
\begin{enumerate}
\item \textbf{Protéger} : s'assurer qu'il n'y a pas de danger pour le secouriste, la victime ou les tiers (circulation, incendie, électricité). Ne déplacer la victime que si danger imminent.
\item \textbf{Évaluer la conscience} : parler fort (« Entendez-vous ? »), taper dans les mains, secouer doucement les épaules. \textbf{Pas de réponse} $\rightarrow$ \cle{inconsciente}.
\item \textbf{Libérer les voies aériennes \& évaluer la respiration} : basculer la tête en arrière (menton relevé, front appuyé), ouvrir la bouche. Observer le thorax pendant \textbf{au plus \SI{10}{s}} : \textbf{Regarder} (mouvements), \textbf{Écouter} (bruits), \textbf{Sentir} (flux d'air sur la joue).
\item \textbf{Conclure} :
\begin{itemize}
\item Inconsciente + \textbf{Ne respire pas} (ou \emph{gasp} seulement) $\rightarrow$ \textbf{ARRÊT CARDIAQUE}.
\item Inconsciente + \textbf{Respire normalement} $\rightarrow$ \textbf{PERTE DE CONNAISSANCE} (voir Fiche 4).
\end{itemize}
\item \textbf{Alerter immédiatement} : composer le \textbf{119} (SAMU) ou \textbf{118} (Pompiers). Préciser : lieu exact (ex. « Marché central, Yaoundé, près de l'entrée principale »), nombre de victimes, état (inconsciente, ne respire pas), gestes en cours. Mettre le téléphone en \textbf{haut-parleur} pour garder les mains libres.
\item \textbf{Secourir sans délai} : débuter la \textbf{RCP} (massage cardiaque + insufflations) et envoyer un témoin chercher un \textbf{DAE} (Défibrillateur Automatisé Externe) si disponible à proximité.
\end{enumerate}
\textbf{Réflexe clé} : \emph{Inconscience + Absence de respiration normale = Arrêt cardiaque = RCP immédiate}. Ne jamais attendre, ne jamais chercher le pouls.
\end{methode}

\begin{exoresolu}[Effondrement au marché central de Yaoundé]
Au marché central de Yaoundé, un homme de 55 ans s'effondre brutalement devant vous. Il ne répond pas quand vous l'interpellez et le secouez par les épaules. Penché au-dessus de son visage, vous n'observez aucun mouvement du thorax et ne sentez aucun souffle pendant \SI{10}{s}. Un témoin affirme : « Il a peut-être juste glissé, attendons qu'il se réveille. »
\begin{enumerate}
\item Quel est votre diagnostic de secourisme ? Justifiez-le par deux signes précis.
\item Expliquez pourquoi l'attitude du témoin est dangereuse, en vous appuyant sur une donnée physiologique.
\item Citez, dans l'ordre, les trois premières actions que vous entreprenez.
\end{enumerate}
\tcblower
\textbf{Solution détaillée}
\begin{enumerate}
\item \textbf{Diagnostic : arrêt cardiaque.} Deux signes le confirment : (a) la victime est \textbf{inconsciente} (absence de réponse à la voix et aux secousses des épaules) ; (b) elle \textbf{ne respire pas} : aucun mouvement thoracique, aucun souffle perçu pendant les \SI{10}{s} d'observation. L'association inconscience + arrêt respiratoire signe l'arrêt cardiaque : le cœur ne propulse plus de sang, l'apport d'\ce{O2} cérébral est nul.
\item \textbf{Danger de l'attente.} En arrêt cardiaque, le débit sanguin est nul. Le cerveau, grand consommateur d'\ce{O2} (\SI{20}{\%} de la consommation totale pour \SI{2}{\%} de la masse), ne stocke pas l'oxygène. Des \textbf{lésions neuronales irréversibles apparaissent dès \SIrange{3}{5}{min}} d'anoxie totale. De plus, la probabilité de succès de la défibrillation (cause majeure : fibrillation ventriculaire) diminue d'environ \textbf{\SI{10}{\%} par minute} de retard. Attendre « qu'il se réveille » condamne la victime : sans RCP ni DAE, la survie est quasi nulle après \SI{10}{min}.
\item \textbf{Les trois premières actions, dans l'ordre :}
\begin{enumerate}
\item \textbf{Alerter} les secours (119 ou 118) ou désigner un témoin précis (« Vous, en chemise bleue, appelez le 119 ! ») en indiquant le lieu exact ;
\item \textbf{Débuter immédiatement le massage cardiaque externe} : 30 compressions thoraciques au centre de la poitrine (sternum, moitié inférieure), fréquence \SIrange{100}{120}{min^{-1}}, profondeur \SI{5}{cm} ;
\item \textbf{Poursuivre la RCP} (cycles 30/2) et faire rechercher un DAE, jusqu'à l'arrivée des secours, la reprise d'une respiration normale ou l'épuisement du secouriste.
\end{enumerate}
\end{enumerate}
\textbf{Conclusion} : devant une victime inconsciente qui ne respire pas, le diagnostic d'arrêt cardiaque est posé en \textless \SI{10}{s} ; seule une action immédiate (alerte + RCP) peut sauver la vie.
\end{exoresolu}

\courssub{Les gestes qui sauvent en cas d'arrêt cardiaque : la RCP}

\begin{propriete}[Principes hémodynamiques de la RCP]
Le massage cardiaque externe (MCE) comprime le cœur entre le sternum et la colonne vertébrale (mécanisme de la \cle{pompe thoracique}) : la compression systolique éjecte le sang vers l'aorte et les artères pulmonaires ; la relâche (diastole) permet le remplissage cardiaque par retour veineux. Une RCP de qualité restitue \textbf{25 à 30 \%} du débit cardiaque normal (\SI{5}{L.min^{-1}} au repos), soit \SIrange{1,25}{1,5}{L.min^{-1}}. Ce flux résiduel : (1) oxygène le cerveau (retardant l'anoxie), (2) perfuse le myocarde via les coronaires (maintenant la fibrillation ventriculaire « fine » sensible au choc électrique).
\end{propriete}

\begin{popfigure}
\begin{tikzpicture}[scale=0.7, every node/.style={font=\small}]
% Thorax cross-section
\draw[thick, fill=poppinkL] (0,0) ellipse (3.5 and 2.5); % Rib cage
\draw[thick, fill=popred!20] (0,-0.2) ellipse (1.8 and 1.2); % Heart
\draw[thick, fill=popblue!30] (-2.5,0) -- (-2.5,1.5) -- (2.5,1.5) -- (2.5,0) -- cycle; % Sternum top
\draw[<->, thick, popdark] (0,1.5) -- (0,2.5) node[midway, right] {Profondeur \SI{5}{cm}};
\draw[->, thick, poporange, dashed] (0,2.5) -- (0,1.5) node[midway, left] {Compression};
\draw[->, thick, popgreen, dashed] (0,1.5) -- (0,2.5) node[midway, right] {Relâche complète};
% Hands position
\draw[fill=popcream, draw=popdark, thick] (-0.8,2.8) ellipse (0.6 and 0.3); % Hand 1 heel
\draw[fill=popcream, draw=popdark, thick] (0.8,2.8) ellipse (0.6 and 0.3); % Hand 2 heel
\draw[thick] (-0.8,2.5) -- (-0.8,3.5); % Arm 1
\draw[thick] (0.8,2.5) -- (0.8,3.5); % Arm 2
\node[popdark, align=center] at (0,3.8) {Épaules verticales\\au-dessus des mains\\Bras tendus};
\node[popdark, align=center] at (0,-2.5) {Centre de la poitrine\\(Sternum, moitié inférieure)};
\end{tikzpicture}
\legende{Coupe schématique du thorax lors du massage cardiaque : position des mains (talon au centre du sternum), direction de la compression (verticale, \SI{5}{cm}), importance de la relâche complète pour le remplissage veineux.}
\end{popfigure}

\courssub{Fiche méthode 2 : Pratiquer la réanimation cardio-pulmonaire (RCP)}

\begin{methode}[La RCP de l'adulte en 6 points techniques]
\begin{enumerate}
\item \textbf{Installer} la victime à plat dos sur une surface \textbf{dure} (sol, planche). Jamais sur un lit ou matelas mou (perte d'efficacité).
\item \textbf{Positionner les mains} : talon d'une main au \textbf{centre de la poitrine} (sternum, moitié inférieure, entre les mamelons), l'autre main par-dessus, doigts entrelacés (ne pas appuyer sur les côtes ni sur le processus xiphoïde). Bras tendus, épaules à la verticale des mains (utiliser le poids du corps).
\item \textbf{Comprimer} : enfoncer le sternum de \textbf{\SI{5}{cm} (entre 5 et 6 cm)} chez l'adulte. Fréquence : \textbf{100 à 120 compressions par minute}. \textbf{Laisser le thorax reprendre son expansion complète} entre deux compressions (pas d'appui résiduel).
\item \textbf{Ventiler (si formé et équipé)} : après 30 compressions, \textbf{2 insufflations} bouche-à-bouche (ou bouche-à-nez/masque). Bascule de la tête (menton relevé), pincement du nez, inspiration normale, insufflation \SI{1}{s} jusqu'à soulèvement thoracique visible. \textbf{Si non formé ou réticent : compressions seules continues (\"Hands-only CPR\")}, sans interruption.
\item \textbf{Relayer} toutes les \textbf{\SI{2}{min}} (ou \SI{5}{cycles} 30/2) si un second secouriste est présent. Le changement doit durer \textless \SI{5}{s} pour ne pas interrompre le flux coronarien.
\item \textbf{Utiliser le DAE} dès son arrivée : allumer, coller électrodes (schéma sur boîtier : sous clavicule droite et apex gauche), \textbf{ne toucher personne} pendant l'analyse et la délivrance du choc. Reprendre RCP immédiatement après le choc (ou si « pas de choc conseillé »).
\end{enumerate}
\textbf{Formule à retenir} : \cle{RCP = 30 compressions / 2 insufflations} ; \cle{100--120 compressions/min} ; \cle{profondeur \approx \SI{5}{cm}} ; \cle{relais toutes les 2 min}.
\end{methode}

\begin{exoresolu}[Efficacité du massage cardiaque : un calcul qui sauve]
Lors d'un match de football à Douala, un joueur de 24 ans s'effondre, inconscient et sans respiration. Deux secouristes formés se relaient pour le masser. Le cœur au repos d'un adulte débite environ \SI{5}{L.min^{-1}} de sang ; un massage cardiaque externe bien exécuté ne restitue que 25 à \SI{30}{\%} de ce débit.
\begin{enumerate}
\item Calculez le débit sanguin minimal et maximal assuré par la RCP chez ce joueur.
\item Les secouristes massent à 110 compressions/min pendant \SI{4}{min} avant l'arrivée des secours. Combien de compressions ont-ils réalisées ?
\item Expliquez pourquoi, malgré un débit inférieur au débit normal, la RCP reste vitale pour le cerveau et le cœur.
\end{enumerate}
\tcblower
\textbf{Solution détaillée}
\begin{enumerate}
\item Le débit normal est $Q_0 = \SI{5}{L.min^{-1}}$. La RCP en restitue 25 à \SI{30}{\%} :
\begin{align*}
Q_{\min} &= 5 \times \dfrac{25}{100} = \SI{1,25}{L.min^{-1}}\\
Q_{\max} &= 5 \times \dfrac{30}{100} = \SI{1,5}{L.min^{-1}}
\end{align*}
La RCP assure donc un débit compris entre \SI{1,25}{L.min^{-1}} et \SI{1,5}{L.min^{-1}}.
\item Nombre de compressions : $N = 110 \times 4 = 440$ compressions (en supposant un massage continu ; en pratique, les pauses d'insufflation réduisent légèrement ce total à environ 400--420).
\item \textbf{Justification physiologique.} Même réduit au quart du débit normal, ce flux sanguin : (a) apporte encore de l'\ce{O2} aux \textbf{cellules cérébrales}, retardant l'apparition des lésions irréversibles qui surviennent en \SIrange{3}{5}{min} d'anoxie totale ; (b) perfuse le \textbf{myocarde} lui-même via les artères coronaires (perfusion diastolique), maintenant le cœur en état de « fibrillation fine » récupérable par le choc du DAE. Sans RCP, le débit est nul : le cerveau meurt et le myocarde devient réfractaire à la défibrillation (asystolie ou fibrillation grossière). La RCP « achète du temps » en attendant le DAE et les secours médicalisés.
\end{enumerate}
\textbf{Conclusion} : avec un débit de seulement \SIrange{1,25}{1,5}{L.min^{-1}}, la RCP ne remplace pas le cœur, mais elle maintient la vie cérébrale et la récupérabilité cardiaque : c'est le maillon central de la chaîne de survie.
\end{exoresolu}

\courssub{L'étouffement total : mécanisme et gestes selon la victime}

\begin{definition}[Étouffement total (obstruction complète des voies aériennes)]
Blocage brutal du flux d'air au niveau du larynx ou de la trachée par un corps étranger (aliment, objet). La victime \textbf{ne peut ni parler, ni tousser, ni respirer, ni émettre le moindre son}. Elle porte instinctivement les mains à la gorge (\cle{signe universel de l'étouffement}), devient anxieuse, puis cyanosée (bleutée lèvres/ongles), et perd connaissance en \SIrange{1}{2}{min} si non secourue.
\end{definition}

\begin{propriete}[Principe commun des manœuvres d'expulsion]
Toutes les manœuvres visent à créer une \textbf{augmentation brutale de la pression intrathoracique} (toux artificielle) pour expulser le corps étranger :
\begin{itemize}
\item \textbf{Claques dorsales} : vibration mécanique + onde de pression rétrograde.
\item \textbf{Compressions abdominales (Heimlich)} : élévation brutale du diaphragme $\rightarrow$ compression des poumons $\rightarrow$ flux expiratoire violent.
\item \textbf{Compressions thoraciques} : compression directe de la cage thoracique $\rightarrow$ même effet, utilisée quand l'abdomen est contre-indiqué (nourrisson, femme enceinte, obésité massive).
\end{itemize}
\end{propriete}

\begin{popfigure}
\begin{tikzpicture}[scale=0.8, every node/.style={font=\small}]
% Adult Heimlich
\begin{scope}[xshift=-5cm]
\draw[thick, fill=popblueL] (0,0) rectangle (3,4);
\draw[->, thick, poporange] (1.5,1) -- (1.5,0.2) node[below] {Poing : dessus nombril, sous sternum};
\draw[->, thick, poporange] (1.5,1) -- (1.5,1.8) node[above] {Tractions : vers soi \& vers le haut};
\node[align=center, popdark] at (1.5, -0.8) {Adulte / Enfant > 1 an\\(Manœuvre de Heimlich)};
\end{scope}
% Infant Back Blows
\begin{scope}[xshift=1cm]
\draw[thick, fill=popgreenL] (0,0) rectangle (3,4);
\draw[->, thick, poporange] (1.5,3.5) -- (1.5,4.3) node[above] {5 claques entre omoplates};
\draw[thick, fill=popcream] (1.5,2) ellipse (0.8 and 0.5); % Infant on forearm
\node[align=center, popdark] at (1.5, -0.8) {Nourrisson < 1 an\\(Tête déclive, pas abdominales)};
\end{scope}
% Pregnant Chest Thrusts
\begin{scope}[xshift=7cm]
\draw[thick, fill=poppurpleL] (0,0) rectangle (3,4);
\draw[->, thick, poporange] (1.5,2.5) -- (1.5,1.7) node[below] {Poing sur sternum};
\draw[->, thick, poporange] (1.5,2.5) -- (1.5,3.3) node[above] {Tractions vers soi};
\node[align=center, popdark] at (1.5, -0.8) {Femme enceinte / Obèse\\(Compressions thoraciques)};
\end{scope}
\end{tikzpicture}
\legende{Adaptation des manœuvres d'expulsion selon la victime : Heimlich (abdominales) pour l'adulte ; claques dorsales + compressions thoraciques (2 doigts) pour le nourrisson ; compressions thoraciques (poing sur sternum) pour la femme enceinte.}
\end{popfigure}

\courssub{Fiche méthode 3 : Secourir une victime d'étouffement total}

\begin{methode}[Étouffement total : reconnaître puis agir selon la victime]
\textbf{Reconnaissance formelle} : victime consciente mais \textbf{incapable de parler, tousser, respirer, émettre un son} + signe mains à la gorge + cyanose. \emph{Si la victime tousse efficacement : l'encourager à tousser, ne pas intervenir.}
\begin{enumerate}
\item \textbf{Adulte ou enfant de plus de 1 an} :
\begin{enumerate}
\item Donner jusqu'à \textbf{5 claques dans le dos} : victime penchée en avant (gravité), paume de la main entre les deux omoplates.
\item Si échec : \textbf{5 compressions abdominales} (Heimlich) : poing fermé pouce vers l'intérieur, placé \textbf{au-dessus du nombril, bien sous le processus xiphoïde} (sternum) ; l'autre main saisit le poing ; tractions vigoureuses \textbf{vers soi et vers le haut}.
\item Alterner 5 claques / 5 compressions jusqu'à expulsion du corps étranger ou perte de connaissance.
\end{enumerate}
\item \textbf{Nourrisson (moins de 1 an)} :
\begin{enumerate}
\item Le placer \textbf{à plat ventre sur l'avant-bras}, tête plus basse que le tronc (tête déclive), soutenir la mâchoire (pouce/index) sans comprimer le cou.
\item Donner \textbf{5 claques fermes entre les omoplates} (paume de la main).
\item Si échec : le retourner sur le dos (toujours tête basse), effectuer \textbf{5 compressions thoraciques} avec \textbf{deux doigts} (index/majeur) sur le sternum (ligne intermamillaire), profondeur \SI{4}{cm} environ.
\item Alterner 5 claques / 5 compressions thoraciques.
\end{enumerate}
\textbf{Interdiction formelle} : \textbf{jamais de compressions abdominales} chez le nourrisson (foie volumineux dépassant le rebord costal $\rightarrow$ risque rupture hépatique hémorragique).
\item \textbf{Femme enceinte (ou personne obèse)} :
\begin{enumerate}
\item 5 claques dans le dos (victime penchée en avant).
\item Si échec : \textbf{compressions thoraciques} (poing sur le sternum, moitié inférieure, tractions vers soi) \textbf{à la place des compressions abdominales}.
\end{enumerate}
\textbf{Justification} : l'utérus gravide occupe l'abdomen (risque fœtal, utérin, inefficacité diaphragme haut).
\item \textbf{Si la victime perd connaissance} : l'allonger au sol, \textbf{alerter les secours (119/118)}, commencer la \textbf{RCP} (les compressions thoraciques de la RCP peuvent expulser le corps étranger). \textbf{Ne pas faire de doigté à l'aveugle dans la bouche.}
\end{enumerate}
\end{methode}

\begin{exoresolu}[Trois situations d'étouffement, trois conduites différentes]
Lors d'une fête familiale à Bafoussam, trois incidents surviennent : (A) un oncle de 60 ans s'étouffe avec un morceau de viande, ne parle plus et porte les mains à sa gorge ; (B) un nourrisson de 8 mois s'étouffe avec un petit jouet et n'émet plus aucun son ; (C) une tante enceinte de 7 mois s'étouffe avec une cacahuète et ne respire plus.
\begin{enumerate}
\item Pour chaque victime, précisez la manœuvre de premier secours adaptée et justifiez la différence de conduite.
\item Expliquez pourquoi on ne pratique jamais de compressions abdominales chez le nourrisson ni chez la femme enceinte.
\end{enumerate}
\tcblower
\textbf{Solution détaillée}
\begin{enumerate}
\item \textbf{Victime A (adulte)} : étouffement total confirmé (silence, mains à la gorge, cyanose). Conduite : jusqu'à 5 claques dans le dos (victime penchée en avant) ; si échec, 5 compressions abdominales (manœuvre de Heimlich : poing au-dessus du nombril, sous le sternum, tractions vers soi et vers le haut). Alterner jusqu'à expulsion. La claque crée une vibration/décollement ; la compression abdominale élève le diaphragme $\rightarrow$ toux artificielle.
\item \textbf{Victime B (nourrisson)} : placer à plat ventre sur l'avant-bras, tête déclive (gravité), 5 claques entre les omoplates ; si échec, retourner, 5 compressions thoraciques à deux doigts sur le sternum. Alterner.
\item \textbf{Victime C (femme enceinte)} : 5 claques dans le dos puis, en cas d'échec, \textbf{compressions thoraciques} (poing sur le sternum, tractions vers soi) \textbf{au lieu des compressions abdominales}.
\item \textbf{Justifications anatomiques.}
\begin{itemize}
\item \textbf{Nourrisson} : le foie est volumineux et dépasse le rebord costal ; une poussée abdominale brutale risque une \textbf{rupture hépatique hémorragique} mortelle. La position tête déclive sur l'avant-bras exploite la gravité pour l'expulsion.
\item \textbf{Femme enceinte} : l'utérus gravide remplit l'abdomen. La manœuvre de Heimlich comprimerait l'utérus (risque \textbf{placenta praevia, détresse fœtale, rupture utérine}) et serait inefficace car le diaphragme est déjà refoulé vers le haut par l'utérus. Les compressions thoraciques produisent le même effet expulsif (pression intrathoracique) sans ces risques.
\end{itemize}
\end{enumerate}
\textbf{Conclusion} : le principe (claques dorsales puis pressions expulsives alternées) est constant, mais la \emph{cible anatomique} des pressions s'adapte à la victime : abdomen (adulte), thorax (nourrisson, femme enceinte).
\end{exoresolu}

\courssub{La perte de connaissance : reconnaissance et position latérale de sécurité}

\begin{definition}[Perte de connaissance (Syncope, Malaise vagal, etc.)]
État transitoire d'inconscience (perte de la vigilance, absence de réponse aux stimulations) \textbf{avec conservation de la respiration normale} et de la circulation (pouls présent). La victime « respire » (thorax mobile, rythme régulier, souffle perçu). Ce n'est \textbf{pas} un arrêt cardiaque.
\end{definition}

\begin{propriete}[Risques vitaux de l'inconscient couché sur le dos]
Deux mécanismes menacent l'inconscient en décubitus dorsal :
\begin{enumerate}
\item \textbf{Obstruction des voies aériennes par la langue} : le relâchement musculaire (hypotonie) fait retomber la base de la langue en arrière, obstruant le pharynx (\"langue en arrière\").
\item \textbf{Inhalation (fausse route) de liquides} : salive, sang, vomissements s'écoulent passivement vers le pharynx et la trachée (pas de réflexe de toux ni de déglutition) $\rightarrow$ asphyxie ou pneumonie d'inhalation.
\end{enumerate}
La \cle{Position Latérale de Sécurité (PLS)} traite ces deux risques simultanément.
\end{propriete}

\begin{popfigure}
\begin{tikzpicture}[scale=0.7, every node/.style={font=\small}]
% Step 1: Prep
\begin{scope}[yshift=0]
\node[align=center, popdark] at (0,2) {Étape 1 : Préparation};
\draw[thick, fill=popblueL] (-2,0) rectangle (2,1.5);
\draw[thick] (-0.5,0.75) -- (-0.5,1.5); % Arm at right angle
\draw[thick] (0.5,0) -- (1.5,1); % Opposite leg bent
\draw[thick] (-1.5,0.75) -- (-2.5,1.5); % Opposite hand to cheek
\node[popdark, align=center] at (0,-0.5) {Bras proche angle droit\\Jambe opposée fléchie\\Main opposée sur joue};
\end{scope}
% Step 2: Roll
\begin{scope}[yshift=-3.5]
\node[align=center, popdark] at (0,2) {Étape 2 : Basculement};
\draw[thick, fill=popgreenL] (-2,0) rectangle (2,1.5);
\draw[thick, ->, poporange] (0,0.75) -- (2,0.75) node[right] {Tirer sur jambe fléchie};
\draw[thick] (-1,0.5) arc (180:0:1 and 0.5); % Body rolling
\node[popdark, align=center] at (0,-0.5) {Rouler vers soi\\en gardant tête alignée};
\end{scope}
% Step 3: Final PLS
\begin{scope}[yshift=-7]
\node[align=center, popdark] at (0,2) {Étape 3 : PLS Stabilisée};
\draw[thick, fill=poporangeL] (-2,0) rectangle (2,1.5);
\draw[thick] (0.5,0) -- (1.5,1); % Upper leg bent 90
\draw[thick] (-1.5,1) -- (-0.5,1.5); % Head tilted back
\draw[->, thick, popblue] (-1,1.2) -- (-1,0.5) node[below] {Bouche vers le sol};
\node[popdark, align=center] at (0,-0.5) {Jambe sup. fléchie 90° (stabilité)\\\textbf{Tête en hyper-extension, bouche basse}};
\end{scope}
\end{tikzpicture}
\legende{Réalisation de la Position Latérale de Sécurité (PLS) en trois temps clés. La stabilité repose sur la jambe supérieure fléchie à 90° ; l'hyper-extension de la tête et la bouche vers le bas assurent la pérméabilité des voies aériennes et l'écoulement des liquides.}
\end{popfigure}

\courssub{Fiche méthode 4 : Reconnaître une perte de connaissance et mettre en PLS}

\begin{methode}[Perte de connaissance : diagnostic et Position Latérale de Sécurité]
\begin{enumerate}
\item \textbf{Confirmer l'inconscience} : interpellation + secousses épaules $\rightarrow$ pas de réponse.
\item \textbf{Vérifier la respiration} : bascule tête en arrière (libération voies aériennes), observer thorax \SI{10}{s} max. \textbf{Respiration normale présente} (mouvements réguliers, souffle senti).
\item \textbf{Diagnostic} : \textbf{Perte de connaissance} (Inconscient + Respire). \textbf{Pas de RCP.}
\item \textbf{Alerter} les secours (119 / 118) : lieu, état (inconsciente, respire), contexte (diabétique, traumatisme, etc.).
\item \textbf{Mettre en PLS} (voir schéma) :
\begin{enumerate}
\item S'agenouiller à côté ; bras le plus proche en angle droit (paume vers le haut).
\item Ramener la jambe opposée fléchie (pied au sol) ; main opposée contre la joue (paume contre joue).
\item Tirer sur la jambe fléchie pour rouler la victime vers soi (tête alignée sur le tronc).
\item Ajuster : jambe supérieure fléchie à \textbf{90°} (hanche et genou) pour la stabilité ; \textbf{basculer la tête en arrière} (menton relevé) ; \textbf{bouche ouverte vers le bas}.
\end{enumerate}
\item \textbf{Surveiller en continu} la respiration (regarder, écouter, sentir) jusqu'à l'arrivée des secours. \textbf{Si la respiration s'arrête} : retourner sur le dos $\rightarrow$ RCP immédiate.
\end{enumerate}
\textbf{But de la PLS} : langue libérée (tombe en avant) + écoulement des liquides (gravité) = voies aériennes pérennes.
\end{methode}

\begin{exoresolu}[Malaise en classe de Terminale]
Pendant un cours, votre camarade Aïcha, diabétique connue, s'affaisse sur sa table. Elle ne répond pas quand vous l'appelez, mais son thorax se soulève régulièrement et vous sentez son souffle. Un élève propose de lui faire boire du jus sucré « pour la réveiller » ; un autre veut commencer le massage cardiaque.
\begin{enumerate}
\item Posez le diagnostic de secourisme et justifiez-le.
\item Expliquez pourquoi les deux propositions des élèves sont dangereuses.
\item Décrivez la conduite à tenir correcte, dans l'ordre.
\end{enumerate}
\tcblower
\textbf{Solution détaillée}
\begin{enumerate}
\item \textbf{Diagnostic : perte de connaissance (inconscience avec respiration conservée).} Aïcha ne répond pas aux stimulations (inconsciente) mais respire normalement (thorax mobile, souffle perceptible, rythme régulier) : le cœur fonctionne, il n'y a pas arrêt cardiaque.
\item \textbf{Analyse des deux erreurs graves :}
\begin{itemize}
\item \emph{Faire boire du jus} : une victime inconsciente a perdu ses réflexes de protection des voies aériennes (déglutition, toux). Le liquide passe dans la trachée (\emph{fausse route}) $\rightarrow$ \textbf{asphyxie aiguë} ou \textbf{pneumonie d'inhalation} sévère. \textbf{On ne donne JAMAIS à boire ni à manger à un inconscient.}
\item \emph{Massage cardiaque} : la RCP est indiquée \textbf{uniquement} si la victime ne respire pas. Masser un cœur qui bat efficacement est inutile, fatigue le secouriste, et risque de provoquer \textbf{fractures costales, contusions pulmonaires, lésions hépatiques/spléniques}.
\end{itemize}
\item \textbf{Conduite à tenir, dans l'ordre :}
\begin{enumerate}
\item \textbf{Libérer les voies aériennes} : basculer doucement la tête d'Aïcha en arrière (menton relevé).
\item \textbf{Alerter} les secours (119) et prévenir l'infirmier de l'établissement (contexte diabète $\rightarrow$ hypoglycémie probable).
\item \textbf{Mettre en PLS} : la rouler sur le côté, jambe supérieure fléchie à 90° pour la stabilité, tête basculée en arrière, bouche vers le bas.
\item \textbf{Surveiller la respiration} sans interruption jusqu'à l'arrivée des secours, prêt à débuter la RCP si la respiration s'arrête.
\end{enumerate}
\end{enumerate}
\textbf{Conclusion} : inconsciente mais respirante, la victime relève de la PLS et de la surveillance, jamais du massage cardiaque ni de l'alimentation orale.
\end{exoresolu}

\courssub{Synthèse : la conduite à tenir du citoyen-secouriste}

\begin{methode}[L'arbre de décision du citoyen-secouriste : Triage vital]
Devant toute victime inanimée (ou situation d'urgence), appliquer systématiquement l'algorithme à deux questions vitales :
\begin{enumerate}
\item \textbf{Question 1 : La victime répond-elle ?} (Interpeller fort, secouer épaules)
\begin{itemize}
\item \textbf{OUI} $\rightarrow$ \textbf{Consciente} : Interroger (maladie, traumatisme), rassurer, traiter le problème visible (étouffement, hémorragie, douleur), \textbf{Alerter} si nécessaire (119/118). Surveiller l'évolution.
\item \textbf{NON} $\rightarrow$ \textbf{Inconsciente} $\rightarrow$ Passer à la Question 2.
\end{itemize}
\item \textbf{Question 2 : Respire-t-elle normalement ?} (Bascule tête en arrière, observer thorax \textbf{\SI{10}{s} max})
\begin{itemize}
\item \textbf{OUI} (thorax mobile, souffle régulier) $\rightarrow$ \textbf{PERTE DE CONNAISSANCE} : \textbf{Alerter (119/118)} + \textbf{PLS} + \textbf{Surveillance continue} de la respiration.
\item \textbf{NON} (immobile, ou \emph{gasp} seulement) $\rightarrow$ \textbf{ARRÊT CARDIAQUE} : \textbf{Alerter (119/118)} + \textbf{RCP immédiate (30/2, 100-120/min, 5 cm)} + \textbf{DAE} dès dispo.
\end{itemize}
\end{enumerate}
\textbf{Réflexes transversaux (Savoir-être)} : \textbf{Protéger} avant tout (danger incendie, route, électricité) ; \textbf{Ne jamais déplacer} sauf danger imminent ; \textbf{Ne jamais donner à boire/manger} à un inconscient ; \textbf{Rester auprès} de la victime ; \textbf{Faire preuve d'empathie} (paroles apaisantes, contact physique rassurant dès reprise de conscience) ; \textbf{Transmettre} les infos aux secours (heure, gestes faits, antécédents connus).
\end{methode}

\begin{exoresolu}[Accident de la route sur l'axe Yaoundé--Bafoussam]
Témoin d'un accident, vous approchez trois victimes après avoir sécurisé la zone (triangle, gilet, coupure contact). Victime 1 : consciente, elle hurle et se plaint d'une douleur à la jambe. Victime 2 : ne répond pas, mais respire calmement. Victime 3 : ne répond pas et n'a aucun mouvement respiratoire.
\begin{enumerate}
\item Attribuez à chaque victime la conduite à tenir prioritaire, en justifiant par les signes observés.
\item Classez ces victimes par ordre de priorité d'intervention et justifiez ce classement.
\end{enumerate}
\tcblower
\textbf{Solution détaillée}
\begin{enumerate}
\item Attribution de la conduite à tenir par victime :
\begin{enumerate}
\item \textbf{Victime 1} : \textbf{Consciente} (parle, exprime douleur) $\rightarrow$ Pas de geste vital immédiat. Conduite : rassurer, immobiliser le membre douloureux (si suspicion fracture), couvrir, surveiller, transmettre infos aux secours.
\item \textbf{Victime 2} : \textbf{Inconsciente mais respirant normalement} $\rightarrow$ \textbf{Perte de connaissance}. Conduite : bascule tête en arrière $\rightarrow$ \textbf{Alerter (119)} $\rightarrow$ \textbf{PLS} $\rightarrow$ Surveillance continue respiration.
\item \textbf{Victime 3} : \textbf{Inconsciente et sans respiration} (ou gasps) $\rightarrow$ \textbf{Arrêt cardiaque}. Conduite : \textbf{Alerter (119)} $\rightarrow$ \textbf{RCP immédiate} (30 compressions / 2 insufflations, 100-120/min, 5 cm) $\rightarrow$ Recherche/Utilisation \textbf{DAE}.
\end{enumerate}
\item \textbf{Ordre de priorité d'intervention : Victime 3 $\rightarrow$ Victime 2 $\rightarrow$ Victime 1.}
\textbf{Justification physiologique et éthique (triage) :}
\begin{itemize}
\item \textbf{Victime 3 (Priorité 1 - Urgence vitale absolue)} : Arrêt cardiaque = débit nul. Lésions cérébrales irréversibles en \SIrange{3}{5}{min}. Survie chute de \SI{10}{\%}/\textit{min}. Elle \textbf{doit} être massée \textbf{immédiatement} (ou par un relais désigné si vous êtes seul, vous massez la V3 pendant qu'un témoin alerte et gère les autres).
\item \textbf{Victime 2 (Priorité 2 - Urgence vitale potentielle)} : Inconscience = risque obstruction voies aériennes (langue, vomissements) et risque évolution vers arrêt cardiaque. La PLS et la surveillance sont urgentes mais moins « immédiates » que la RCP sur un cœur arrêté.
\item \textbf{Victime 1 (Priorité 3 - Urgence relative)} : Consciente = voies aériennes pérennes, cerveau oxygéné, circulation efficace. La douleur (fracture ?) est invalidante mais n'engage pas le pronostic vital \textit{à la minute}.
\end{itemize}
\end{enumerate}
\textbf{Conclusion} : le triage repose sur deux questions (Répond-elle ? Respire-t-elle ?). L'absence de respiration chez un inconscient constitue \textbf{toujours} la priorité absolue.
\end{exoresolu}

\courssub{Fiche méthode 5 : Rédiger une réponse de secourisme au Bac}

\begin{methode}[Structurer sa copie comme un secouriste agit : les 5 règles d'or]
\begin{enumerate}
\item \textbf{Nommer le diagnostic} avec le vocabulaire \textbf{exact du programme} : « \cle{arrêt cardiaque} », « \cle{perte de connaissance} », « \cle{étouffement total} ». \textbf{Interdiction} : termes vagues (\"malaise\", \"évanouissement\", \"inconfort\") sans définition précise.
\item \textbf{Justifier par les signes observés} : citer \textbf{explicitement} les deux critères (conscience, respiration) tels qu'ils apparaissent dans l'énoncé (ex. : « La victime ne répond pas à la voix ni aux secousses $\rightarrow$ inconsciente ; son thorax ne se soulève pas $\rightarrow$ ne respire pas »).
\item \textbf{Donner les gestes dans l'ordre chronologique strict} : Alerter (119/118) \textbf{avant} ou \textbf{pendant} le secours ; PLS \textbf{avant} surveillance ; 5 claques \textbf{avant} 5 compressions ; 30 compressions \textbf{avant} 2 insufflations.
\item \textbf{Chiffrer les paramètres techniques} quand l'énoncé le demande (ou pour montrer la maîtrise) : \textbf{30/2}, \textbf{100--120 compressions/min}, \textbf{\SI{5}{cm} profondeur}, \textbf{5 claques / 5 compressions}, \textbf{\SI{10}{s} observation}, \textbf{relais 2 min}.
\item \textbf{Conclure par la finalité physiologique} : « La RCP maintient un débit résiduel (\SIrange{1,25}{1,5}{L.min^{-1}}) oxygénant le cerveau et perfusant les coronaires en attendant le DAE » ; « La PLS libère les voies aériennes (langue) et empêche la fausse route » ; « Les manœuvres d'étouffement créent une toux artificielle par surpression intrathoracique ».
\end{enumerate}
\end{methode}

\begin{exoresolu}[Question de synthèse type Bac]
Définissez la chaîne de survie et expliquez, en vous appuyant sur les données physiologiques de l'arrêt cardiaque, pourquoi chacun de ses maillons doit être enclenché le plus tôt possible.
\tcblower
\textbf{Solution détaillée}
\textbf{Définition.} La \cle{chaîne de survie} est l'enchaînement ordonné de quatre maillons indissociables :
\begin{enumerate}
\item \textbf{Reconnaissance précoce} de l'arrêt cardiaque et \textbf{Alerte} des secours (119/118) ;
\item \textbf{RCP précoce} par les témoins (massage + insufflations) ;
\item \textbf{Défibrillation précoce} par DAE (public ou secours) ;
\item \textbf{Soins médicalisés avancés} (réanimation avancée, coronarographie, soins intensifs en USIC).
\end{enumerate}

\textbf{Justification physiologique de la précocité de chaque maillon.}
\begin{itemize}
\item \textbf{Maillon 1 (Reconnaissance/Alerte) :} En arrêt cardiaque, le débit sanguin est \textbf{nul}. L'apport d'\ce{O2} cérébral cesse. Les neurones, incapables de stocker l'\ce{O2} et très consommateurs (\SI{20}{\%} \ce{O2} corporel), subissent des \textbf{lésions irréversibles dès \SIrange{3}{5}{min}}. L'alerte précoce conditionne l'arrivée du DAE (maillon 3) et de l'équipe médicalisée (maillon 4).
\item \textbf{Maillon 2 (RCP) :} La RCP ne restitue que \textbf{25 à 30 \%} du débit normal (\SIrange{1,25}{1,5}{L.min^{-1}} pour $Q_0=\SI{5}{L.min^{-1}}$). Elle ne réanime pas seule, mais elle \textbf{ralentit l'anoxie cérébrale} et, crucialement, \textbf{perfuse le myocarde} (coronaires) maintenant la fibrillation ventriculaire (FV) sous forme « fine » (amplitude faible, fréquence élevée) qui est \textbf{sensible au choc électrique}. Sans RCP, la FV dégénère en asystolie ou FV grossière (réfractaire au choc) en \SIrange{4}{6}{min}.
\item \textbf{Maillon 3 (DAE) :} La cause n°1 de l'ACA de l'adulte est la \textbf{fibrillation ventriculaire}. Le \textbf{seul traitement efficace} est le \textbf{choc électrique} (dépolarisation massive $\rightarrow$ reprise rythme sinusal par nœud sinusal). La probabilité de succès du choc \textbf{décroît de \SI{10}{\%} par minute} de retard. Sans RCP, elle est quasi nulle à \SI{10}{min}.
\item \textbf{Maillon 4 (Soins avancés) :} Traite la cause (infarctus, trouble électrolytique, etc.), stabilise l'hémodynamique, prévient la récidive (hypothermie thérapeutique, coronarographie).
\end{itemize}
\textbf{Interdépendance :} Une alerte tardive retarde le DAE ; une RCP tardive rend le myocarde réfractaire au choc ; un DAE tardif laisse le cerveau en anoxie prolongée. La rapidité d'exécution de \emph{chaque} maillon conditionne l'efficacité de toute la chaîne.
\textbf{Conclusion} : la chaîne de survie n'est forte que si ses maillons s'enclenchent dans les \textbf{premières minutes}. C'est pourquoi la formation de \textbf{tous les citoyens} aux gestes qui sauvent (RCP, DAE, PLS, étouffement) est un enjeu majeur de santé publique, y compris au Cameroun où les délais d'arrivée des secours médicalisés (SAMU/Pompiers) peuvent dépasser \SI{15}{min} en zone urbaine et bien plus en zone rurale.
\end{exoresolu}

\begin{attention}[Les 5 erreurs qui coûtent des points au Bac (et des vies dans la réalité)]
\begin{enumerate}
\item \textbf{Confondre perte de connaissance et arrêt cardiaque.} Le critère discriminant unique est la \textbf{respiration} : Inconscient \textbf{qui respire} $\rightarrow$ \textbf{PLS} ; Inconscient \textbf{qui ne respire pas} $\rightarrow$ \textbf{RCP}. Proposer un massage cardiaque pour une victime qui respire est une \textbf{erreur grave} (invalidation de la réponse, geste dangereux).
\item \textbf{Oublier l'alerte des secours (119/118)} ou la placer \textbf{après} les gestes de secours. L'alerte fait partie intégrante de la conduite à tenir (maillon 1). Il faut citer le \textbf{numéro}, les \textbf{infos à transmettre} (lieu exact, nombre, état, gestes en cours) et l'usage du \textbf{haut-parleur}.
\item \textbf{Donner des valeurs techniques fausses ou absentes.} Au Bac, « masser le cœur » ne suffit pas. Il faut écrire : \textbf{30 compressions / 2 insufflations}, \textbf{100 à 120 compressions/min}, \textbf{profondeur \SI{5}{cm} (5-6 cm)}, \textbf{centre de la poitrine (sternum moitié inférieure)}. Idem étouffement : \textbf{5 claques / 5 compressions}.
\item \textbf{Appliquer la manœuvre de Heimlich (compressions abdominales) à toutes les victimes.} \textbf{Interdit} chez le \textbf{nourrisson} (risque rupture hépatique) et la \textbf{femme enceinte} (utérus gravide, inefficacité). On substitue des \textbf{compressions thoraciques}. Oublier cette adaptation est une \textbf{faute de cours éliminatoire}.
\item \textbf{Proposer des gestes dangereux ou non codifiés :} donner à boire/manger à un inconscient (fausse route) ; chercher le pouls carotidien (perte de temps, non fiable) ; interrompre la RCP pour « vérifier si ça marche » (arrêt flux coronarien) ; déplacer une victime sans danger imminent (aggravation traumatisme rachidien). Tout geste proposé doit être \textbf{sûr, codifié, justifié physiologiquement}.
\end{enumerate}
\end{attention}

\newpage

\coursec{Exercices}

\courssub{Je vérifie mes connaissances}

\begin{aretenir}[Rappels indispensables : La Chaîne de Survie]
La \cle{chaîne de survie} (ERC/AHA 2021) comporte 4 maillons indissociables, dont la rapidité d'exécution conditionne le pronostic vital :
\begin{enumerate}
    \item \textbf{Alerte précoce} : Appel immédiat du 112 (ou 119/118 au Cameroun) dès la reconnaissance de l'arrêt cardiaque.
    \item \textbf{RCP précoce} : Massage cardiaque externe (MCE) commencé dans les 1\textsuperscript{re} minute, idéalement par un témoin formé (PSC1).
    \item \textbf{Défibrillation précoce} : Utilisation d'un DAE (Défibrillateur Automatisé Externe) public le plus vite possible (cible < 3--5 min).
    \item \textbf{Prise en charge médicale spécialisée précoce} : Intervention du SMUR / réanimation (médicaments, ventilation invasive, soins post-réanimation).
\end{enumerate}
\textbf{Règle d'or :} Tout maillon manqué ou retardé réduit drastiquement la survie. \emph{Le temps, c'est du myocarde, c'est du cerveau.}
\end{aretenir}

\begin{exercice}[QCM : Identification des situations d'urgence]
Parmi les situations suivantes, laquelle ou lesquelles nécessitent \textbf{immédiatement} le début d'une réanimation cardio-pulmonaire (RCP) avec compressions thoraciques ?
\begin{enumerate}
    \item Une victime inconsciente, qui ne respire pas, après une électrocution domestique à Douala.
    \item Une victime consciente, qui tousse violemment après avoir avalé un os de poisson lors d'un repas de \emph{ndolé}.
    \item Une victime inconsciente, qui respire normalement (mouvements thoraciques réguliers, 14 cycles/min), trouvée au bord de la route après un accident de moto-taxi.
    \item Une victime inconsciente, qui ne respire pas, mais présente des mouvements respiratoires agoniaques (gasps lents, bruyants, irréguliers) après un malaise cardiaque.
    \item Un nourrisson de 8 mois, conscient, qui ne peut plus émettre de son, ne tousse plus et devient cyanotique après avoir mis un grain de maïs dans sa bouche.
\end{enumerate}
\emph{Pour chaque situation, justifiez votre choix en citant les critères de reconnaissance de l'arrêt cardiaque ou de l'étouffement total.}
\end{exercice}

\begin{exercice}[Vrai ou Faux justifié : Chaîne de survie et physiologie]
Indiquez si les affirmations suivantes sont vraies ou fausses. Justifiez chaque réponse par un argument physiologique ou procédural précis.
\begin{enumerate}
    \item Le massage cardiaque externe (MCE) permet de restituer un débit cardiaque équivalent à 60\,\% du débit normal au repos.
    \item La défibrillation précoce est le seul maillon de la chaîne de survie capable de traiter une fibrillation ventriculaire (FV) ou une tachycardie ventriculaire (TV) sans pouls.
    \item Chez la femme enceinte en fin de grossesse (au-delà de 20 SA), les compressions thoraciques doivent être réalisées au niveau de l'appendice xiphoïde pour éviter de comprimer le fœtus.
    \item La position latérale de sécurité (PLS) est indiquée pour toute victime inconsciente, qu'elle respire ou non.
    \item Chez le nourrisson de moins d'un an, la manœuvre de Heimlich (compressions abdominales) est formellement contre-indiquée en raison du risque de lésion hépatique ou splénique.
\end{enumerate}
\end{exercice}

\begin{attention}[Pièges classiques du QCM et V/F]
\begin{itemize}
    \item \textbf{Gasps $\neq$ Respiration :} Les gasps (respiration agoniaque) sont des mouvements saccadés, bruyants, lents ($<$ 6/min), signe d'ischémie cérébrale sévère \textbf{équivalant à l'apnée}. $\rightarrow$ \textbf{Déclenchent la RCP}.
    \item \textbf{MCE $\neq$ 60\% DC :} Le MCE manuel génère $\approx$ 30\,\% du débit cardiaque normal (mécanisme pompe thoracique). Seule la RCP de haute qualité (fréquence, profondeur, relâchement, fraction de compression $>$ 80\,\%) approche ce chiffre.
    \item \textbf{Femme enceinte :} Compressions au \textbf{centre de la poitrine} (moitié inférieure du sternum), \emph{pas} sur l'appendice xiphoïde (risque hépatectomie). Déclive utérin gauche manuel si 2\textsuperscript{e} secouriste.
    \item \textbf{PLS :} \textbf{Uniquement si respiration normale.} Inconscient + apnée/gasps = Arrêt cardiaque = RCP.
    \item \textbf{Nourrisson :} Heimlich (abdominal) \textbf{interdit} $<$ 1 an (foie/splène fragiles, cage thoracique souple). $\rightarrow$ \textbf{5 tapes dorsales + 5 compressions thoraciques} (2 doigts, ligne intermammaire).
\end{itemize}
\end{attention}

\begin{exercice}[Définitions et seuils critiques]
Définissez précisément les termes suivants en utilisant le vocabulaire scientifique attendu au Baccalauréat :
\begin{enumerate}
    \item Arrêt cardiaque (critères cliniques de reconnaissance).
    \item Étouffement total (ou obstruction complète des voies aériennes supérieures).
    \item Gasps (ou respiration agoniaque).
    \item Position latérale de sécurité (PLS) : objectif physiologique principal.
    \item Chaîne de survie : citez les 4 maillons dans l'ordre chronologique.
\end{enumerate}
Indiquez les valeurs seuils critiques (avec unités) pour :
\begin{enumerate}
    \setcounter{enumi}{5}
    \item La fréquence des compressions thoraciques chez l'adulte.
    \item La profondeur des compressions thoraciques chez l'adulte.
    \item Le ratio compressions / insufflations pour un secouriste seul chez l'adulte.
    \item Le délai maximal au-delà duquel les lésions cérébrales deviennent irréversibles en l'absence de RCP.
\end{enumerate}
\end{exercice}

\begin{aretenir}[Seuils critiques de la RCP Adulte (ERC 2021 / Guidelines Bac)]
\begin{center}
\begin{tabular}{|l|c|c|}
\hline
\textbf{Paramètre} & \textbf{Valeur cible} & \textbf{Justification physiologique} \\
\hline
Fréquence & 100 à 120 / min & Optimise le débit systolique sans compromettre le remplissage diastolique \\
\hline
Profondeur & 5 à 6 cm & Génère une pression de perfusion coronarienne (PPC) suffisante ($>$ 15--20 mmHg) \\
\hline
Relâchement & Complet (poitrine remonte) & Permet le retour veineux (précharge) et la perfusion coronaire (diastole) \\
\hline
Ratio (1 secouriste) & 30 compressions / 2 insufflations & Minimise les interruptions ; 2 insufflations $\approx$ 4--5 s \\
\hline
Ratio (2 secouristes) & 15 compressions / 2 insufflations (enfant/nourrisson) & Besoins ventilatoires relatifs plus élevés chez l'enfant \\
\hline
Fraction de compression (CCF) & $>$ 80\,\% & Temps mains sur thorax / temps total ; chaque interruption chute PPC $\approx$ 0 \\
\hline
Délai lésions cérébrales & 3 à 5 min sans flux & Réserve O$_2$ cérébrale épuisée $\rightarrow$ ischémie irréversible (noyau gris) \\
\hline
\end{tabular}
\end{center}
\end{aretenir}

\begin{exercice}[Calcul direct : Impact du délai d'intervention]
On admet que, \textbf{sans RCP}, la probabilité de survie après un arrêt cardiaque extra-hospitalier diminue d'environ 10\,\% par minute (modèle exponentiel $S(t) = 100 \times 0,9^t$). \textbf{Avec RCP immédiate de qualité}, la décroissance est ralentie à $\approx$ 3--4\,\% par minute ($S_{\text{RCP}}(t) \approx 100 \times 0,96^t$).
\begin{enumerate}
    \item Calculez la probabilité de survie théorique \textbf{sans RCP} si elle est commencée 4 minutes après l'effondrement ($t=0$ : 100\,\%).
    \item Calculez la probabilité de survie théorique \textbf{avec RCP immédiate} (délai 0 min) maintenue jusqu'à l'arrivée des secours médicalisés (SMUR) à 8 minutes.
    \item En déduire l'intérêt majeur de la formation du grand public au secourisme (PSC1) dans le contexte camerounais où le délai moyen d'arrivée des secours peut dépasser 15 minutes en zone urbaine dense (Yaoundé, Douala) et 30 minutes en zone rurale.
\end{enumerate}
\end{exercice}

\courssub{Je m'entraîne}

\begin{methode}[Conduite type : Arrêt Cardiaque Adulte (Témoin unique, DAE dispo)]
\begin{enumerate}
    \item \textbf{Sécuriser} : Danger ? (Courant, circulation, feu). \cle{Ne pas devenir une victime de plus}.
    \item \textbf{Évaluer conscience} : Appel verbal + stimulation tactile (épaules). \cle{Pas de réponse = Inconscient}.
    \item \textbf{Alerter (Phone First Adulte)} : Appeler 112/119 \textbf{avant} de commencer la RCP (sauf noyade/asphyxie/enfant : 1 min RCP puis alerte). Envoyer chercher DAE si proche.
    \item \textbf{Libérer voies aériennes} : Déclive menton / élévation menton (pas d'hyperextension si traumatisme suspecté).
    \item \textbf{Évaluer respiration} (10 sec max) : \textbf{Regarder} (thorax/abdomen), \textbf{Écouter}, \textbf{Sentir} (joue). \cle{Gasps = Pas de respiration normale}.
    \item \textbf{Si pas de respiration normale} : \textbf{Débuter RCP immédiatement} (30 compressions / 2 insufflations).
    \item \textbf{MCE technique} : Talon d'une main au centre de la poitrine (ligne intermammaire), 2\textsuperscript{e} main par-dessus, bras tendus, verticalité, 5--6 cm, 100--120/min, relâchement total.
    \item \textbf{DAE arrivé} : Allumer $\rightarrow$ Coller électrodes (schéma sur boîtier) $\rightarrow$ \textbf{Ne plus toucher la victime} (Analyse) $\rightarrow$ Si « Choc recommandé » : \textbf{Éloigner tout le monde} $\rightarrow$ Appuyer « Choc » $\rightarrow$ \textbf{Reprendre MCE immédiatement} 2 min (5 cycles) $\rightarrow$ Réanalyse.
    \item \textbf{Continuer} jusqu'à : Reprise activité vitale (mouvements, toux, respiration normale), Arrivée secours médicalisés, Épuisement secouriste, Danger.
\end{enumerate}
\end{methode}

\begin{exercice}[Conduite à tenir : Arrêt cardiaque chez l'adulte]
Vous êtes témoin d'un malaise sur le terrain de handball du lycée : un élève de Terminale D s'effondre brutalement pendant l'échauffement. Il est au sol, ne bouge pas, ne répond pas à votre appel ni à la stimulation tactile (épaules). Vous ouvrez les voies aériennes (déclive menton) et observez : pas de mouvement thoracique, pas de bruit respiratoire, pas de souffle sur votre joue pendant 10 secondes. Vous notez cependant des mouvements respiratoires saccadés, bruyants, toutes les 10-15 secondes.
\begin{enumerate}
    \item Quel diagnostic posez-vous ? Justifiez en distinguant respiration normale et gasps.
    \item Rédigez la conduite à tenir complète et séquencée (numérotez les étapes) depuis l'alerte jusqu'à l'arrivée des secours, en précisant la technique du MCE (position des mains, fréquence, profondeur, relâchement) et la gestion du DAE si disponible dans l'établissement.
    \item Pourquoi ne doit-on \textbf{pas} pratiquer de bouche-à-bouche si l'on n'est pas équipé (masque, filtre) et que l'on est seul ? Quelle alternative recommandée par les guidelines internationales (ERC/AHA) ?
\end{enumerate}
\end{exercice}

\begin{exercice}[Cas pratique : Étouffement total - Différenciation selon la victime]
Complétez le tableau comparatif suivant en justifiant chaque geste par l'anatomie ou la physiologie de la victime.
\begin{center}
\begin{tabularx}{\linewidth}{|l|X|X|X|}
\hline
\rowcolor{popblueL}
\textbf{Critère} & \textbf{Nourrisson ($<$ 1 an)} & \textbf{Adulte / Enfant $>$ 1 an} & \textbf{Femme enceinte (ou obèse)} \\
\hline
Signes cliniques d'obstruction totale & & & \\
\hline
1\textsuperscript{ère} manœuvre (si conscient) & & & \\
\hline
2\textsuperscript{ème} manœuvre (si 1\textsuperscript{ère} échoue) & & & \\
\hline
Manœuvre \textbf{interdite} & Compressions abdominales (Heimlich) & Aucune (si protocole respecté) & Compressions abdominales classiques \\
\hline
Justification anatomique de l'interdiction & & & \\
\hline
Conduite si la victime devient inconsciente & & & \\
\hline
\end{tabularx}
\end{center}
\end{exercice}

\begin{methode}[Algorithme Étouffement Total Conscient]
\begin{enumerate}
    \item \textbf{Reconnaître} : Signe de l'étranglement (mains au cou), \textbf{parole impossible}, \textbf{toux absente/inefficace}, \textbf{stridor} (inspiration bruyante), \textbf{cyanose}, angoisse. \cle{Demander : "Vous étouffez ?" $\rightarrow$ Hocher la tête = Oui}.
    \item \textbf{Appeler à l'aide} (ne pas quitter la victime).
    \item \textbf{Adulte/Enfant $>$ 1 an} : 5 \textbf{tapes dorsales} (paume, entre omoplates, victim penché en avant) $\rightarrow$ Si échec : 5 \textbf{compressions abdominales} (Heimlich : poing au-dessus nombril, sous xiphoïde, traction postéro-supérieure). Alterner jusqu'à expulsion ou inconscience.
    \item \textbf{Nourrisson $<$ 1 an} : 5 \textbf{tapes dorsales} (avant-bras sur cuisse, tête plus bas que corps, paume entre omoplates) $\rightarrow$ Si échec : 5 \textbf{compressions thoraciques} (2 doigts, ligne intermammaire, profondeur 1/3 AP). \textbf{Jamais abdominal}. Alterner.
    \item \textbf{Femme enceinte / Obèse} : 5 \textbf{tapes dorsales} $\rightarrow$ Si échec : 5 \textbf{compressions thoraciques} (haut sternum, comme MCE) $\rightarrow$ Alterner.
    \item \textbf{Devient inconscient} : Allonger au sol $\rightarrow$ \textbf{Appeler secours} $\rightarrow$ \textbf{Débuter RCP} (30 compressions). \textbf{Avant chaque insufflation : regarder dans la bouche} (retirer corps étranger visible si accessible, pas de doigt à l'aveugle).
\end{enumerate}
\end{methode}

\begin{exercice}[Mise en œuvre de la PLS : Technique et surveillance]
Une victime adulte est trouvée inconsciente dans la cour de l'école après une chute. Elle respire normalement (fréquence 16/min, amplitude correcte). Aucun traumatisme rachidien n'est suspecté.
\begin{enumerate}
    \item Décrivez pas à pas la technique de mise en PLS (côté droit) pour un secouriste seul, en insistant sur : la position du bras droit (angle), la main gauche (joue), la cuisse gauche (angle), le basculement, l'ouverture de la bouche, l'inclinaison de la tête en arrière.
    \item Quels sont les 3 objectifs physiologiques de la PLS ?
    \item Quelle surveillance continue devez-vous assurer pendant l'attente des secours ? À quelle fréquence devez-vous réévaluer la respiration ?
    \item Si la victime vomit pendant qu'elle est en PLS, quelle action immédiate devez-vous réaliser ?
\end{enumerate}
\end{exercice}

\begin{aretenir}[PLS : Technique standardisée (Côté droit) \& Surveillance]
\textbf{Étapes (Secouriste seul, côté droit) :}
\begin{enumerate}
    \item Victime sur le dos, jambes tendues. \textbf{Bras droit} : perpendiculaire au corps, coude fléchi 90° (main vers tête).
    \item \textbf{Bras gauche} : traverser le thorax, \textbf{dos de la main gauche contre joue droite} de la victime (maintien tête).
    \item \textbf{Jambe gauche} : fléchir genou, pied au sol (point d'appui).
    \item \textbf{Basculer} : Tirer sur genou gauche $\rightarrow$ Victime roule sur le côté droit (garder main gauche sous joue pour contrôler tête).
    \item \textbf{Ajuster} : Cuisse gauche fléchie à 90° (hanche/genou) pour stabiliser. Tête en \textbf{hyperextension modérée} (menton relevé) $\rightarrow$ \textbf{Ouverture bouche} (pouce/indice) pour drainage.
\end{enumerate}
\textbf{3 Objectifs physiologiques :} 1. \cle{Permeabilité VAS} (langue ne tombe pas en arrière). 2. \cle{Drainage gravitaire} (salive, sang, vomissures $\rightarrow$ extérieur). 3. \cle{Prévention fausse route} (poumons protégés).
\textbf{Surveillance :} Respiration (regarder/écouter/sentir) \textbf{toutes les 1 à 2 minutes}. Vérifier conscience (parler/toucher).
\textbf{Vomissements en PLS :} \textbf{Maintenir en PLS}, ouvrir la bouche, laisser s'écouler (gravité), essuyer si nécessaire. \textbf{Ne pas redresser}.
\end{aretenir}

\begin{exercice}[Alerte aux secours : Message structuré]
Vous êtes seul témoin d'un accident de la circulation sur l'autoroute Yaoundé-Douala (pk 45). Il y a deux victimes :
\begin{itemize}
    \item Victime 1 (conducteur) : Inconsciente, ne respire pas (arrêt cardiaque confirmé).
    \item Victime 2 (passager) : Inconsciente, respire normalement, saignement abondant au niveau de la cuisse gauche (fracture ouverte fémorale suspectée).
\end{itemize}
\begin{enumerate}
    \item Rédigez le message d'alerte complet en respectant la structure standardisée (Qui ? Où ? Quoi ? Combien ? État ? Mesures prises ?).
    \item Justifiez l'ordre de priorité de vos actions : quelle est la \textbf{première action} à réaliser seul face à un adulte en arrêt cardiaque ? Pourquoi la victime 2 est-elle gérée après l'alerte et le début de la RCP sur la victime 1 ?
    \item Si un deuxième témoin arrive pendant que vous faites le MCE sur la victime 1, comment répartissez-vous les tâches ?
\end{enumerate}
\end{exercice}

\begin{methode}[Message d'Alerte Structuré (Norme Européenne / Cameroun)]
\textbf{Numéro :} 112 (GSM/International), 119 (Pompiers Cameroun), 118 (SAMU Yaoundé).
\textbf{Script mémorisable :} \cle{« \textbf{OÙ} ? \textbf{QUOI} ? \textbf{COMBIEN} ? \textbf{ÉTAT} ? \textbf{MESURES} ? \textbf{NE RACCROCHEZ PAS LE PREMIER.} »}
\begin{enumerate}
    \item \textbf{OÙ} : Localisation précise (PK 45 autoroute Ydé-Dla, sens Ydé$\rightarrow$Dla, repères visibles).
    \item \textbf{QUOI} : Accident circulation, 2 véhicules ? (préciser).
    \item \textbf{COMBIEN} : 2 victimes.
    \item \textbf{ÉTAT} : V1 : Inconsciente, \textbf{ne respire pas} (Arrêt Cardiaque \textbf{PRIORITÉ ABSOLUE}). V2 : Inconsciente, respire, hémorragie cuisse gauche (comprimée).
    \item \textbf{MESURES} : RCP en cours sur V1 (MCE 30:2). PLS + Compression manuelle hémorragie V2.
    \item \textbf{VOTRE NOM / NUMÉRO} : Pour rappel.
\end{enumerate}
\textbf{Règle priorisation (Triage simple) :} \cle{Arrêt cardiaque (Arrêt circulatoire) > Hémorragie massive > Inconscient respirant (PLS) > Blessés conscients}.
\end{methode}

\courssub{Je me prépare au Bac}

\begin{exercice}[Sujet type Bac : Analyse de la chaîne de survie et modélisation]
\textbf{Contexte :} Au Cameroun, les maladies cardiovasculaires sont la deuxième cause de mortalité après le paludisme. La mort subite de l'adulte survient le plus souvent en pré-hospitalier. La « chaîne de survie » comprend 4 maillons : 1. Alerte précoce, 2. RCP précoce, 3. Défibrillation précoce, 4. Prise en charge médicale spécialisée précoce.

\textbf{Document 1 :} Courbe de survie en fonction du délai de défibrillation (données ERC 2021).
\begin{popfigure}
\begin{tikzpicture}
\begin{axis}[
    width=\linewidth, height=8cm,
    xlabel={Délai avant défibrillation (min)},
    ylabel={Taux de survie à la sortie d'hôpital (\%)},
    xmin=0, xmax=12, ymin=0, ymax=80,
    xtick={0,2,4,6,8,10,12}, ytick={0,10,20,30,40,50,60,70},
    grid=major, grid style={popline},
    axis lines=middle,
    enlargelimits={abs=0.5},
    tick label style={font=\small},
    label style={font=\small},
]
\addplot[popcurve, very thick, domain=0:10, samples=100] {70*exp(-0.25*x)};
\addplot[poporange, dashed, thick] coordinates {(0,70) (10,5.75)};
\addplot[popgreen, mark=*, only marks, mark size=2.5] coordinates {(1,55) (3,33) (5,20) (7,12) (9,7)};
\legend{Modèle $S(t)=70e^{-0.25t}$, $S(0)=70\%$ (théorique), Données cliniques observées}
\end{axis}
\end{tikzpicture}
\legende{Évolution du taux de survie selon le délai de défibrillation. La courbe continue modélise $S(t) = 70 e^{-0,25t}$ ($t$ en min). Les points verts représentent des données cliniques observées (ERC 2021). La courbe en pointillés montre la survie maximale théorique si défibrillation immédiate.}
\end{popfigure}

\textbf{Document 2 :} Extrait des recommandations ERC 2021.
« Le massage cardiaque externe (MCE) maintient un flux sanguin cérébral minimal (environ 30\,\% du normal) grâce au mécanisme de la pompe thoracique. L'interruption du MCE pour insufflation ou analyse DAE fait chuter la pression de perfusion coronarienne (PPC) à zéro ; sa remontée prend plusieurs compressions. »

\begin{enumerate}
    \item \textbf{Analyse de données :} À partir du modèle mathématique $S(t) = 70 e^{-0,25t}$ (avec $t$ en minutes), calculez le taux de survie théorique pour un délai de 3 min et de 8 min. Comparez avec les points expérimentaux du Document 1. Commentez la validité du modèle.
    \item \textbf{Calcul de perte :} Déterminez la perte moyenne de chances de survie par minute (en \% par minute) entre $t=0$ et $t=4$ min à partir d'une moyenne sur l'intervalle. Interprétez ce résultat pour la stratégie de formation du public.
    \item \textbf{Physiologie :} En vous appuyant sur le Document 2 et vos connaissances, expliquez pourquoi :
    \begin{enumerate}
        \item Le MCE doit être commencé \textbf{immédiatement} et \textbf{interrompu le moins possible}.
        \item La défibrillation précoce (maillon 3) est inefficace si le maillon 2 (RCP) n'a pas été réalisé pendant l'attente du DAE.
    \end{enumerate}
    \item \textbf{Contexte camerounais :} Proposez deux mesures concrètes, adaptées au contexte local (écoles, marchés, motos-taxis, délais secours), pour renforcer les deux premiers maillons de la chaîne de survie.
\end{enumerate}
\end{exercice}

\begin{exercice}[Sujet type Bac : Cas clinique intégré - Arrêt cardiaque vs Étouffement vs PLS]
Trois situations distinctes surviennent lors d'une kermesse scolaire à Bafoussam. Pour chaque situation, vous devez : \textbf{1. Poser le diagnostic}, \textbf{2. Justifier par les signes cliniques}, \textbf{3. Rédiger la conduite à tenir complète et ordonnée} (y compris l'alerte et la prise en charge spécifique).

\textbf{Situation A :} Un homme de 55 ans, connu hypertendu, s'assoit brutalement en tenant sa poitrine. Il devient pâle, transpire abondamment, ne répond plus. Vous vous approchez : il ne réagit pas à la parole ni au pincement du trapèze. Voies aériennes dégagées. Aucune ventilation perçue en 10 sec (ni thoracique, ni abdominale, ni souffle). Pas de pouls carotidien palpé (secouriste non médecin). Pas de DAE sur site. Secours à 15 min.

\textbf{Situation B :} Un enfant de 3 ans court en mangeant une arachide. Il s'arrête net, porte les mains à la gorge (signe de l'étranglement), ne peut pas parler, ne tousse pas, devient bleu (cyanose lèvres/ongles), regard effrayé. Il est conscient.

\textbf{Situation C :} Une jeune fille de 17 ans s'effondre après être restée debout longtemps au soleil (38°C). Elle est inconsciente au sol. Vous l'appelez, la stimulez : pas de réponse. Vous ouvrez la bouche : pas de corps étranger. Vous libérez les voies aériennes : vous entendez un bruit de ronflement (râle) à l'inspiration, vous voyez la cage thoracique se soulever régulièrement (18 mov/min), vous sentez le souffle sur votre joue. Pouls radial présent (80/min, régulier). Pas de traumatisme crânien apparent.

\begin{enumerate}
    \item Pour chaque situation, remplissez le tableau de synthèse suivant :
    \begin{center}
    \begin{tabularx}{\linewidth}{|l|X|X|X|}
    \hline
    \rowcolor{popblueL}
    & \textbf{Situation A} & \textbf{Situation B} & \textbf{Situation C} \\
    \hline
    Diagnostic & & & \\
    \hline
    Signes cliniques clés (justification) & & & \\
    \hline
    Conduite à tenir (étapes numérotées) & & & \\
    \hline
    \end{tabularx}
    \end{center}
    \item Situation A : Au bout de 2 minutes de RCP efficace, un DAE arrive. L'appareil indique « Choc recommandé ». Décrivez la séquence exacte des actions (sécurité, choc, reprise RCP).
    \item Situation B : L'enfant devient inconscient pendant que vous effectuez les manœuvres. Quelle modification immédiate de la conduite à tenir ?
    \item Situation C : Pourquoi la PLS est-elle réalisée du côté droit de préférence (convention pédagogique) ? Quel risque spécifique surveillez-vous en priorité chez cette victime (contexte : chaleur, station debout prolongée) ?
\end{enumerate}
\end{exercice}

\begin{exercice}[Sujet type Bac : Rédaction argumentée - Physiologie de la réanimation]
\textbf{Sujet :} « Montrer, à partir des besoins en dioxygène du cerveau et des mécanismes hémodynamiques du massage cardiaque externe, pourquoi la réanimation cardio-pulmonaire (RCP) doit être entreprise immédiatement, poursuivie sans interruption majeure, et associée à une défibrillation la plus précoce possible en cas de fibrillation ventriculaire. »

Rédigez une réponse structurée en trois parties argumentées, en mobilisant vos connaissances de physiologie (métabolisme cérébral, autorégulation, pression de perfusion coronarienne) et de secourisme (mécanismes du MCE, rôle du DAE, gestion des interruptions). Utilisez un schéma fonctionnel simplifié (boucle de régulation ou flux) pour illustrer votre démonstration.
\begin{enumerate}
    \item \textbf{Besoins cérébraux et fenêtre thérapeutique :} Quantifiez la consommation d'O$_2$ cérébrale (en mL/100g/min ou \% du total), la réserve en O$_2$ (hémoglobine, myoglobine, dissous), et le délai d'apparition de l'ischémie irréversible. Expliquez pourquoi le « temps cerveau » est le facteur limitant.
    \item \textbf{Mécanique du MCE et hémodynamique :} Décrivez les deux théories (pompe cardiaque vs pompe thoracique) expliquant le flux sanguin pendant le MCE. Définissez la Pression de Perfusion Coronarienne (PPC = Pression Aortique Diastolique - Pression Atriale Droite / Pression VG Diastolique). Expliquez pourquoi l'interruption des compressions fait chuter la PPC à zéro et pourquoi sa remontée est lente.
    \item \textbf{Synergie RCP - Défibrillation :} Expliquez pourquoi une FV « non traitée » évolue vers une asystolie. Montrez que la RCP « prépare » le myocarde à la défibrillation en maintenant un métabolisme aérobie minimal (substrats, pH, ATP). Concluez sur l'optimisation du ratio « temps mains sur le thorax / temps total » (Fraction de Compression Thoracique > 80\,\%).
\end{enumerate}
\end{exercice}

\begin{savaistu}[Secourisme au Cameroun : Realités et Acteurs]
\begin{itemize}
    \item \textbf{Numéros d'urgence :} \cle{112} (GSM, standard international, redirige vers CPC), \cle{119} (Pompiers / Sécurité Civile), \cle{118} (SAMU Yaoundé). \emph{Conseil : Enregistrer 112 et 119 dans « Favoris » téléphone.}
    \item \textbf{Formation PSC1 (Prévention et Secours Civiques niveau 1) :} Dispensée par la \textbf{Croix-Rouge Camerounaise}, la \textbf{Protection Civile}, certains lycées (club secourisme) et entreprises. Durée : 7--10 h. Diplôme reconnu.
    \item \textbf{DAE publics :} Encore rares (aéroports, certains hôtels/banques, ambassades). Plaidoyer pour installation dans les \textbf{gares routières, marchés (Mokolo, Central), agences de voyage (Voyages Express, etc.), lycées}.
    \item \textbf{Motos-taxis (Benskin) :} Premiers témoins fréquents. Formation ciblée (casque, gilet, trousse, PSC1) = levier majeur pour maillon 1 \& 2.
    \item \textbf{Bonne Samaritaine (Loi) :} Toute personne qui porte assistance à une personne en danger ne peut être poursuivie pour les conséquences de son acte (sauf faute intentionnelle ou caractérisée). \cle{Ne pas hésiter à agir}.
\end{itemize}
\end{savaistu}

\newpage

\coursec{Corrigés des exercices}

\coursec{Corrigés des exercices — Pratique du secourisme pour la prévention des accidents cardiovasculaires}

\begin{corrige}[Exercice 1 — QCM : Identification des situations d'urgence]
\textbf{Réponses : les situations 1 et 4 nécessitent une RCP immédiate.}

\begin{enumerate}
    \item \textbf{RCP immédiate — OUI.} Critères de l'arrêt cardiaque réunis : victime \cle{inconsciente} (ne répond pas) ET \cle{absence de respiration}. L'électrocution est une cause classique d'arrêt cardiaque (fibrillation ventriculaire par passage du courant à travers le myocarde). \textbf{Condition préalable :} couper la source de courant avant tout contact (sécurité du secouriste), puis alerter et masser.
    \item \textbf{RCP — NON.} La victime est \textbf{consciente} et \textbf{tousse violemment} : il s'agit d'une obstruction \emph{partielle} des voies aériennes avec toux efficace. La toux est le mécanisme le plus efficace d'expulsion (pics de pression intrathoracique élevés). Conduite : encourager à tousser, surveiller, ne rien faire d'autre. Toute manœuvre (tapes, Heimlich) serait ici injustifiée et dangereuse.
    \item \textbf{RCP — NON.} Victime inconsciente mais \textbf{respirant normalement} (\SI{14}{cycles/min}, régulier) : ce n'est pas un arrêt cardiaque, c'est une \cle{perte de connaissance}. Conduite : alerte + \textbf{PLS} (sauf suspicion de traumatisme rachidien — ici accident de moto-taxi : minimiser la mobilisation, maintien tête-cou, PLS adaptée si nécessaire pour libérer les voies aériennes).
    \item \textbf{RCP immédiate — OUI.} Les \cle{gasps} (respiration agoniaque : mouvements lents, bruyants, irréguliers, $<$ \SI{6}{min}) ne sont \textbf{pas une respiration efficace} : ils traduisent une ischémie cérébrale majeure et équivalent à une apnée. Inconscient + gasps = \textbf{arrêt cardiaque}. C'est le piège classique : environ \SI{40}{\%} des arrêts cardiaques présentent des gasps dans les premières minutes.
    \item \textbf{RCP — NON (pour l'instant).} Nourrisson \textbf{conscient} avec obstruction \textbf{totale} (silence, pas de toux, cyanose) : conduite = \textbf{5 tapes dorsales + 5 compressions thoraciques} en alternance (jamais de Heimlich $<$ 1 an). La RCP ne débute que si le nourrisson devient \textbf{inconscient}.
\end{enumerate}
\end{corrige}

\begin{corrige}[Exercice 2 — Vrai ou Faux justifié]
\begin{enumerate}
    \item \textbf{FAUX.} Le MCE manuel, même bien réalisé, ne restitue qu'environ \textbf{\SI{30}{\%} du débit cardiaque normal} (mécanisme de pompe thoracique/cardiaque). C'est un flux de \emph{survie}, suffisant pour préserver le cerveau et le cœur quelques dizaines de minutes, mais pas un débit physiologique. Les \SI{60}{\%} ne sont approchés que par des dispositifs mécaniques dans des conditions expérimentales.
    \item \textbf{VRAI.} La fibrillation ventriculaire (FV) et la tachycardie ventriculaire sans pouls sont des troubles du rythme \textbf{« choquables »} : seul un \textbf{choc électrique externe} (défibrillation) peut resynchroniser les cellules myocardiques (dépolarisation massive simultanée suivie d'une période réfractaire permettant au nœud sinusal de reprendre le commandement). Ni le MCE ni les insufflations ne convertissent une FV ; le MCE ne fait que \emph{maintenir} le myocarde en attente du choc.
    \item \textbf{FAUX.} Les compressions se font toujours au \textbf{centre de la poitrine} (moitié inférieure du sternum), jamais sur l'appendice xiphoïde (risque de lacération hépatique). Chez la femme enceinte $>$ \SI{20}{SA}, la mesure spécifique est le \textbf{déplacement manuel de l'utérus vers la gauche} (ou déclive latérale gauche) pour décomprimer la veine cave inférieure et restaurer le retour veineux.
    \item \textbf{FAUX.} La PLS n'est indiquée que pour la victime \textbf{inconsciente qui respire normalement}. Si la victime inconsciente ne respire pas (ou gasps), c'est un arrêt cardiaque : elle doit rester \textbf{sur le dos} pour la RCP. Mettre en PLS une victime en arrêt cardiaque ferait perdre un temps vital.
    \item \textbf{VRAI.} Chez le nourrisson $<$ 1 an, le foie et la rate sont volumineux relativement à l'abdomen, peu protégés par une cage thoracique souple ; les compressions abdominales (Heimlich) exposent à des \textbf{lésions hépatiques ou spléniques}. Protocole : 5 tapes dorsales puis 5 compressions \emph{thoraciques} (2 doigts, ligne intermammaire).
\end{enumerate}
\end{corrige}

\begin{corrige}[Exercice 3 — Définitions et seuils critiques]
\textbf{Définitions :}
\begin{enumerate}
    \item \textbf{Arrêt cardiaque :} cessation brutale de l'activité mécanique efficace du cœur, diagnostiquée cliniquement par la triade : victime \textbf{inconsciente} (pas de réponse à la parole et à la stimulation), \textbf{absence de respiration normale} (apnée ou gasps) après libération des voies aériennes, et — pour les secouristes formés — \textbf{absence de pouls carotidien}.
    \item \textbf{Étouffement total :} obstruction complète des voies aériennes supérieures par un corps étranger, se traduisant chez la victime consciente par : impossibilité de parler, toux absente ou inefficace, signe de l'étranglement (mains à la gorge), stridor inspiratoire puis cyanose ; sans manœuvre, évolue vers la perte de connaissance puis l'arrêt cardiaque par hypoxie.
    \item \textbf{Gasps :} mouvements respiratoires agoniaques, lents ($<$ \SI{6}{min}), saccadés, bruyants, inefficaces, reflétant l'ischémie du tronc cérébral ; ils \textbf{équivalent à une apnée} et imposent la RCP.
    \item \textbf{PLS :} position latérale stable sur le côté, dont l'objectif physiologique principal est d'assurer la \textbf{perméabilité des voies aériennes} et le \textbf{drainage gravitaire} des liquides (salive, sang, vomissures) vers l'extérieur, prévenant la fausse route chez la victime inconsciente qui respire.
    \item \textbf{Chaîne de survie :} 1. Alerte précoce ; 2. RCP précoce ; 3. Défibrillation précoce ; 4. Prise en charge médicale spécialisée précoce.
\end{enumerate}
\textbf{Seuils critiques :}
\begin{enumerate}
    \setcounter{enumi}{5}
    \item Fréquence des compressions : \textbf{\SIrange{100}{120}{min^{-1}}}.
    \item Profondeur : \textbf{\SIrange{5}{6}{cm}} chez l'adulte (sans excéder \SI{6}{cm}).
    \item Ratio (secouriste seul, adulte) : \textbf{30 compressions / 2 insufflations}.
    \item Délai d'irréversibilité cérébrale sans flux sanguin : \textbf{3 à 5 minutes} (réserve d'O$_2$ épuisée, mort neuronale du noyau gris).
\end{enumerate}
\end{corrige}

\begin{corrige}[Exercice 4 — Calcul direct : impact du délai d'intervention]
\begin{enumerate}
    \item \textbf{Sans RCP, $t = \SI{4}{min}$ :}
    \[ S(4) = 100 \times 0{,}9^4 = 100 \times 0{,}6561 = \SI{65,61}{\%} \]
    (Vérification : $0{,}9^2 = 0{,}81$ ; $0{,}81^2 = 0{,}6561$.) La survie est déjà tombée à $\approx$ \textbf{\SI{66}{\%}} en seulement \SI{4}{minutes}.
    \item \textbf{Avec RCP immédiate maintenue jusqu'à $t = \SI{8}{min}$ :}
    \[ S_{\text{RCP}}(8) = 100 \times 0{,}96^8 \]
    Calcul : $0{,}96^2 = 0{,}9216$ ; $0{,}96^4 = 0{,}9216^2 \approx 0{,}8493$ ; $0{,}96^8 = 0{,}8493^2 \approx 0{,}7214$.
    \[ S_{\text{RCP}}(8) \approx \SI{72,1}{\%} \]
    \item \textbf{Interprétation :} à \SI{8}{minutes}, la victime massée dès la première minute conserve $\approx$ \textbf{\SI{72}{\%}} de chances de survie, contre $100 \times 0{,}9^8 \approx \SI{43}{\%}$ sans RCP ; à \SI{15}{minutes} (délai urbain camerounais), la survie sans RCP tombe à $0{,}9^{15} \approx \SI{21}{\%}$, et à \SI{30}{min} (zone rurale) à $0{,}9^{30} \approx \SI{4}{\%}$, alors qu'une RCP de qualité maintient un flux cérébral minimal prolongeant la fenêtre thérapeutique jusqu'à l'arrivée du DAE et du SMUR. \textbf{Conclusion :} dans un contexte où les secours médicalisés sont lents, le \textbf{témoin formé (PSC1) est le maillon décisif} : la formation du grand public (lycéens, motos-taxis, commerçants) est la seule stratégie capable de compenser structurellement la longueur des délais d'intervention.
\end{enumerate}
\end{corrige}

\begin{corrige}[Exercice 5 — Conduite à tenir : arrêt cardiaque chez l'adulte]
\begin{enumerate}
    \item \textbf{Diagnostic : arrêt cardiaque.} Justification : victime \textbf{inconsciente} (pas de réponse verbale ni tactile) + \textbf{absence de respiration normale} pendant \SI{10}{s} d'observation. Les mouvements saccadés, bruyants, espacés de \SIrange{10}{15}{s} (soit \SIrange{4}{6}{min}) sont des \cle{gasps} (respiration agoniaque), signes d'ischémie cérébrale, \textbf{équivalents à une apnée} : ils ne doivent pas retarder la RCP.
    \item \textbf{Conduite à tenir séquencée :}
    \begin{enumerate}
        \item \textbf{Sécuriser} la zone (terrain, pas de danger particulier ici).
        \item \textbf{Alerter immédiatement} (adulte, cause vraisemblablement cardiaque : \emph{phone first}) : appeler le 112/119 en haut-parleur, ou désigner nommément un élève : « Toi, appelle le 119, dis qu'il y a un arrêt cardiaque au lycée, et reviens me confirmer ». Envoyer un autre élève \textbf{chercher le DAE} de l'établissement et prévenir l'infirmier/surveillant.
        \item \textbf{Débuter le MCE} : victime sur le dos, plan dur ; talon d'une main au \textbf{centre de la poitrine} (ligne intermammaire, moitié inférieure du sternum), deuxième main par-dessus, doigts entrelacés, \textbf{bras tendus}, épaules à la verticale des mains ; \textbf{profondeur \SIrange{5}{6}{cm}}, \textbf{fréquence \SIrange{100}{120}{min^{-1}}}, \textbf{relâchement complet} entre chaque compression (la poitrine remonte).
        \item Après 30 compressions : \textbf{2 insufflations} (déclive menton, pincer le nez, \SI{1}{s} par insufflation, soulèvement thoracique visible) si équipé ou volontaire ; sinon \textbf{MCE continu}.
        \item \textbf{DAE arrivé :} allumer $\rightarrow$ suivre les instructions vocales $\rightarrow$ coller les électrodes (thorax nu, une sous clavicule droite, une latérale gauche) $\rightarrow$ \textbf{ne plus toucher la victime} pendant l'analyse $\rightarrow$ si « choc recommandé » : crier « Écartez-vous ! », vérifier visuellement, appuyer sur « choc » $\rightarrow$ \textbf{reprendre le MCE immédiatement} pour \SI{2}{min} (5 cycles de 30:2) avant réanalyse.
        \item \textbf{Relais} toutes les \SI{2}{min} si un second secouriste est disponible (la fatigue dégrade profondeur et fréquence).
        \item \textbf{Poursuivre} jusqu'à : reprise d'une respiration normale/mouvements, arrivée des secours médicalisés, épuisement total ou danger.
    \end{enumerate}
    \item \textbf{Bouche-à-bouche sans protection :} il expose le secouriste au risque de \textbf{transmission d'infections} (VIH, hépatites, tuberculose — enjeu réel au Cameroun) et provoque des \textbf{interruptions prolongées} des compressions chez un secouriste seul non entraîné. Les guidelines ERC/AHA recommandent alors la \textbf{RCP « mains seules » (hands-only)} : compressions thoraciques \emph{continues}, sans insufflations — l'air résiduel et les petits mouvements passifs suffisent pendant les premières minutes chez l'adulte en arrêt d'origine cardiaque.
\end{enumerate}
\end{corrige}

\begin{corrige}[Exercice 6 — Étouffement total : tableau comparatif]
\begin{center}
\begin{tabularx}{\linewidth}{|l|X|X|X|}
\hline
\rowcolor{popblueL}
\textbf{Critère} & \textbf{Nourrisson ($<$ 1 an)} & \textbf{Adulte / Enfant $>$ 1 an} & \textbf{Femme enceinte / obèse} \\
\hline
Signes d'obstruction totale & Silence (pas de cri ni de toux), cyanose, agitation puis hypotonie, signe d'étranglement impossible (mains trop petites) & Signe de l'étranglement (mains à la gorge), parole impossible, toux absente/inefficace, stridor, cyanose, angoisse & Identiques à l'adulte \\
\hline
1\textsuperscript{ère} manœuvre & 5 tapes dorsales : nourrisson à califourchon sur l'avant-bras, tête plus basse que le tronc, menton soutenu, claques fermes entre les omoplates & 5 tapes dorsales : victime penchée en avant, talon de la main entre les omoplates & 5 tapes dorsales identiques \\
\hline
2\textsuperscript{ème} manœuvre & 5 compressions thoraciques : 2 doigts sur la ligne intermammaire, profondeur $\approx$ 1/3 du diamètre antéro-postérieur ($\approx$ \SI{4}{cm}) & 5 compressions abdominales (Heimlich) : poing au-dessus du nombril, sous le xiphoïde, traction franche vers l'arrière et le haut & 5 compressions thoraciques : mains sur le sternum (technique du MCE), poussée franche vers l'arrière \\
\hline
Justification de l'interdiction & Foie et rate volumineux, peu protégés par une cage thoracique souple : risque de rupture hépatique/splénique & — & L'utérus gravide déplace le diaphragme et interdit la prise abdominale ; risque de traumatisme utéro-placentaire et d'inefficacité mécanique \\
\hline
Si inconscience & Allonger au sol, alerter, débuter la RCP (ratio 15:2 si 2 secouristes formés, sinon 30:2) ; regarder dans la bouche avant chaque insufflation et retirer le corps étranger \emph{uniquement s'il est visible} & Idem : RCP 30:2, inspection buccale avant insufflations, jamais de doigt à l'aveugle & Idem ; déclive utérine gauche si possible pendant la RCP \\
\hline
\end{tabularx}
\end{center}
\textbf{Principe physiologique commun :} tapes dorsales = vibrations + gravité ; compressions (thoraciques ou abdominales) = élévation brutale de la pression intrathoracique qui chasse l'air résiduel et expulse le corps étranger (« toux artificielle »).
\end{corrige}

\begin{corrige}[Exercice 7 — Mise en œuvre de la PLS]
\begin{enumerate}
    \item \textbf{Technique (côté droit, secouriste seul) :}
    \begin{enumerate}
        \item S'agenouiller à côté de la victime, jambes de celle-ci tendues.
        \item Placer le \textbf{bras droit} de la victime perpendiculaire à son corps, coude fléchi à 90°, paume vers le haut.
        \item Amener le \textbf{bras gauche} en travers du thorax et placer le \textbf{dos de sa main gauche contre sa joue droite} ; maintenir cette main (elle protège et contrôle la tête).
        \item Saisir la \textbf{cuisse gauche} au-dessus du genou, fléchir la jambe (pied au sol).
        \item \textbf{Basculer} la victime vers soi en tirant sur le genou gauche : elle roule sur le côté droit, la tête reposant sur le dos de la main.
        \item Ajuster : cuisse gauche fléchie à 90° (hanche et genou) pour stabiliser le bassin ; mettre la tête en \textbf{légère extension} (menton relevé) pour ouvrir les voies aériennes ; \textbf{ouvrir la bouche} pour permettre l'écoulement des liquides.
    \end{enumerate}
    \item \textbf{3 objectifs physiologiques :} (1) maintenir la \textbf{perméabilité des voies aériennes} (la langue ne bascule pas en arrière contre le pharynx) ; (2) assurer le \textbf{drainage gravitaire} de la salive, du sang et des vomissures vers l'extérieur ; (3) \textbf{prévenir l'inhalation} (fausse route) et l'asphyxie.
    \item \textbf{Surveillance :} réévaluer la \textbf{respiration toutes les 1 à 2 minutes} (regarder le thorax, écouter, sentir le souffle) et la conscience ; être prêt à retourner la victime sur le dos et débuter la RCP si la respiration s'arrête ou devient agonique.
    \item \textbf{Vomissement en PLS :} \textbf{maintenir la victime en PLS} (ne jamais la redresser), ouvrir davantage la bouche, laisser s'écouler par gravité, dégager éventuellement les restes visibles avec un linge ; la position latérale fait précisément office de protection contre l'inhalation.
\end{enumerate}
\end{corrige}

\begin{corrige}[Exercice 8 — Alerte aux secours : message structuré]
\begin{enumerate}
    \item \textbf{Message d'alerte (au 119, haut-parleur) :} « Bonjour, je suis [nom], témoin d'un \textbf{accident de la circulation} sur l'autoroute Yaoundé--Douala au \textbf{PK 45, sens Yaoundé vers Douala}, devant [repère visible]. Il y a \textbf{deux victimes} : la première est \textbf{inconsciente et ne respire pas — arrêt cardiaque — je suis en train de la masser} ; la deuxième est inconsciente mais respire, avec un \textbf{saignement abondant de la cuisse gauche} que je vais comprimer. J'ai besoin d'une ambulance médicalisée en urgence. Mon numéro est le [numéro]. \textbf{Je ne raccroche pas, je reste en ligne.} »
    \item \textbf{Priorisation :} la \textbf{première action} du témoin seul face à un adulte en arrêt cardiaque est d'\textbf{appeler les secours (phone first)} — en haut-parleur pour ne pas retarder le MCE — car l'arrêt cardiaque d'origine cardiaque nécessite la défibrillation, que seuls les secours (ou un DAE) apporteront. La victime 1 est prioritaire selon la règle de triage : \textbf{arrêt cardiaque $>$ hémorragie massive $>$ inconscient respirant} : sans circulation, la mort cérébrale survient en \SIrange{3}{5}{min}, alors que la victime 2 respire et son hémorragie, bien que grave, laisse un délai d'action plus long ; de plus, comprimer une hémorragie ne sert à rien si l'on laisse mourir la victime réanimable.
    \item \textbf{Répartition avec un deuxième témoin :} le témoin 2 \textbf{prend en charge la victime 2} : compression manuelle ferme du saignement (linge/gant si possible) puis mise en PLS du côté opposé au saignement si l'état le permet, tout en surveillant la respiration ; moi je \textbf{poursuis le MCE} sur la victime 1 ; on se \textbf{relaie toutes les 2 min} sur le massage si possible, et le témoin 2 guide les secours à leur arrivée.
\end{enumerate}
\end{corrige}

\begin{corrige}[Exercice 9 — Sujet type Bac : chaîne de survie et modélisation]
\textbf{1. Analyse de données (4 pts).}
\[ S(3) = 70\,e^{-0{,}25 \times 3} = 70\,e^{-0{,}75} = 70 \times 0{,}4724 \approx \SI{33,1}{\%} \]
\[ S(8) = 70\,e^{-0{,}25 \times 8} = 70\,e^{-2} = 70 \times 0{,}1353 \approx \SI{9,5}{\%} \]
Comparaison : à $t=\SI{3}{min}$, le point clinique ($\approx \SI{33}{\%}$) coïncide presque exactement avec le modèle ; à $t=\SI{5}{min}$, modèle $70e^{-1{,}25} \approx \SI{20,1}{\%}$ vs donnée \SI{20}{\%} : excellent accord. Le modèle exponentiel $S(t)=70e^{-0{,}25t}$ est donc \textbf{validé} par les données cliniques sur l'intervalle 0--10 min (écart faible, décroissance de $\approx \SI{22}{\%}$ par minute : $1-e^{-0{,}25} \approx 0{,}221$).

\textbf{2. Perte moyenne entre 0 et 4 min (3 pts).}
$S(0) = \SI{70}{\%}$ ; $S(4) = 70\,e^{-1} = 70 \times 0{,}3679 \approx \SI{25,75}{\%}$.
\[ \text{Perte moyenne} = \frac{70 - 25{,}75}{4} \approx \SI{11,1}{\%}\ \text{par minute} \]
Interprétation : chaque minute de retard avant défibrillation coûte $\approx$ \textbf{11 points de survie} ; il faut donc former le public à \textbf{reconnaître et alerter instantanément}, et rapprocher les DAE des lieux de vie.

\textbf{3. Physiologie (6 pts).}
\begin{enumerate}
    \item Le MCE maintient un flux cérébral $\approx$ \SI{30}{\%} du normal, juste au-dessus du seuil de survie neuronale ; sans lui, les réserves d'O$_2$ s'épuisent en \SIrange{3}{5}{min}. Or chaque \textbf{interruption} des compressions fait chuter la \textbf{pression de perfusion coronarienne (PPC) à zéro} (Doc. 2), et sa remontée exige plusieurs compressions : les pauses répétées privent le cœur et le cerveau de perfusion, d'où l'exigence d'une fraction de compression $>$ \SI{80}{\%}.
    \item En FV, le myocarde consomme ses réserves d'ATP et s'acidifie ; sans RCP, la FV « fine » dégénère en \textbf{asystolie non choquable} en quelques minutes, et même si le DAE arrive ensuite, le choc tombe sur un myocarde énergétiquement épuisé qui ne peut reprendre une activité mécanique. Le MCE, en maintenant une perfusion coronaire minimale (apport d'O$_2$ et de substrats, élimination partielle des déchets), \textbf{prépare le myocarde à répondre au choc} : les maillons 2 et 3 sont synergiques, non substituables.
\end{enumerate}

\textbf{4. Contexte camerounais (3 pts).} Exemples attendus (2 mesures justifiées) :
\begin{itemize}
    \item \textbf{Formation PSC1 obligatoire} dans les lycées (clubs de secourisme, EPS) et pour les \textbf{motos-taxis} et conducteurs de gares routières, premiers témoins des accidents — renforce les maillons 1 et 2.
    \item \textbf{Installation de DAE publics} dans les marchés (Mokolo, marché central), gares, stades et lycées, avec signalétique claire et enregistrement des numéros 112/119 dans les téléphones — renforce les maillons 1 et 3.
\end{itemize}
\textbf{Points de vigilance (barème /16) :} calculs avec valeurs exactes intermédiaires ; citer le Doc. 2 (PPC, chute à zéro) ; ne pas écrire que le MCE « relance le cœur » ; distinguer rythme choquable/non choquable ; mesures proposées réalistes et reliées aux maillons.
\end{corrige}

\begin{corrige}[Exercice 10 — Sujet type Bac : cas clinique intégré]
\textbf{1. Tableau de synthèse :}
\begin{center}
\begin{tabularx}{\linewidth}{|l|X|X|X|}
\hline
\rowcolor{popblueL}
& \textbf{Situation A} & \textbf{Situation B} & \textbf{Situation C} \\
\hline
Diagnostic & Arrêt cardiaque (sur malaise thoracique : IDM probable) & Étouffement total par corps étranger (arachide), enfant conscient & Perte de connaissance (malaise vagal / coup de chaleur), victime qui respire \\
\hline
Signes clés & Inconscient + pas de respiration en \SI{10}{s} + pas de pouls ; contexte douleur thoracique, sueurs, HTA & Signe de l'étranglement, silence (ni parole ni toux), cyanose, conscience conservée & Inconsciente MAIS respiration régulière \SI{18}{min}, souffle perçu, pouls présent \SI{80}{min} \\
\hline
Conduite & 1. Alerter 119 (haut-parleur) ; 2. MCE immédiat 30:2, \SIrange{5}{6}{cm}, \SIrange{100}{120}{min^{-1}}, relâchement complet ; 3. Relais toutes les \SI{2}{min} ; 4. Poursuivre \SI{15}{min} jusqu'aux secours & 1. Demander « tu t'étouffes ? » ; 2. Alerter ; 3. 5 tapes dorsales (penché en avant) ; 4. Si échec : 5 compressions abdominales (Heimlich, enfant $>$ 1 an) ; 5. Alterner jusqu'à expulsion & 1. Alerter ; 2. PLS côté droit ; 3. Desserrer les vêtements, ombrager, ventiler ; 4. Surveiller respiration toutes les 1--2 min jusqu'aux secours \\
\hline
\end{tabularx}
\end{center}
\textbf{2. Situation A — séquence DAE :} allumer le DAE $\rightarrow$ coller les électrodes sur thorax nu pendant que le MCE continue $\rightarrow$ annoncer « \textbf{Ne touchez plus la victime !} » pendant l'analyse $\rightarrow$ « Choc recommandé » : crier « \textbf{Écartez-vous !} », vérifier visuellement que personne ne touche la victime $\rightarrow$ appuyer sur le bouton choc $\rightarrow$ \textbf{reprendre immédiatement le MCE} pour \SI{2}{minutes} (5 cycles 30:2) sans vérifier le pouls $\rightarrow$ nouvelle analyse à la demande de l'appareil.

\textbf{3. Situation B — perte de conscience :} cesser les manœuvres debout, \textbf{allonger l'enfant au sol}, alerter (si pas déjà fait), \textbf{débuter la RCP} (30 compressions / 2 insufflations ; 15:2 si deux secouristes formés), et \textbf{regarder dans la bouche avant chaque insufflation} pour retirer le corps étranger \emph{uniquement s'il est visible} — jamais de doigt à l'aveugle. Les compressions thoraciques de la RCP génèrent des pics de pression pouvant expulser le corps étranger.

\textbf{4. Situation C :} le côté droit est une \textbf{convention pédagogique} (standardisation de l'enseignement et de l'évaluation ; le côté gauche est physiologiquement acceptable). Risque prioritaire à surveiller : la \textbf{dégradation de la respiration} (arrêt respiratoire $\rightarrow$ retour sur le dos et RCP) et, dans ce contexte de chaleur (\SI{38}{\celsius}), l'aggravation vers un \textbf{coup de chaleur} (hyperthermie majeure) : ombrager, déshabiller partiellement, rafraîchir, surveiller en continu.

\textbf{Barème indicatif (/20) :} tableau 9 pts (1 par case), séquence DAE 4 pts, modification B 3 pts, question C 4 pts. \textbf{Points de vigilance :} justifier chaque diagnostic par les \emph{signes} (pas seulement le nom) ; interdire le Heimlich $<$ 1 an ; rappeler que la PLS exige une respiration normale.
\end{corrige}

\begin{corrige}[Exercice 11 — Sujet type Bac : rédaction argumentée]
\textbf{Corrigé structuré (barème /20 : I = 6 pts, II = 7 pts, III = 5 pts, schéma = 2 pts).}

\textbf{I. Besoins cérébraux et fenêtre thérapeutique.} Le cerveau (2\,\% de la masse corporelle) consomme $\approx$ \textbf{20\,\% du dioxygène} de l'organisme au repos, soit $\approx$ \SI{3,5}{mL O2/100g/min} ; il est dépourvu de réserves significatives (pas de myoglobine, glycogène quasi nul) et son métabolisme est exclusivement aérobie. À l'arrêt circulatoire, l'O$_2$ dissous et l'O$_2$ lié à l'hémoglobine cérébrale sont épuisés en quelques secondes ; la chute d'ATP provoque la défaillance des pompes Na$^+$/K$^+$, la dépolarisation neuronale, la libération de glutamate excitotoxique et l'entrée calcique : au-delà de \textbf{3 à 5 minutes sans flux}, les lésions deviennent \textbf{irréversibles}. Le « temps cerveau » est donc le facteur limitant : toute la stratégie de réanimation vise à \emph{substituer artificiellement} un flux minimal avant cette échéance.

\textbf{II. Mécanique du MCE et hémodynamique.} Deux mécanismes complémentaires expliquent le flux généré par le MCE :
\begin{itemize}
    \item \textbf{Théorie de la pompe cardiaque} : la compression sternale écrase le ventricule gauche entre sternum et colonne vertébrale ; le sang est éjecté vers l'aorte (les valves empêchant le reflux).
    \item \textbf{Théorie de la pompe thoracique} : la compression élève la pression de toute la cavité thoracique ; le cœur agit alors comme conduit passif, le reflux veineux étant bloqué au niveau jugulaire.
\end{itemize}
La variable décisive est la \textbf{pression de perfusion coronarienne} : $PPC = P_{Ao\,diastolique} - P_{OD}$ (pression auriculaire droite). Le myocarde n'est perfusé que pendant la « diastole » du massage, c'est-à-dire pendant le \textbf{relâchement} — d'où l'exigence d'un relâchement complet. À chaque interruption des compressions, la PPC \textbf{chute à zéro} ; sa reconstruction exige plusieurs compressions consécutives (remplissage progressif du compartiment artériel). Les pauses répétées (insufflations lentes, analyses DAE mal gérées) privent donc cœur et cerveau de perfusion : c'est la justification physiologique de la \textbf{fraction de compression $>$ 80\,\%} et du ratio 30:2.

\textbf{III. Synergie RCP--défibrillation.} En fibrillation ventriculaire, le myocarde se contracte de façon anarchique (400--600 dépolarisations/min) sans éjection : la consommation d'O$_2$ reste élevée alors que la perfusion coronaire est nulle ; l'ATP s'épuise, le pH chute (anaérobie), et la FV « grossière » (choquable) dégénère en FV fine puis en \textbf{asystolie}, rythme \textbf{non choquable} de très mauvais pronostic. Le MCE, en maintenant une perfusion coronaire $\approx$ \SI{30}{\%} du normal, préserve un métabolisme aérobie minimal (apport d'O$_2$ et de glucose, wash-out partiel du lactate, maintien du pH et de l'ATP) : il \textbf{« prépare » le myocarde à répondre au choc} électrique, qui dépolarise simultanément toutes les cellules et permet au nœud sinusal de reprendre le commandement. D'où la conclusion opérationnelle : RCP \textbf{immédiate}, \textbf{quasi continue} (maximiser le temps « mains sur le thorax »), et défibrillation \textbf{la plus précoce possible} — les deux premiers maillons conditionnant l'efficacité du troisième.

\begin{popfigure}
\begin{tikzpicture}[font=\small, >=Stealth, node distance=1.1cm]
\node[draw=popblue, fill=popblueL, rounded corners, align=center, minimum width=3.2cm] (ac) {Arrêt cardiaque\\ (FV) : flux = 0};
\node[draw=poporange, fill=poporangeL, rounded corners, align=center, right=2.2cm of ac, minimum width=3.2cm] (mce) {MCE 100--120/min\\ 5--6 cm, relâchement\\ complet};
\node[draw=popgreen, fill=popgreenL, rounded corners, align=center, right=2.2cm of mce, minimum width=3.2cm] (ppc) {PPC $>$ 15 mmHg\\ ATP / pH myocarde\\ préservés};
\node[draw=poppurple, fill=poppurpleL, rounded corners, align=center, below=1.4cm of ppc, minimum width=3.2cm] (dae) {Choc DAE précoce\\ ($<$ 3--5 min)};
\node[draw=popdark, fill=popcream, rounded corners, align=center, left=2.2cm of dae, minimum width=3.2cm] (rv) {Reprise activité\\ cardiaque efficace};
\draw[->, very thick, poporange] (ac) -- (mce) node[midway, above]{immédiat};
\draw[->, very thick, popgreen] (mce) -- (ppc) node[midway, above]{sans interruption};
\draw[->, very thick, poppurple] (ppc) -- (dae);
\draw[->, very thick, popdark] (dae) -- (rv);
\draw[->, dashed, thick, poppink] (rv.north) to[bend right=20] node[midway, left, align=center]{sinon reprise\\ MCE 2 min} (ppc.west);
\end{tikzpicture}
\legende{Schéma fonctionnel de la synergie RCP--défibrillation : le MCE immédiat et continu maintient la pression de perfusion coronarienne (PPC), préservant le métabolisme myocardique jusqu'au choc du DAE ; en l'absence de reprise d'activité, le cycle MCE 2 min -- analyse est répété.}
\end{popfigure}

\textbf{Points de vigilance :} exiger les valeurs chiffrées (20\,\% de l'O$_2$, 3--5 min, 30\,\% du débit, PPC $>$ 15--20 mmHg, CCF $>$ 80\,\%) ; interdire la formulation « le massage relance le cœur » ; la définition de la PPC doit être littérale ; le schéma doit comporter les flèches de causalité et la boucle de reprise du MCE après le choc.
\end{corrige}

\newpage

\coursec{Fiche bilan}

\begin{popfigure}
\begin{center}\tcbox[colback=popink,colframe=popink,coltext=white,arc=9pt,boxsep=4pt,left=6pt,right=6pt]{\bfseries\small Secourisme \& Accidents Cardiovasculaires}\end{center}
\vspace{-2pt}
\begin{tcbraster}[raster columns=3,raster equal height=rows,raster column skip=6pt,raster row skip=6pt,enhanced,blankest]
\begin{tcolorbox}[enhanced,colback=white,colframe=popblue,boxrule=0.9pt,arc=3pt,title={Chaîne de survie (4 maillons)},fonttitle=\bfseries\scriptsize,coltitle=white,colbacktitle=popblue,left=4pt,right=4pt,top=2pt,bottom=2pt,boxsep=1pt,toptitle=1.5pt,bottomtitle=1.5pt,halign=left]
\begin{itemize}[nosep,leftmargin=*,itemsep=1pt,font=\scriptsize]
\item Alerte précoce
\item RCP immédiate
\item Défibrillation
\item Soins avancés
\end{itemize}
\end{tcolorbox}
\begin{tcolorbox}[enhanced,colback=white,colframe=poporange,boxrule=0.9pt,arc=3pt,title={Arrêt cardiaque},fonttitle=\bfseries\scriptsize,coltitle=white,colbacktitle=poporange,left=4pt,right=4pt,top=2pt,bottom=2pt,boxsep=1pt,toptitle=1.5pt,bottomtitle=1.5pt,halign=left]
\begin{itemize}[nosep,leftmargin=*,itemsep=1pt,font=\scriptsize]
\item Reconnaissance
\item RCP 30/2
\item DAE
\end{itemize}
\end{tcolorbox}
\begin{tcolorbox}[enhanced,colback=white,colframe=popgreen,boxrule=0.9pt,arc=3pt,title={Étouffement total},fonttitle=\bfseries\scriptsize,coltitle=white,colbacktitle=popgreen,left=4pt,right=4pt,top=2pt,bottom=2pt,boxsep=1pt,toptitle=1.5pt,bottomtitle=1.5pt,halign=left]
\begin{itemize}[nosep,leftmargin=*,itemsep=1pt,font=\scriptsize]
\item Adulte/Enfant
\item Nourrisson
\item Femme enceinte
\end{itemize}
\end{tcolorbox}
\begin{tcolorbox}[enhanced,colback=white,colframe=poppurple,boxrule=0.9pt,arc=3pt,title={Perte de connaissance},fonttitle=\bfseries\scriptsize,coltitle=white,colbacktitle=poppurple,left=4pt,right=4pt,top=2pt,bottom=2pt,boxsep=1pt,toptitle=1.5pt,bottomtitle=1.5pt,halign=left]
\begin{itemize}[nosep,leftmargin=*,itemsep=1pt,font=\scriptsize]
\item Reconnaissance
\item PLS
\end{itemize}
\end{tcolorbox}
\begin{tcolorbox}[enhanced,colback=white,colframe=poppink,boxrule=0.9pt,arc=3pt,title={Conduite PAS},fonttitle=\bfseries\scriptsize,coltitle=white,colbacktitle=poppink,left=4pt,right=4pt,top=2pt,bottom=2pt,boxsep=1pt,toptitle=1.5pt,bottomtitle=1.5pt,halign=left]
\begin{itemize}[nosep,leftmargin=*,itemsep=1pt,font=\scriptsize]
\item Protéger
\item Alerter
\item Secourir
\end{itemize}
\end{tcolorbox}
\begin{tcolorbox}[enhanced,colback=white,colframe=popgold,boxrule=0.9pt,arc=3pt,title={Prévention},fonttitle=\bfseries\scriptsize,coltitle=white,colbacktitle=popgold,left=4pt,right=4pt,top=2pt,bottom=2pt,boxsep=1pt,toptitle=1.5pt,bottomtitle=1.5pt,halign=left]
\begin{itemize}[nosep,leftmargin=*,itemsep=1pt,font=\scriptsize]
\item Facteurs risque
\item Hygiène de vie
\end{itemize}
\end{tcolorbox}
\end{tcbraster}

\legende{Carte mentale récapitulative de la leçon 14 : Secourisme et accidents cardiovasculaires.}
\end{popfigure}

\begin{bilan}

\textbf{Définitions clés}
\begin{itemize}
  \item \cle{Arrêt cardiaque (AC)} : Arrêt brutal et inopiné de la fonction de pompe du cœur $\rightarrow$ arrêt de la circulation sanguine, perte de connaissance immédiate, absence de respiration normale (ou respiration agonique).
  \item \cle{Étouffement total} : Obstruction complète des voies aériennes par un corps étranger $\rightarrow$ victime incapable de parler, tousser, respirer ; signe de la main à la gorge (signe universel).
  \item \cle{Perte de connaissance (Pdc)} : Interruption brutale et transitoire de l'état de conscience, avec conservation de la respiration et de la circulation (pouls carotidien présent). La victime ne répond pas, ne réagit pas aux stimulations.
  \item \cle{Chaîne de survie} : Suite de 4 maillons indispensables pour maximiser les chances de survie : 1. Alerte précoce (112/119), 2. RCP précoce, 3. Défibrillation précoce (DAE), 4. Prise en charge médicale avancée.
  \item \cle{Position Latérale de Sécurité (PLS)} : Position standardisée du corps d'une victime inconsciente qui respire, permettant le drainage des voies aériennes (langue, vomissements) et l'attente des secours en sécurité.
  \item \cle{RCP (Réanimation Cardio-Pulmonaire)} : Association de compressions thoraciques (massage cardiaque) et d'insufflations (bouche-à-bouche ou masque) pour maintenir une circulation et une oxygénation minimales. Ratio universel \textbf{30 compressions / 2 insufflations}.
  \item \cle{DAE (Défibrillateur Automatisé Externe)} : Appareil portable analysant le rythme cardiaque et délivrant un choc électrique si nécessaire (fibrillation ventriculaire, tachycardie ventriculaire sans pouls). Utilisable par tout témoin.
\end{itemize}

\textbf{Algorithmes de prise en charge (à connaître par cœur)}
\begin{center}
\begin{tabularx}{\linewidth}{|l|X|}\hline
\textbf{Situation} & \textbf{Étapes clés (Ordre impératif)} \\ \hline
\rowcolor{popblueL}
\textbf{Arrêt Cardiaque Adulte} &
1. \textbf{Sécuriser} les lieux. \\
2. \textbf{Évaluer} : Inconscient ? Ne respire pas (ou mal) ? $\rightarrow$ \textbf{AC confirmé}. \\
3. \textbf{Alerter} : 112 / 119 (Cameroon) / 1500 (SAMU) + \textbf{Demander un DAE}. \\
4. \textbf{RCP} : 30 compressions (centre thorax, 5--6 cm, 100--120/min) / 2 insufflations (1 sec chacune, élévation thorax). \textbf{Ne pas s'arrêter} sauf DAE/Reprise/Épuisement. \\
5. \textbf{DAE} : Allumer $\rightarrow$ Poser électrodes (schéma sur boîtier) $\rightarrow$ Laisser analyser $\rightarrow$ Choc si indiqué $\rightarrow$ Reprendre RCP immédiat. \\ \hline
\rowcolor{poporangeL}
\textbf{Étouffement Total Adulte/Enfant $>1$ an} &
1. \textbf{Encourager à tousser} (si partiel). Si total (pas de toux, pas de voix, cyanose) : \\
2. \textbf{5 claques dorsales} (talon main, entre omoplates, buste penché). \\
3. \textbf{5 compressions abdominales (Heimlich)} : Poing fermé au-dessus nombril, sous apophyse xiphoïde, traction violente postéro-supérieure. \\
4. Alterner 5 claques / 5 compressions jusqu'à expulsion ou perte de connaissance. \\
\textit{Si Pdc} $\rightarrow$ Allonger, Alerter, RCP (regarder bouche avant insufflations). \\ \hline
\rowcolor{popgreenL}
\textbf{Étouffement Total Nourrisson $<1$ an} &
1. \textbf{5 claques dorsales} : Nourrisson sur avant-bras, tête plus bas que corps, main tenant mâchoire. \\
2. \textbf{5 compressions thoraciques} : 2 doigts (index/majeur) sur sternum (ligne intermammaire), profondeur 1/3 diamètre antéro-postérieur ($\approx$ 4 cm). \\
3. Alterner jusqu'à expulsion ou Pdc. \textbf{Jamais de compression abdominale}. \\ \hline
\rowcolor{poppurpleL}
\textbf{Étouffement Femme enceinte / Obèse} &
Remplacer compressions abdominales par \textbf{compressions thoraciques} (poing sur sternum, comme RCP mais conscient). \\
Claques dorsales identiques. \\ \hline
\rowcolor{poppinkL}
\textbf{Perte de Connaissance (Victime qui respire)} &
1. \textbf{Libérer voies aériennes} (dégager bouche, menton en avant). \\
2. \textbf{Vérifier respiration} (Voir, Entendre, Sentir $\le$ 10 sec). \\
3. Si respire normalement $\rightarrow$ \textbf{PLS} : \\
\quad $\bullet$ Bras droit perpendiculaire, main sous joue gauche. \\
\quad $\bullet$ Jambe gauche fléchie, pied au sol. \\
\quad $\bullet$ Basculer vers soi (tirer genou gauche), stabiliser (hanche/genou 90°). \\
\quad $\bullet$ Tête en extension (menton en avant), bouche ouverte vers le bas. \\
4. \textbf{Alerter} (si pas fait), \textbf{Surveiller} respiration en continu. \\ \hline
\end{tabularx}
\end{center}

\textbf{Méthodes express}
\begin{itemize}
  \item \textbf{Méthode PAS} : \cle{Protéger} (balisage, risques) $\rightarrow$ \cle{Alerter} (112/119 : Où ? Quoi ? Combien ? État ? Ne pas raccrocher) $\rightarrow$ \cle{Secourir} (gestes adaptés).
  \item \textbf{Reconnaissance AC} : \cle{TAPER - CRIER} (éveil) $\rightarrow$ \cle{REGARDER - ENTENDRE - SENTIR} (respiration $\le$ 10 s). \textbf{Doute = AC}.
  \item \textbf{Qualité RCP} : \textbf{Pousser fort, pousser vite, laisser remonter, ne pas interrompre}. Changement de secouriste tous les 2 min (5 cycles) si possible.
  \item \textbf{DAE} : Ne pas toucher la victime pendant l'analyse/choc. Sécher thorax si mouillé. Éloigner patchs médicamenteux/pacemaker ($\ge$ 8 cm).
\end{itemize}

\end{bilan}

\begin{aretenir}[Checklist avant le Bac]
\begin{itemize}
  \item $\square$ Je cite les 4 maillons de la \cle{chaîne de survie} dans l'ordre.
  \item $\square$ Je sais reconnaître un \cle{arrêt cardiaque} : Inconscience + Absence respiration normale (ou respiration agonique).
  \item $\square$ Je connais les paramètres de la \cle{RCP adulte} : Ratio 30/2, Fréquence 100--120/min, Profondeur 5--6 cm, Décompression complète.
  \item $\square$ Je sais utiliser un \cle{DAE} : Allumer $\rightarrow$ Coller électrodes (schéma) $\rightarrow$ Ne pas toucher $\rightarrow$ Choc $\rightarrow$ RCP immédiat.
  \item $\square$ Je différencie \cle{étouffement partiel} (toux efficace, parole possible) et \cle{total} (pas de toux, pas de voix, cyanose, signe gorge).
  \item $\square$ Je réalise la manœuvre pour \cle{étouffement adulte/enfant} : 5 claques dorsales / 5 compressions abdominales (Heimlich).
  \item $\square$ Je réalise la manœuvre pour \cle{étouffement nourrisson} : 5 claques dorsales / 5 compressions thoraciques (2 doigts). \textbf{Pas d'abdominal}.
  \item $\square$ J'adapte pour \cle{femme enceinte/obèse} : Compressions thoraciques (sternum) au lieu d'abdominales.
  \item $\square$ Je reconnais une \cle{perte de connaissance} : Inconscient MAIS respiration normale présente (pouls carotidien palpable).
  \item $\square$ Je sais mettre en \cle{PLS} dans l'ordre exact (bras, joue, genou, bascule, extension tête, bouche basse).
  \item $\square$ Je connais les numéros d'urgence au Cameroun : \textbf{112} (GSM/Standard), \textbf{119} (Pompiers), \textbf{1500} (SAMU Yaoundé).
  \item $\square$ J'adopte l'attitude \cle{citoyen-secouriste} : Empathie, calme, protection, alerte précise, gestes sans délai, surveillance continue.
\end{itemize}
\end{aretenir}

\findecours

\end{document}
